\documentclass[11pt,a4paper]{article}

% Packages
\usepackage[utf8]{inputenc}
\usepackage{amsmath}
\usepackage{amssymb}
\usepackage{physics}
\usepackage{graphicx}
\usepackage{mathtools}  
\usepackage{geometry}
\usepackage{mathrsfs}
\usepackage{braket}  % This already provides \ket, \bra, \braket
\usepackage{hyperref}
\usepackage{enumitem}
\usepackage{tcolorbox}

% Page setup
\geometry{margin=1in}
\setlength{\parindent}{0pt}
\setlength{\parskip}{6pt}

% Custom commands for QM (only adding what's not in braket package)
\newcommand{\expect}[1]{\left\langle #1 \right\rangle}

\title{Quantum Mechanics Exam Notes}
\author{Tianshuo Wang}
\date{\today}

\begin{document}

\maketitle
\tableofcontents
\newpage


% ============================================
% ============================================
\newpage
\section{Dirac Delta Function and Gaussian Integrals}

\subsection{The Dirac Delta Function}

\subsubsection{Essential Properties}

\textbf{Sifting property (most important):}
\begin{equation}
\boxed{\int_{-\infty}^{\infty} f(x)\delta(x - a)\, dx = f(a)}
\tag{D.1}
\end{equation}

\textbf{Normalization:}
\begin{equation}
\int_{-\infty}^{\infty} \delta(x - a)\, dx = 1
\tag{D.2}
\end{equation}

\textbf{Symmetry:}
\begin{equation}
\delta(-x) = \delta(x), \quad \delta(x - a) = \delta(a - x)
\tag{D.3}
\end{equation}

\textbf{Scaling:}
\begin{equation}
\boxed{\delta(ax) = \frac{1}{|a|}\delta(x)}
\tag{D.4}
\end{equation}

\textbf{Delta function of a function:}
\begin{equation}
\boxed{\delta(g(x)) = \sum_{i}\frac{\delta(x - x_i)}{|g'(x_i)|}}
\tag{D.5}
\end{equation}
where $x_i$ are simple zeros of $g(x)$.

\textbf{Product property:}
\begin{equation}
f(x)\delta(x - a) = f(a)\delta(x - a), \quad x\delta(x) = 0
\tag{D.6}
\end{equation}

\textbf{Derivative:}
\begin{equation}
\int_{-\infty}^{\infty} f(x)\delta'(x - a)\, dx = -f'(a), \quad x\delta'(x) = -\delta(x)
\tag{D.7}
\end{equation}

\subsubsection{Three-Dimensional Delta Function}

\textbf{Cartesian:}
\begin{equation}
\boxed{\delta^3(\mathbf{r} - \mathbf{r}_0) = \delta(x - x_0)\delta(y - y_0)\delta(z - z_0)}
\tag{D.8}
\end{equation}

\textbf{Spherical (radially symmetric):}
\begin{equation}
\delta^3(\mathbf{r}) = \frac{\delta(r)}{4\pi r^2}
\tag{D.9}
\end{equation}

\subsubsection{Important Representations}

\textbf{Fourier representation (crucial for QM):}
\begin{equation}
\boxed{\delta(x - x') = \frac{1}{2\pi}\int_{-\infty}^{\infty} e^{ik(x - x')}\, dk}
\tag{D.10}
\end{equation}

\textbf{3D Fourier:}
\begin{equation}
\delta^3(\mathbf{r} - \mathbf{r}') = \frac{1}{(2\pi)^3}\int e^{i\mathbf{k}\cdot(\mathbf{r} - \mathbf{r}')}\, d^3\mathbf{k}
\tag{D.11}
\end{equation}

\subsection{Gaussian Integrals}

\subsubsection{Fundamental Formulas}

\textbf{Basic Gaussian:}
\begin{equation}
\boxed{\int_{-\infty}^{\infty} e^{-ax^2}\, dx = \sqrt{\frac{\pi}{a}} \quad (a > 0)}
\tag{G.1}
\end{equation}

\textbf{With linear term:}
\begin{equation}
\boxed{\int_{-\infty}^{\infty} e^{-ax^2 + bx}\, dx = \sqrt{\frac{\pi}{a}}e^{b^2/(4a)}}
\tag{G.2}
\end{equation}

\textbf{Complex version (for QM):}
\begin{equation}
\boxed{\int_{-\infty}^{\infty} e^{-ax^2 + ibx}\, dx = \sqrt{\frac{\pi}{a}}e^{-b^2/(4a)}}
\tag{G.3}
\end{equation}

\textbf{Even powers:}
\begin{align}
\int_{-\infty}^{\infty} x^2 e^{-ax^2}\, dx &= \frac{1}{2}\sqrt{\frac{\pi}{a^3}} \tag{G.4a} \\
\int_{-\infty}^{\infty} x^{2n}e^{-ax^2}\, dx &= \frac{(2n-1)!!}{2^n a^n}\sqrt{\frac{\pi}{a}} \tag{G.4b}
\end{align}

\textbf{Fourier transform of Gaussian:}
\begin{equation}
\boxed{\int_{-\infty}^{\infty} e^{-ax^2}e^{ikx}\, dx = \sqrt{\frac{\pi}{a}}e^{-k^2/(4a)}}
\tag{G.5}
\end{equation}

\textbf{Multi-dimensional:}
\begin{equation}
\int_{\mathbb{R}^n} e^{-a|\mathbf{x}|^2}\, d^n\mathbf{x} = \left(\frac{\pi}{a}\right)^{n/2}
\tag{G.6}
\end{equation}

\subsection{QM Applications - Quick Reference}

\textbf{Normalization constant for Gaussian:}
\begin{equation}
\psi(x) = Ae^{-\alpha x^2} \quad \Rightarrow \quad A = \left(\frac{2\alpha}{\pi}\right)^{1/4}
\tag{A.1}
\end{equation}

\textbf{Harmonic oscillator ground state:}
\begin{equation}
\psi_0(x) = \left(\frac{m\omega}{\pi\hbar}\right)^{1/4}e^{-m\omega x^2/(2\hbar)}, \quad \langle x^2\rangle = \frac{\hbar}{2m\omega}
\tag{A.2}
\end{equation}

\textbf{Momentum-position transformation:}
\begin{equation}
\langle x'|p'\rangle = \frac{1}{\sqrt{2\pi\hbar}}e^{ip'x'/\hbar}
\tag{A.3}
\end{equation}

% ============================================
\newpage
\section{Jacobian and Change of Variables}

\subsection{Jacobian Determinant}

\subsubsection{Two Dimensions}

\textbf{Definition:}
\begin{equation}
\boxed{J = \frac{\partial(x, y)}{\partial(u, v)} = \begin{vmatrix}
\frac{\partial x}{\partial u} & \frac{\partial x}{\partial v} \\[6pt]
\frac{\partial y}{\partial u} & \frac{\partial y}{\partial v}
\end{vmatrix} = \frac{\partial x}{\partial u}\frac{\partial y}{\partial v} - \frac{\partial x}{\partial v}\frac{\partial y}{\partial u}}
\tag{J.1}
\end{equation}

\textbf{Change of variables:}
\begin{equation}
\boxed{\iint_R f(x, y)\, dx\, dy = \iint_{R'} f(x(u,v), y(u,v))\left|J\right|\, du\, dv}
\tag{J.2}
\end{equation}

\subsubsection{Three Dimensions}

\textbf{Definition:}
\begin{equation}
\boxed{J = \frac{\partial(x, y, z)}{\partial(u, v, w)} = \begin{vmatrix}
\frac{\partial x}{\partial u} & \frac{\partial x}{\partial v} & \frac{\partial x}{\partial w} \\[6pt]
\frac{\partial y}{\partial u} & \frac{\partial y}{\partial v} & \frac{\partial y}{\partial w} \\[6pt]
\frac{\partial z}{\partial u} & \frac{\partial z}{\partial v} & \frac{\partial z}{\partial w}
\end{vmatrix}}
\tag{J.3}
\end{equation}

\textbf{Change of variables:}
\begin{equation}
\boxed{\iiint_V f(x, y, z)\, dx\, dy\, dz = \iiint_{V'} f(x,y,z)\left|J\right|\, du\, dv\, dw}
\tag{J.4}
\end{equation}

\subsection{Essential Coordinate Systems}

\subsubsection{Polar Coordinates (2D)}

\textbf{Transformation:} $x = r\cos\theta$, $y = r\sin\theta$

\textbf{Jacobian:} $J = r$

\textbf{Area element:}
\begin{equation}
\boxed{dA = r\, dr\, d\theta}
\tag{J.5}
\end{equation}

\subsubsection{Cylindrical Coordinates (3D)}

\textbf{Transformation:} $x = \rho\cos\phi$, $y = \rho\sin\phi$, $z = z$

\textbf{Jacobian:} $J = \rho$

\textbf{Volume element:}
\begin{equation}
\boxed{dV = \rho\, d\rho\, d\phi\, dz}
\tag{J.6}
\end{equation}

\subsubsection{Spherical Coordinates (3D)}

\textbf{Transformation:}
\begin{equation}
x = r\sin\theta\cos\phi, \quad y = r\sin\theta\sin\phi, \quad z = r\cos\theta
\tag{J.7}
\end{equation}

\textbf{Jacobian:} $J = r^2\sin\theta$

\textbf{Volume element:}
\begin{equation}
\boxed{dV = r^2\sin\theta\, dr\, d\theta\, d\phi}
\tag{J.8}
\end{equation}

\textbf{Ranges:} $r \geq 0$, $0 \leq \theta \leq \pi$, $0 \leq \phi < 2\pi$

\subsection{Key Properties}

\textbf{Inverse Jacobian:}
\begin{equation}
J_{uv \to xy} = \frac{1}{J_{xy \to uv}}
\tag{J.9}
\end{equation}

\textbf{Chain rule:}
\begin{equation}
\frac{\partial(x,y)}{\partial(s,t)} = \frac{\partial(x,y)}{\partial(u,v)} \cdot \frac{\partial(u,v)}{\partial(s,t)}
\tag{J.10}
\end{equation}

\subsection{QM Applications}

\textbf{Spherical normalization:}
\begin{equation}
\int_0^{2\pi}\int_0^{\pi}\int_0^{\infty} |\psi(r, \theta, \phi)|^2 r^2\sin\theta\, dr\, d\theta\, d\phi = 1
\tag{Q.1}
\end{equation}

\textbf{Radial probability (spherically symmetric):}
\begin{equation}
P(r)\, dr = |\psi(r)|^2 \cdot 4\pi r^2\, dr
\tag{Q.2}
\end{equation}

\textbf{Surface area of sphere:}
\begin{equation}
\int_0^{2\pi}\int_0^{\pi} r^2\sin\theta\, d\theta\, d\phi = 4\pi r^2
\tag{Q.3}
\end{equation}

\subsection{Common Exam Examples}

\textbf{Example 1: Gaussian in 2D}
\begin{equation}
\int_{-\infty}^{\infty}\int_{-\infty}^{\infty} e^{-(x^2 + y^2)}\, dx\, dy = \int_0^{2\pi}\int_0^{\infty} e^{-r^2} r\, dr\, d\theta = \pi
\tag{E.1}
\end{equation}

\textbf{Example 2: Volume of sphere}
\begin{equation}
V = \int_0^{2\pi}\int_0^{\pi}\int_0^R r^2\sin\theta\, dr\, d\theta\, d\phi = \frac{4\pi R^3}{3}
\tag{E.2}
\end{equation}

\textbf{Example 3: Hydrogen atom normalization}
\begin{equation}
\psi_{100} = \frac{1}{\sqrt{\pi a_0^3}}e^{-r/a_0} \quad \text{(normalized)}
\tag{E.3}
\end{equation}
% ============================================
\newpage
\section{Kets, Bras, and Operators}

\subsection{State Vectors (Kets)}

\textbf{Key Concept:} Physical states are represented by state vectors called \textbf{kets}: $\ket{\alpha}$

\textbf{Properties:}
\begin{itemize}
\item Kets can be added: $\ket{\alpha} + \ket{\beta} = \ket{\gamma}$
\item Scalar multiplication: $c\ket{\alpha} = \ket{\alpha}c$
\item $\ket{\alpha}$ and $c\ket{\alpha}$ (with $c \neq 0$) represent the \textit{same physical state}
\item Vector space: \textbf{Hilbert space} (dimension set by the system: 2 for spin-$\frac{1}{2}$, infinite for position/momentum)
\end{itemize}

\subsection{Operators and Eigenkets}

\textbf{Observables} (momentum, spin, etc.) are represented by \textbf{operators} $A$ that act on kets from the left:
$$A\ket{\alpha} = \text{another ket}$$

\textbf{Eigenvalue Equation:}
\begin{equation}
A\ket{a'} = a'\ket{a'}
\tag{1.19}
\end{equation}

\begin{itemize}
\item $\ket{a'}$ = eigenket (eigenstate)
\item $a'$ = eigenvalue (measurement result)
\item Set of eigenvalues: $\{a', a'', a''', \ldots\}$
\end{itemize}

\textbf{Example - Spin-1/2:}
\begin{align}
S_z\ket{S_z; +} &= \frac{\hbar}{2}\ket{S_z; +} \tag{1.20} \\
S_z\ket{S_z; -} &= -\frac{\hbar}{2}\ket{S_z; -} \notag
\end{align}

\textbf{Eigenket Expansion:} Any state can be expanded in eigenkets:
\begin{equation}
\ket{\alpha} = \sum_{a'} c_{a'}\ket{a'}
\tag{1.22}
\end{equation}

\subsection{Bra Space and Dual Correspondence}

\textbf{Bra-Ket Duality:}
\begin{itemize}
\item Every ket $\ket{\alpha}$ has a corresponding bra $\bra{\alpha}$ in the dual space
\item \textbf{Important:} $c\ket{\alpha} \xleftrightarrow{\text{DC}} c^*\bra{\alpha}$ (complex conjugate!)
\end{itemize}

\subsection{Inner Product}

\textbf{Definition:} $\langle\beta|\alpha\rangle$ (bra on left, ket on right)

\textbf{Key Properties:}
\begin{align}
\langle\beta|\alpha\rangle &= \langle\alpha|\beta\rangle^* \tag{1.26} \\
\langle\alpha|\alpha\rangle &\geq 0 \tag{1.27}
\end{align}

\textbf{Orthogonality:} $\ket{\alpha}$ and $\ket{\beta}$ are orthogonal if:
\begin{equation}
\langle\alpha|\beta\rangle = 0 \tag{1.28}
\end{equation}

\textbf{Normalization:}
\begin{equation}
\ket{\tilde{\alpha}} = \frac{1}{\sqrt{\langle\alpha|\alpha\rangle}}\ket{\alpha} \quad \Rightarrow \quad \langle\tilde{\alpha}|\tilde{\alpha}\rangle = 1
\tag{1.30, 1.31}
\end{equation}

\textbf{Norm:} $\sqrt{\langle\alpha|\alpha\rangle}$ is the norm of $\ket{\alpha}$

\subsection{Operator Multiplication}

\textbf{Key Properties:}
\begin{itemize}
\item \textbf{Non-commutative:} In general, $XY \neq YX$ \hfill (1.41)
\item \textbf{Associative:} $X(YZ) = (XY)Z = XYZ$ \hfill (1.42)
\end{itemize}

\textbf{Associativity with kets and bras:}
\begin{equation}
X(Y\ket{\alpha}) = (XY)\ket{\alpha} = XY\ket{\alpha}, \quad (\bra{\beta}X)Y = \bra{\beta}(XY) = \bra{\beta}XY
\tag{1.43}
\end{equation}

\textbf{Hermitian conjugate of product:}
\begin{equation}
(XY)^\dagger = Y^\dagger X^\dagger \quad \text{(reverse order!)}
\tag{1.44}
\end{equation}

\subsection{Outer Product}

\textbf{Definition:} Multiply $\ket{\beta}$ and $\bra{\alpha}$ in that order:
\begin{equation}
\ket{\beta}\bra{\alpha} = \text{an operator (NOT a number!)}
\tag{1.46}
\end{equation}

\textbf{Compare:}
\begin{itemize}
\item Inner product: $\langle\beta|\alpha\rangle$ = complex number
\item Outer product: $\ket{\beta}\bra{\alpha}$ = operator
\end{itemize}

\textbf{Illegal products:}
\begin{itemize}
\item $\ket{\alpha}X$ or $X\bra{\alpha}$ (operator on wrong side)
\item $\ket{\alpha}\ket{\beta}$ or $\bra{\alpha}\bra{\beta}$ (two kets/bras from same space)
\end{itemize}

\subsection{The Associative Axiom}

The \textbf{associative axiom of multiplication} allows us to drop parentheses in products of kets, bras, and operators.

\textbf{Outer product acting on a ket:}
\begin{equation}
(\ket{\beta}\bra{\alpha}) \cdot \ket{\gamma} = \ket{\beta}(\langle\alpha|\gamma\rangle)
\tag{1.47, 1.48}
\end{equation}
The outer product $\ket{\beta}\bra{\alpha}$ acts as an operator that "rotates" $\ket{\gamma}$ into the direction of $\ket{\beta}$.

\textbf{Matrix elements:} Compact notation for operator between bra and ket:
\begin{equation}
\langle\beta|X|\alpha\rangle \equiv (\bra{\beta}) \cdot (X\ket{\alpha}) = (\bra{\beta}X) \cdot (\ket{\alpha})
\tag{1.51, 1.52}
\end{equation}

\textbf{Important identity:}
\begin{equation}
\langle\beta|X|\alpha\rangle = \langle\alpha|X^\dagger|\beta\rangle^*
\tag{1.53}
\end{equation}

\textbf{For Hermitian operators} ($X = X^\dagger$):
\begin{equation}
\langle\beta|X|\alpha\rangle = \langle\alpha|X|\beta\rangle^*
\tag{1.54}
\end{equation}

\textbf{Hermitian conjugate of outer product:}
\begin{equation}
X = \ket{\beta}\bra{\alpha} \quad \Rightarrow \quad X^\dagger = \ket{\alpha}\bra{\beta}
\tag{1.49, 1.50}
\end{equation}
% ============================================
% ============================================
\newpage
\section{Base Kets, Orthonormality, and Matrix Representations}

\subsection{Eigenkets of an Observable}

\textbf{Theorem 1:} For a Hermitian operator $A$:
\begin{enumerate}
\item The eigenvalues of $A$ are \textbf{real}
\item Eigenkets corresponding to different eigenvalues are \textbf{orthogonal}
\end{enumerate}

\textbf{Proof:}
Start with the eigenvalue equation:
\begin{equation}
A\ket{a'} = a'\ket{a'}
\tag{1.55}
\end{equation}

Since $A$ is Hermitian ($A = A^\dagger$):
\begin{equation}
\langle a''|A = a''^*\langle a''|
\tag{1.56}
\end{equation}

Multiply (1.55) by $\langle a''|$ on left, (1.56) by $\ket{a'}$ on right, and subtract:
\begin{equation}
(a' - a''^*)\langle a''|a'\rangle = 0
\tag{1.57}
\end{equation}

\textbf{Case 1:} If $a' = a''$, then $a' = a'^*$ (eigenvalues are real) \hfill (1.58)

\textbf{Case 2:} If $a' \neq a''$, then $\langle a''|a'\rangle = 0$ (orthogonality) \hfill (1.59)

\textbf{Orthonormal set:} Normalize the eigenkets so that $\{\ket{a'}\}$ form an orthonormal set:
\begin{equation}
\langle a''|a'\rangle = \delta_{a''a'}
\tag{1.60}
\end{equation}

\subsection{Eigenkets as Base Kets}

Normalized eigenkets of $A$ form a complete orthonormal set. Any arbitrary ket can be expanded:
\begin{equation}
\ket{\alpha} = \sum_{a'} c_{a'}\ket{a'}
\tag{1.61}
\end{equation}

\textbf{Finding expansion coefficients:}
Multiply by $\langle a''|$ on the left and use orthonormality (1.60):
\begin{equation}
c_{a'} = \langle a'|\alpha\rangle
\tag{1.62}
\end{equation}

Therefore:
\begin{equation}
\ket{\alpha} = \sum_{a'}\ket{a'}\langle a'|\alpha\rangle
\tag{1.63}
\end{equation}

\textbf{Completeness Relation (Closure):}
Since $\ket{a'}\langle a'|\alpha\rangle$ can be viewed as the operator $\ket{a'}\langle a'|$ acting on $\ket{\alpha}$:
\begin{equation}
\boxed{\sum_{a'}\ket{a'}\langle a'| = 1}
\tag{1.65}
\end{equation}
where $1$ is the identity operator. This is the \textbf{completeness relation} or \textbf{closure}.

\textbf{Application of completeness:}
Insert identity between $\langle\alpha|$ and $\ket{\alpha}$:
\begin{equation}
\langle\alpha|\alpha\rangle = \langle\alpha|\left(\sum_{a'}\ket{a'}\langle a'|\right)|\alpha\rangle = \sum_{a'}|\langle a'|\alpha\rangle|^2
\tag{1.66}
\end{equation}

For normalized $\ket{\alpha}$:
\begin{equation}
\sum_{a'}|c_{a'}|^2 = \sum_{a'}|\langle a'|\alpha\rangle|^2 = 1
\tag{1.67}
\end{equation}

\textbf{Projection Operator:}
The outer product $\ket{a'}\langle a'|$ acting on $\ket{\alpha}$ gives:
\begin{equation}
(\ket{a'}\langle a'|) \cdot \ket{\alpha} = \ket{a'}\langle a'|\alpha\rangle = c_{a'}\ket{a'}
\tag{1.68}
\end{equation}

This selects the component of $\ket{\alpha}$ parallel to $\ket{a'}$. Define the \textbf{projection operator}:
\begin{equation}
\Lambda_{a'} \equiv \ket{a'}\langle a'|
\tag{1.69}
\end{equation}

Completeness relation becomes:
\begin{equation}
\sum_{a'}\Lambda_{a'} = 1
\tag{1.70}
\end{equation}

\subsection{Matrix Representations}

\textbf{Operator as matrix:}
Using completeness relation twice:
\begin{equation}
X = \sum_{a''}\sum_{a'}\ket{a''}\langle a''|X|a'\rangle\langle a'|
\tag{1.71}
\end{equation}

The matrix elements $\langle a''|X|a'\rangle$ form an $N \times N$ square matrix:
\begin{equation}
\underbrace{\langle a''|}_{\text{row}} X \underbrace{|a'\rangle}_{\text{column}}
\tag{1.72}
\end{equation}

Explicitly:
\begin{equation}
X \doteq \begin{pmatrix}
\langle a^{(1)}|X|a^{(1)}\rangle & \langle a^{(1)}|X|a^{(2)}\rangle & \cdots \\
\langle a^{(2)}|X|a^{(1)}\rangle & \langle a^{(2)}|X|a^{(2)}\rangle & \cdots \\
\vdots & \vdots & \ddots
\end{pmatrix}
\tag{1.73}
\end{equation}

where $\doteq$ means "is represented by."

\textbf{Hermitian conjugate in matrix form:}
Using (1.53):
\begin{equation}
\langle a''|X|a'\rangle = \langle a'|X^\dagger|a''\rangle^*
\tag{1.74}
\end{equation}

For \textbf{Hermitian operator} $B$ ($B = B^\dagger$):
\begin{equation}
\langle a''|B|a'\rangle = \langle a'|B|a''\rangle^*
\tag{1.75}
\end{equation}

\textbf{Ket as column vector:}
\begin{equation}
\ket{\alpha} \doteq \begin{pmatrix}
\langle a^{(1)}|\alpha\rangle \\
\langle a^{(2)}|\alpha\rangle \\
\langle a^{(3)}|\alpha\rangle \\
\vdots
\end{pmatrix}
\tag{1.80}
\end{equation}

\textbf{Bra as row vector:}
\begin{equation}
\bra{\gamma} \doteq (\langle\gamma|a^{(1)}\rangle^*, \langle\gamma|a^{(2)}\rangle^*, \langle\gamma|a^{(3)}\rangle^*, \ldots)
\tag{1.83}
\end{equation}

\textbf{Inner product as matrix multiplication:}
\begin{equation}
\langle\beta|\alpha\rangle = \sum_{a'}\langle\beta|a'\rangle\langle a'|\alpha\rangle = (\langle a^{(1)}|\beta\rangle^*, \langle a^{(2)}|\beta\rangle^*, \ldots) \begin{pmatrix}
\langle a^{(1)}|\alpha\rangle \\
\langle a^{(2)}|\alpha\rangle \\
\vdots
\end{pmatrix}
\tag{1.84}
\end{equation}

\textbf{Outer product as matrix:}
\begin{equation}
\ket{\beta}\bra{\alpha} \doteq \begin{pmatrix}
\langle a^{(1)}|\beta\rangle\langle a^{(1)}|\alpha\rangle^* & \langle a^{(1)}|\beta\rangle\langle a^{(2)}|\alpha\rangle^* & \cdots \\
\langle a^{(2)}|\beta\rangle\langle a^{(1)}|\alpha\rangle^* & \langle a^{(2)}|\beta\rangle\langle a^{(2)}|\alpha\rangle^* & \cdots \\
\vdots & \vdots & \ddots
\end{pmatrix}
\tag{1.85}
\end{equation}

\textbf{Observable in its own eigenbasis:}
When eigenkets of $A$ are used as base kets:
\begin{equation}
A = \sum_{a''}\sum_{a'}\ket{a''}\langle a''|A|a'\rangle\langle a'|
\tag{1.86}
\end{equation}

But $\langle a''|A|a'\rangle = \langle a'|A|a'\rangle\delta_{a''a'} = a'\delta_{a''a'}$ (diagonal!), so:
\begin{equation}
\boxed{A = \sum_{a'}a'\ket{a'}\langle a'| = \sum_{a'}a'\Lambda_{a'}}
\tag{1.88}
\end{equation}

This is a \textbf{diagonal matrix} with eigenvalues on the diagonal.

\subsection*{*Why This Matrix Form? A Concrete Derivation}

\textbf{Goal:} Understand why $X = \sum_{a''}\sum_{a'}\ket{a''}\langle a''|X|a'\rangle\langle a'|$

\textbf{Setup:} We have:
\begin{itemize}
\item An operator $X$
\item A basis of eigenkets $\{\ket{a^{(1)}}, \ket{a^{(2)}}, \ket{a^{(3)}}, \ldots\}$
\item An arbitrary ket $\ket{\alpha}$ that we want to transform: $\ket{\gamma} = X\ket{\alpha}$
\end{itemize}

\textbf{Step 1: Expand the input ket $\ket{\alpha}$}

Using the completeness relation:
$$\ket{\alpha} = \sum_{a'}\ket{a'}\langle a'|\alpha\rangle$$

In column vector form, $\ket{\alpha}$ has components $\langle a'|\alpha\rangle$:
$$\ket{\alpha} \doteq \begin{pmatrix}
\langle a^{(1)}|\alpha\rangle \\
\langle a^{(2)}|\alpha\rangle \\
\langle a^{(3)}|\alpha\rangle \\
\vdots
\end{pmatrix}$$

\textbf{Step 2: Apply the operator}

$$\ket{\gamma} = X\ket{\alpha} = X\left(\sum_{a'}\ket{a'}\langle a'|\alpha\rangle\right) = \sum_{a'}X\ket{a'}\langle a'|\alpha\rangle$$

Now here's the key: What is $X\ket{a'}$? We need to express it in our basis!

\textbf{Step 3: Expand $X\ket{a'}$ in the basis}

Insert the completeness relation again:
$$X\ket{a'} = \left(\sum_{a''}\ket{a''}\langle a''|\right)X\ket{a'} = \sum_{a''}\ket{a''}\langle a''|X|a'\rangle$$

So $X$ transforms the basis ket $\ket{a'}$ into a linear combination of all basis kets, with coefficients $\langle a''|X|a'\rangle$.

\textbf{Step 4: Combine everything}

$$\ket{\gamma} = \sum_{a'}X\ket{a'}\langle a'|\alpha\rangle = \sum_{a'}\left(\sum_{a''}\ket{a''}\langle a''|X|a'\rangle\right)\langle a'|\alpha\rangle$$

Rearranging:
$$\ket{\gamma} = \sum_{a''}\ket{a''}\left(\sum_{a'}\langle a''|X|a'\rangle\langle a'|\alpha\rangle\right)$$

\textbf{Step 5: Interpret as matrix multiplication}

The component of $\ket{\gamma}$ along $\ket{a''}$ is:
$$\langle a''|\gamma\rangle = \sum_{a'}\langle a''|X|a'\rangle\langle a'|\alpha\rangle$$

In matrix form:
$$\begin{pmatrix}
\langle a^{(1)}|\gamma\rangle \\
\langle a^{(2)}|\gamma\rangle \\
\langle a^{(3)}|\gamma\rangle \\
\vdots
\end{pmatrix} = \begin{pmatrix}
\langle a^{(1)}|X|a^{(1)}\rangle & \langle a^{(1)}|X|a^{(2)}\rangle & \langle a^{(1)}|X|a^{(3)}\rangle & \cdots \\
\langle a^{(2)}|X|a^{(1)}\rangle & \langle a^{(2)}|X|a^{(2)}\rangle & \langle a^{(2)}|X|a^{(3)}\rangle & \cdots \\
\langle a^{(3)}|X|a^{(1)}\rangle & \langle a^{(3)}|X|a^{(2)}\rangle & \langle a^{(3)}|X|a^{(3)}\rangle & \cdots \\
\vdots & \vdots & \vdots & \ddots
\end{pmatrix} \begin{pmatrix}
\langle a^{(1)}|\alpha\rangle \\
\langle a^{(2)}|\alpha\rangle \\
\langle a^{(3)}|\alpha\rangle \\
\vdots
\end{pmatrix}$$

\textbf{The Meaning:}
\begin{itemize}
\item \textbf{Matrix element $\langle a''|X|a'\rangle$} = "How much of $\ket{a''}$ do you get when $X$ acts on $\ket{a'}$?"
\item \textbf{Row index ($a''$)}: Which output basis ket
\item \textbf{Column index ($a'$)}: Which input basis ket
\item The matrix tells you: "To transform any ket, apply these coefficients to each basis component"
\end{itemize}

\textbf{Concrete Example: Spin-1/2}

For spin operators with basis $\{\ket{+}, \ket{-}\}$:

If $S_x\ket{+} = \frac{\hbar}{2}\ket{-}$, then:
$$\langle +|S_x|+\rangle = 0, \quad \langle -|S_x|+\rangle = \frac{\hbar}{2}$$

These become the first column of the $S_x$ matrix!

$$S_x \doteq \begin{pmatrix}
\langle +|S_x|+\rangle & \langle +|S_x|-\rangle \\
\langle -|S_x|+\rangle & \langle -|S_x|-\rangle
\end{pmatrix} = \frac{\hbar}{2}\begin{pmatrix}
0 & 1 \\
1 & 0
\end{pmatrix}$$

\subsection{Spin-$\frac{1}{2}$ Systems}

\textbf{Base kets:} For spin-$\frac{1}{2}$ systems, use basis $\{\ket{S_z; +}, \ket{S_z; -}\}$, abbreviated as $\{\ket{+}, \ket{-}\}$

\textbf{Identity operator:}
\begin{equation}
1 = \ket{+}\langle +| + \ket{-}\langle -|
\tag{1.89}
\end{equation}

\textbf{Operator $S_z$ in terms of projection operators:}
Using (1.88):
\begin{equation}
S_z = (\hbar/2)[(\ket{+}\langle +|) - (\ket{-}\langle -|)]
\tag{1.90}
\end{equation}

The eigenvalue-eigenket relation:
\begin{equation}
S_z\ket{\pm} = \pm(\hbar/2)\ket{\pm}
\tag{1.91}
\end{equation}

\textbf{Ladder operators (Raising and Lowering):}
\begin{equation}
S_+ \equiv \hbar\ket{+}\langle -|, \quad S_- \equiv \hbar\ket{-}\langle +|
\tag{1.92}
\end{equation}

\textbf{Physical interpretation:}
\begin{itemize}
\item $S_+$ raises spin by one unit of $\hbar$: $S_+\ket{-} = \hbar\ket{+}$, but $S_+\ket{+} = 0$ (null ket)
\item $S_-$ lowers spin by one unit of $\hbar$: $S_-\ket{+} = \hbar\ket{-}$, but $S_-\ket{-} = 0$ (null ket)
\item Both are \textbf{non-Hermitian}
\item Later: $S_\pm = S_x \pm iS_y$
\end{itemize}

\textbf{Matrix representations:}
Convention: Label rows/columns in \textit{descending} order of angular momentum (max first).

\textbf{Basis kets as column vectors:}
\begin{equation}
\ket{+} \doteq \begin{pmatrix} 1 \\ 0 \end{pmatrix}, \quad \ket{-} \doteq \begin{pmatrix} 0 \\ 1 \end{pmatrix}
\tag{1.93a}
\end{equation}

\textbf{Spin operators as matrices:}
\begin{equation}
S_z \doteq \frac{\hbar}{2}\begin{pmatrix} 1 & 0 \\ 0 & -1 \end{pmatrix}, \quad 
S_+ \doteq \hbar\begin{pmatrix} 0 & 1 \\ 0 & 0 \end{pmatrix}, \quad 
S_- \doteq \hbar\begin{pmatrix} 0 & 0 \\ 1 & 0 \end{pmatrix}
\tag{1.93b}
\end{equation}

\textbf{Verification:} Check $S_+\ket{-} = \hbar\ket{+}$ using matrices:
$$S_+\ket{-} = \hbar\begin{pmatrix} 0 & 1 \\ 0 & 0 \end{pmatrix}\begin{pmatrix} 0 \\ 1 \end{pmatrix} = \hbar\begin{pmatrix} 1 \\ 0 \end{pmatrix} = \hbar\ket{+} \quad \checkmark$$
% ============================================
\newpage
\section{Tensor Products in Quantum Mechanics}

\subsection{Definition}

For Hilbert spaces $\mathcal{H}_A$ (dim $d_A$) and $\mathcal{H}_B$ (dim $d_B$):

$$\mathcal{H}_A \otimes \mathcal{H}_B \quad \text{(dimension: } d_A \times d_B\text{)}$$

\textbf{Basis:} If $\{|i\rangle_A\}$ basis for $\mathcal{H}_A$ and $\{|j\rangle_B\}$ basis for $\mathcal{H}_B$:

$$\{|i\rangle_A \otimes |j\rangle_B\} \text{ is basis for } \mathcal{H}_A \otimes \mathcal{H}_B$$

\subsection{Properties}

\textbf{Bilinearity:}
$$(a|v_1\rangle + b|v_2\rangle) \otimes |w\rangle = a(|v_1\rangle \otimes |w\rangle) + b(|v_2\rangle \otimes |w\rangle)$$

\textbf{Inner product:}
$$\langle v_1| \otimes \langle w_1| \cdot |v_2\rangle \otimes |w_2\rangle = (\langle v_1|v_2\rangle)(\langle w_1|w_2\rangle)$$

\textbf{Operator action:}
$$(\hat{A} \otimes \hat{B})(|v\rangle \otimes |w\rangle) = (\hat{A}|v\rangle) \otimes (\hat{B}|w\rangle)$$

\subsection*{Notation}

$$|v\rangle_A \otimes |w\rangle_B = |v\rangle_A|w\rangle_B = |v, w\rangle = |vw\rangle$$

\subsection{Two Spin-1/2 Example}

\textbf{Basis states:}
$$\{|++\rangle, |+-\rangle, |-+\rangle, |--\rangle\}$$

\textbf{General state:}
$$|\psi\rangle = c_1|++\rangle + c_2|+-\rangle + c_3|-+\rangle + c_4|--\rangle$$

\textbf{Separable (product) state:}
$$|\psi\rangle = |\phi\rangle_A \otimes |\chi\rangle_B$$

\textbf{Entangled state (Bell state):}
$$|\Phi^+\rangle = \frac{1}{\sqrt{2}}(|++\rangle + |--\rangle)$$

Cannot be written as product!

\subsection{Matrix Representation (Kronecker Product)}

$$A \otimes B = \begin{pmatrix}
a_{11}B & a_{12}B & \cdots \\
a_{21}B & a_{22}B & \cdots \\
\vdots & \vdots & \ddots
\end{pmatrix}$$

\textbf{Example:}
$$\sigma_z \otimes \sigma_z = \begin{pmatrix}
1 & 0 & 0 & 0 \\
0 & -1 & 0 & 0 \\
0 & 0 & -1 & 0 \\
0 & 0 & 0 & 1
\end{pmatrix}$$

\subsection{Key Formulas}

\textbf{Measurement of subsystem A:}
$$\hat{O}_A = \hat{O} \otimes \mathbb{1}_B$$

\textbf{Total spin operators:}
$$\hat{S}_{\text{total}} = \hat{S}_A \otimes \mathbb{1}_B + \mathbb{1}_A \otimes \hat{S}_B$$

\textbf{Joint probability (any state $|\psi\rangle_{AB}$, product or entangled):}
$$P(a, b) = |\langle a|_A\langle b|_B \cdot |\psi\rangle_{AB}|^2$$
% ============================================
\newpage
\section{Spin Operators and Pauli Matrices}

\subsection{Spin-1/2 Operators}

\subsubsection{Matrix Representations}

\textbf{Basis states in $S_z$ representation:}
\begin{equation}
|+\rangle = \begin{pmatrix} 1 \\ 0 \end{pmatrix}, \quad |-\rangle = \begin{pmatrix} 0 \\ 1 \end{pmatrix}
\tag{S.1}
\end{equation}

\textbf{Spin operators:}
\begin{equation}
\boxed{S_x = \frac{\hbar}{2}\begin{pmatrix} 0 & 1 \\ 1 & 0 \end{pmatrix}, \quad 
S_y = \frac{\hbar}{2}\begin{pmatrix} 0 & -i \\ i & 0 \end{pmatrix}, \quad 
S_z = \frac{\hbar}{2}\begin{pmatrix} 1 & 0 \\ 0 & -1 \end{pmatrix}}
\tag{S.2}
\end{equation}

\textbf{Ladder operators:}
\begin{equation}
S_+ = \hbar\begin{pmatrix} 0 & 1 \\ 0 & 0 \end{pmatrix}, \quad 
S_- = \hbar\begin{pmatrix} 0 & 0 \\ 1 & 0 \end{pmatrix}
\tag{S.3}
\end{equation}

\textbf{Relations:}
\begin{equation}
S_\pm = S_x \pm iS_y, \quad S_+ = S_-^\dagger
\tag{S.4}
\end{equation}

\subsection{Pauli Matrices}

\subsubsection{Definition}

\textbf{Pauli spin matrices:}
\begin{equation}
\boxed{\sigma_x = \begin{pmatrix} 0 & 1 \\ 1 & 0 \end{pmatrix}, \quad 
\sigma_y = \begin{pmatrix} 0 & -i \\ i & 0 \end{pmatrix}, \quad 
\sigma_z = \begin{pmatrix} 1 & 0 \\ 0 & -1 \end{pmatrix}}
\tag{S.5}
\end{equation}

\textbf{Relation to spin operators:}
\begin{equation}
\boxed{\mathbf{S} = \frac{\hbar}{2}\boldsymbol{\sigma}, \quad S_i = \frac{\hbar}{2}\sigma_i}
\tag{S.6}
\end{equation}

\subsubsection{Properties of Pauli Matrices}

\textbf{1. Hermitian:}
\begin{equation}
\sigma_i^\dagger = \sigma_i
\tag{S.7}
\end{equation}

\textbf{2. Traceless:}
\begin{equation}
\text{tr}(\sigma_i) = 0
\tag{S.8}
\end{equation}

\textbf{3. Determinant:}
\begin{equation}
\det(\sigma_i) = -1
\tag{S.9}
\end{equation}

\textbf{4. Square to identity:}
\begin{equation}
\boxed{\sigma_i^2 = I \quad \text{for } i = x, y, z}
\tag{S.10}
\end{equation}

\textbf{5. Anticommutation relations:}
\begin{equation}
\boxed{\{\sigma_i, \sigma_j\} = 2\delta_{ij}I}
\tag{S.11}
\end{equation}

\textbf{6. Commutation relations:}
\begin{equation}
\boxed{[\sigma_i, \sigma_j] = 2i\varepsilon_{ijk}\sigma_k}
\tag{S.12}
\end{equation}

where $\varepsilon_{ijk}$ is the Levi-Civita symbol.

\textbf{Explicit forms:}
\begin{align}
[\sigma_x, \sigma_y] &= 2i\sigma_z \tag{S.13a} \\
[\sigma_y, \sigma_z] &= 2i\sigma_x \tag{S.13b} \\
[\sigma_z, \sigma_x] &= 2i\sigma_y \tag{S.13c}
\end{align}

\textbf{7. Product formula:}
\begin{equation}
\boxed{\sigma_i\sigma_j = \delta_{ij}I + i\varepsilon_{ijk}\sigma_k}
\tag{S.14}
\end{equation}

\textbf{Special cases:}
\begin{align}
\sigma_x\sigma_y &= i\sigma_z \tag{S.15a} \\
\sigma_y\sigma_z &= i\sigma_x \tag{S.15b} \\
\sigma_z\sigma_x &= i\sigma_y \tag{S.15c}
\end{align}

\subsection{Spin Commutation Relations}

\textbf{Fundamental commutators:}
\begin{equation}
\boxed{[S_i, S_j] = i\hbar\varepsilon_{ijk}S_k}
\tag{S.16}
\end{equation}

\textbf{Explicit:}
\begin{align}
[S_x, S_y] &= i\hbar S_z \tag{S.17a} \\
[S_y, S_z] &= i\hbar S_x \tag{S.17b} \\
[S_z, S_x] &= i\hbar S_y \tag{S.17c}
\end{align}

\textbf{Total spin operator:}
\begin{equation}
\mathbf{S}^2 = S_x^2 + S_y^2 + S_z^2 = \frac{3\hbar^2}{4}I
\tag{S.18}
\end{equation}

\textbf{Commutes with all components:}
\begin{equation}
[\mathbf{S}^2, S_i] = 0 \quad \text{for all } i
\tag{S.19}
\end{equation}

\subsection{Spin in Arbitrary Direction}

\subsubsection{Direction Specified by Angles}

\textbf{Unit vector in direction $(\alpha, \beta)$:}

From the figure, with $\alpha$ = azimuthal angle (from $x$-axis in $xy$-plane) and $\beta$ = polar angle (from $z$-axis):

\begin{equation}
\boxed{\hat{\mathbf{n}} = \sin\beta\cos\alpha\,\hat{\mathbf{x}} + \sin\beta\sin\alpha\,\hat{\mathbf{y}} + \cos\beta\,\hat{\mathbf{z}}}
\tag{S.20}
\end{equation}

or in components:
\begin{equation}
\hat{\mathbf{n}} = (n_x, n_y, n_z) = (\sin\beta\cos\alpha, \sin\beta\sin\alpha, \cos\beta)
\tag{S.21}
\end{equation}

\textbf{Alternative notation (common in physics):}
Often written with $\theta$ = polar angle, $\phi$ = azimuthal angle:
\begin{equation}
\hat{\mathbf{n}} = (\sin\theta\cos\phi, \sin\theta\sin\phi, \cos\theta)
\tag{S.22}
\end{equation}

\subsubsection{Spin Operator Along $\hat{\mathbf{n}}$}

\textbf{General formula:}
\begin{equation}
\boxed{\mathbf{S} \cdot \hat{\mathbf{n}} = S_x n_x + S_y n_y + S_z n_z}
\tag{S.23}
\end{equation}

\textbf{In terms of Pauli matrices:}
\begin{equation}
\boxed{\mathbf{S} \cdot \hat{\mathbf{n}} = \frac{\hbar}{2}(\boldsymbol{\sigma} \cdot \hat{\mathbf{n}}) = \frac{\hbar}{2}(n_x\sigma_x + n_y\sigma_y + n_z\sigma_z)}
\tag{S.24}
\end{equation}

\textbf{Matrix form using angles $(\alpha, \beta)$:}
\begin{equation}
\boxed{\mathbf{S} \cdot \hat{\mathbf{n}} = \frac{\hbar}{2}\begin{pmatrix}
\cos\beta & \sin\beta e^{-i\alpha} \\
\sin\beta e^{i\alpha} & -\cos\beta
\end{pmatrix}}
\tag{S.25}
\end{equation}

\textbf{Alternative form:}
\begin{equation}
\boldsymbol{\sigma} \cdot \hat{\mathbf{n}} = \begin{pmatrix}
\cos\beta & \sin\beta e^{-i\alpha} \\
\sin\beta e^{i\alpha} & -\cos\beta
\end{pmatrix}
\tag{S.26}
\end{equation}

\subsubsection{Eigenvalues and Eigenstates}

\textbf{Eigenvalues:}
\begin{equation}
\mathbf{S} \cdot \hat{\mathbf{n}} = \pm\frac{\hbar}{2}
\tag{S.27}
\end{equation}

(Same eigenvalues as any spin component!)

\textbf{Eigenstates (spin up along $\hat{\mathbf{n}}$):}
\begin{equation}
\boxed{|+\rangle_{\hat{\mathbf{n}}} = \cos\frac{\beta}{2}|+\rangle + e^{i\alpha}\sin\frac{\beta}{2}|-\rangle}
\tag{S.28}
\end{equation}

\textbf{Spin down along $\hat{\mathbf{n}}$:}
\begin{equation}
\boxed{|-\rangle_{\hat{\mathbf{n}}} = \sin\frac{\beta}{2}|+\rangle - e^{i\alpha}\cos\frac{\beta}{2}|-\rangle}
\tag{S.29}
\end{equation}

\textbf{In column vector form:}
\begin{equation}
|+\rangle_{\hat{\mathbf{n}}} = \begin{pmatrix} \cos\frac{\beta}{2} \\ e^{i\alpha}\sin\frac{\beta}{2} \end{pmatrix}, \quad
|-\rangle_{\hat{\mathbf{n}}} = \begin{pmatrix} \sin\frac{\beta}{2} \\ -e^{i\alpha}\cos\frac{\beta}{2} \end{pmatrix}
\tag{S.30}
\end{equation}

\subsubsection{Verification}

\textbf{Orthonormality:}
\begin{equation}
{}_{\hat{\mathbf{n}}}\langle +|+\rangle_{\hat{\mathbf{n}}} = 1, \quad {}_{\hat{\mathbf{n}}}\langle -|-\rangle_{\hat{\mathbf{n}}} = 1, \quad {}_{\hat{\mathbf{n}}}\langle +|-\rangle_{\hat{\mathbf{n}}} = 0
\tag{S.31}
\end{equation}

\textbf{Eigenvalue equations:}
\begin{align}
(\mathbf{S} \cdot \hat{\mathbf{n}})|+\rangle_{\hat{\mathbf{n}}} &= +\frac{\hbar}{2}|+\rangle_{\hat{\mathbf{n}}} \tag{S.32a} \\
(\mathbf{S} \cdot \hat{\mathbf{n}})|-\rangle_{\hat{\mathbf{n}}} &= -\frac{\hbar}{2}|-\rangle_{\hat{\mathbf{n}}} \tag{S.32b}
\end{align}

\subsection{Special Cases}

\textbf{Along $z$-axis:} $\beta = 0$, $\alpha$ arbitrary
\begin{equation}
|+\rangle_{\hat{\mathbf{z}}} = |+\rangle, \quad |-\rangle_{\hat{\mathbf{z}}} = |-\rangle
\tag{S.33}
\end{equation}

\textbf{Along $x$-axis:} $\beta = \pi/2$, $\alpha = 0$
\begin{equation}
|+\rangle_{\hat{\mathbf{x}}} = \frac{1}{\sqrt{2}}(|+\rangle + |-\rangle), \quad
|-\rangle_{\hat{\mathbf{x}}} = \frac{1}{\sqrt{2}}(|+\rangle - |-\rangle)
\tag{S.34}
\end{equation}

\textbf{Along $y$-axis:} $\beta = \pi/2$, $\alpha = \pi/2$
\begin{equation}
|+\rangle_{\hat{\mathbf{y}}} = \frac{1}{\sqrt{2}}(|+\rangle + i|-\rangle), \quad
|-\rangle_{\hat{\mathbf{y}}} = \frac{1}{\sqrt{2}}(|+\rangle - i|-\rangle)
\tag{S.35}
\end{equation}

\subsection{Useful Identities}

\textbf{Expectation values in $|+\rangle$ state:}
\begin{align}
\langle +|S_x|+\rangle &= 0 \tag{S.36a} \\
\langle +|S_y|+\rangle &= 0 \tag{S.36b} \\
\langle +|S_z|+\rangle &= \frac{\hbar}{2} \tag{S.36c}
\end{align}

\textbf{Transition amplitudes:}
\begin{equation}
|\langle +|+\rangle_{\hat{\mathbf{n}}}|^2 = \cos^2\frac{\beta}{2}, \quad
|\langle -|+\rangle_{\hat{\mathbf{n}}}|^2 = \sin^2\frac{\beta}{2}
\tag{S.37}
\end{equation}

\textbf{Rotation operator (about $z$-axis by angle $\phi$):}
\begin{equation}
R_z(\phi) = e^{-iS_z\phi/\hbar} = e^{-i\sigma_z\phi/2} = \begin{pmatrix}
e^{-i\phi/2} & 0 \\
0 & e^{i\phi/2}
\end{pmatrix}
\tag{S.38}
\end{equation}

\textbf{General rotation operator:}
\begin{equation}
R_{\hat{\mathbf{n}}}(\theta) = e^{-i(\mathbf{S} \cdot \hat{\mathbf{n}})\theta/\hbar} = \cos\frac{\theta}{2}I - i\sin\frac{\theta}{2}(\boldsymbol{\sigma} \cdot \hat{\mathbf{n}})
\tag{S.39}
\end{equation}

\subsection{Quick Reference}

\begin{tcolorbox}[title=Essential Spin Formulas]
\textbf{Pauli matrices:} $\sigma_i^2 = I$, $\{\sigma_i, \sigma_j\} = 2\delta_{ij}I$

\textbf{Spin-Pauli relation:} $\mathbf{S} = \frac{\hbar}{2}\boldsymbol{\sigma}$

\textbf{Commutators:} $[S_i, S_j] = i\hbar\varepsilon_{ijk}S_k$

\textbf{Spin along $\hat{\mathbf{n}}$:} $\mathbf{S} \cdot \hat{\mathbf{n}} = \frac{\hbar}{2}\begin{pmatrix}
\cos\beta & \sin\beta e^{-i\alpha} \\
\sin\beta e^{i\alpha} & -\cos\beta
\end{pmatrix}$

\textbf{Eigenstate:} $|+\rangle_{\hat{\mathbf{n}}} = \cos\frac{\beta}{2}|+\rangle + e^{i\alpha}\sin\frac{\beta}{2}|-\rangle$
\end{tcolorbox}
% ============================================
\newpage
% ============================================
\section{Measurements, Observables, and the Uncertainty Relations}

\subsection{Measurements and State Collapse}

\textbf{Dirac's Postulate:} "A measurement always causes the system to jump into an eigenstate of the dynamical variable that is being measured."

\textbf{Before measurement:} System in superposition
\begin{equation}
\ket{\alpha} = \sum_{a'} c_{a'}\ket{a'} = \sum_{a'}\ket{a'}\langle a'|\alpha\rangle
\tag{1.94}
\end{equation}

\textbf{After measurement of observable $A$:}
\begin{equation}
\ket{\alpha} \xrightarrow{A \text{ measurement}} \ket{a'}
\tag{1.95}
\end{equation}

The system "jumps into" one of the eigenstates $\ket{a'}$.

\textbf{Exception:} If already in eigenstate $\ket{a'}$:
\begin{equation}
\ket{a'} \xrightarrow{A \text{ measurement}} \ket{a'} \quad \text{(no change)}
\tag{1.96}
\end{equation}

\textbf{Key point:} Measurement \textit{usually changes the state}. When the measurement causes $\ket{\alpha}$ to change into $\ket{a'}$, we say "$A$ is measured to be $a'$" (the eigenvalue).

\subsection{Probability Postulate}

\textbf{Non-degenerate case:} The probability for the system to jump into eigenstate $\ket{a'}$ is:
\begin{equation}
\boxed{\text{Probability for } a' = |\langle a'|\alpha\rangle|^2}
\tag{1.97}
\end{equation}
provided $\ket{\alpha}$ is normalized.

\textbf{Properties:}
\begin{itemize}
\item Probability is always non-negative
\item Probabilities sum to unity: $\sum_{a'}|\langle a'|\alpha\rangle|^2 = \langle\alpha|\alpha\rangle = 1$
\item Orthogonal states are mutually exclusive: if $\langle a''|a'\rangle = 0$ with $a'' \neq a'$, then measuring eigenvalue $a'$ excludes getting $a''$
\end{itemize}

\subsection{Degenerate Case}

\textbf{Degeneracy:} When multiple linearly independent eigenstates share the same eigenvalue $a'$:
\begin{equation}
A\ket{a'; i} = a'\ket{a'; i}, \quad i = 1, 2, \ldots, g_{a'}
\end{equation}
where $g_{a'}$ is the \textbf{degeneracy} (dimension of the eigenspace).

\textbf{Example:} For orbital angular momentum $L^2$ and $L_z$:
\begin{itemize}
\item Eigenvalue of $L^2$ is $\ell(\ell+1)\hbar^2$
\item For given $\ell$, there are $(2\ell + 1)$ states with different $m$ values: $\ket{\ell, m}$ with $m = -\ell, -\ell+1, \ldots, \ell-1, \ell$
\item All these states have the same $L^2$ eigenvalue → $(2\ell+1)$-fold degeneracy
\end{itemize}

\textbf{Probability for degenerate eigenvalue:}
When eigenvalue $a'$ is $g_{a'}$-fold degenerate, the probability of measuring $a'$ is:
\begin{equation}
\boxed{\text{Probability for eigenvalue } a' = \sum_{i=1}^{g_{a'}}|\langle a'; i|\alpha\rangle|^2}
\tag{degenerate}
\end{equation}

Sum over all eigenstates with the same eigenvalue!

\textbf{State after measurement:}
In the degenerate case, measurement does NOT collapse to a unique eigenstate. Instead:
\begin{equation}
\ket{\alpha} \xrightarrow{A \text{ measurement yields } a'} \frac{\sum_{i=1}^{g_{a'}}\ket{a'; i}\langle a'; i|\alpha\rangle}{\sqrt{\sum_{i=1}^{g_{a'}}|\langle a'; i|\alpha\rangle|^2}}
\end{equation}

The state collapses to the \textbf{degenerate subspace} corresponding to eigenvalue $a'$.

\textbf{Projection operator for degenerate subspace:}
\begin{equation}
P_{a'} = \sum_{i=1}^{g_{a'}}\ket{a'; i}\langle a'; i|
\end{equation}

This projects onto the entire degenerate subspace with eigenvalue $a'$.

\textbf{Key insight:} 
\begin{itemize}
\item Non-degenerate: Measurement determines the state completely
\item Degenerate: Measurement determines only the eigenvalue; the state is in the degenerate subspace but specific superposition within that subspace remains undetermined
\item To fully specify the state, need to measure additional compatible observables that lift the degeneracy
\end{itemize}

\textbf{Example - Angular momentum:}
\begin{itemize}
\item Measuring $L^2$ alone gives eigenvalue $\ell(\ell+1)\hbar^2$ but doesn't specify which $\ket{\ell, m}$ state
\item Must also measure $L_z$ to get specific state $\ket{\ell, m}$
\item $L^2$ and $L_z$ are \textbf{compatible observables} (they commute: $[L^2, L_z] = 0$)
\end{itemize}

\subsection{Expectation Value}

\textbf{Definition:} The expectation value (average measured value) of observable $A$ in state $\ket{\alpha}$:
\begin{equation}
\boxed{\langle A \rangle \equiv \langle\alpha|A|\alpha\rangle}
\tag{1.98}
\end{equation}

Sometimes written as $\langle A \rangle_\alpha$ to emphasize the state.

\textbf{Derivation (including degeneracy):} 
\begin{align}
\langle A \rangle &= \sum_{a'} \sum_{i=1}^{g_{a'}} \langle\alpha|a'; i\rangle\langle a'; i|A|\alpha\rangle \notag \\
&= \sum_{a'} \sum_{i=1}^{g_{a'}} a' |\langle a'; i|\alpha\rangle|^2 \notag \\
&= \sum_{a'} a' \underbrace{\left(\sum_{i=1}^{g_{a'}}|\langle a'; i|\alpha\rangle|^2\right)}_{\text{probability for eigenvalue } a'}
\tag{1.99}
\end{align}

\textbf{Critical distinction:} 
\begin{itemize}
\item \textbf{Eigenvalues} of $A$: discrete values $\{\hbar/2, -\hbar/2\}$ for $S_z$
\item \textbf{Expectation value} of $A$: can be \textit{any} real number between eigenvalues (e.g., $0.273\hbar$ for $S_z$)
\end{itemize}

\subsection{Selective Measurement (Filtration)}

\textbf{Concept:} A measurement that selects only one eigenstate $\ket{a'}$ and rejects all others (like blocking unwanted spin components in Stern-Gerlach).

\textbf{Non-degenerate case:} Apply the projection operator $\Lambda_{a'}$ to $\ket{\alpha}$:
\begin{equation}
\Lambda_{a'}\ket{\alpha} = \ket{a'}\langle a'|\alpha\rangle
\tag{1.100}
\end{equation}

\textbf{Degenerate case:} Apply the projection operator for the degenerate subspace:
\begin{equation}
P_{a'}\ket{\alpha} = \sum_{i=1}^{g_{a'}}\ket{a'; i}\langle a'; i|\alpha\rangle
\end{equation}

This filters all components in the degenerate subspace with eigenvalue $a'$.

\textbf{Physical interpretation:}
\begin{itemize}
\item Input: arbitrary state $\ket{\alpha}$
\item Device selects eigenvalue $a'$ (possibly degenerate)
\item Output: state in eigenspace of $a'$ (normalized)
\item All eigenstates with $a'' \neq a'$ are blocked/rejected
\end{itemize}

\textbf{Schwinger's approach:} Developed entire QM formalism based on selective measurements, using measurement symbol $M(a')$ which is equivalent to $\Lambda_{a'}$ or $P_{a'}$.

\subsection{Spin-$\frac{1}{2}$ Systems, Once Again}

Using sequential Stern-Gerlach experiments and quantum measurement postulates, we can determine the eigenkets and operators for all spin components.

\subsubsection{Determining $S_x$ Eigenkets}

\textbf{Experimental fact:} When $S_x+$ beam passes through SGẑ apparatus, it splits equally into $|S_z; +\rangle$ and $|S_z; -\rangle$.

This means:
\begin{equation}
|\langle +|S_x; +\rangle| = |\langle -|S_x; +\rangle| = \frac{1}{\sqrt{2}}
\tag{1.101}
\end{equation}

Therefore, the $S_x+$ eigenket is:
\begin{equation}
|S_x; +\rangle = \frac{1}{\sqrt{2}}|+\rangle + \frac{1}{\sqrt{2}}e^{i\delta_1}|-\rangle
\tag{1.102}
\end{equation}

with $\delta_1$ real (overall phase immaterial; coefficient of $|+\rangle$ chosen real and positive by convention).

\textbf{Orthogonality requirement:} Since $S_x+$ and $S_x-$ are mutually exclusive:
\begin{equation}
|S_x; -\rangle = \frac{1}{\sqrt{2}}|+\rangle - \frac{1}{\sqrt{2}}e^{i\delta_1}|-\rangle
\tag{1.103}
\end{equation}

\textbf{Constructing $S_x$ operator:}
Using the spectral decomposition (1.88):
\begin{align}
S_x &= \frac{\hbar}{2}[(|S_x; +\rangle\langle S_x; +|) - (|S_x; -\rangle\langle S_x; -|)] \notag \\
&= \frac{\hbar}{2}[e^{-i\delta_1}(|+\rangle\langle -|) + e^{i\delta_1}(|-\rangle\langle +|)]
\tag{1.104}
\end{align}

Note: $S_x$ is Hermitian, as it must be.

\subsubsection{Determining $S_y$ Eigenkets}

\textbf{Similar argument for $S_y$:}
\begin{equation}
|S_y; \pm\rangle = \frac{1}{\sqrt{2}}|+\rangle \pm \frac{1}{\sqrt{2}}e^{i\delta_2}|-\rangle
\tag{1.105}
\end{equation}

\begin{equation}
S_y = \frac{\hbar}{2}[e^{-i\delta_2}(|+\rangle\langle -|) + e^{i\delta_2}(|-\rangle\langle +|)]
\tag{1.106}
\end{equation}

\subsubsection{Determining Phase Relationships}

\textbf{Additional constraint:} Sequential SGx̂ followed by SGŷ experiment gives:
\begin{equation}
|\langle S_y; \pm|S_x; +\rangle| = |\langle S_y; \pm|S_x; -\rangle| = \frac{1}{\sqrt{2}}
\tag{1.107}
\end{equation}

(Rotational invariance of physics)

Substituting (1.103) and (1.105) into (1.107):
\begin{equation}
\frac{1}{2}|1 \pm e^{i(\delta_1-\delta_2)}| = \frac{1}{\sqrt{2}}
\tag{1.108}
\end{equation}

This is satisfied only if:
\begin{equation}
\delta_2 - \delta_1 = \frac{\pi}{2} \quad \text{or} \quad -\frac{\pi}{2}
\tag{1.109}
\end{equation}

\textbf{Key insight:} Matrix elements of $S_x$ and $S_y$ cannot both be real. If $S_x$ is real, then $S_y$ must be purely imaginary (and vice versa). Complex numbers are essential in QM!

\textbf{Convention:} Choose $\delta_1 = 0$ (making $S_x$ real) and $\delta_2 = \pi/2$ (right-handed coordinate system).

\subsubsection{Final Results}

\textbf{Eigenkets:}
\begin{align}
|S_x; \pm\rangle &= \frac{1}{\sqrt{2}}|+\rangle \pm \frac{1}{\sqrt{2}}|-\rangle \tag{1.110a} \\
|S_y; \pm\rangle &= \frac{1}{\sqrt{2}}|+\rangle \pm \frac{i}{\sqrt{2}}|-\rangle \tag{1.110b}
\end{align}

\textbf{Operators:}
\begin{align}
S_x &= \frac{\hbar}{2}[(|+\rangle\langle -|) + (|-\rangle\langle +|)] \tag{1.111a} \\
S_y &= \frac{\hbar}{2}[-i(|+\rangle\langle -|) + i(|-\rangle\langle +|)] \tag{1.111b}
\end{align}

\textbf{Ladder operators:}
\begin{equation}
\boxed{S_\pm = S_x \pm iS_y}
\tag{1.112}
\end{equation}

These agree with the earlier definitions (1.92): $S_+ = \hbar|+\rangle\langle -|$ and $S_- = \hbar|-\rangle\langle +|$.

\subsubsection{Commutation and Anticommutation Relations}

\textbf{Commutator:} (using Einstein summation convention)
\begin{equation}
\boxed{[S_i, S_j] = i\varepsilon_{ijk}\hbar S_k}
\tag{1.113}
\end{equation}

where $\varepsilon_{ijk}$ is the Levi-Civita symbol:
\begin{itemize}
\item $\varepsilon_{ijk} = +1$ for even permutations of $(1,2,3)$: $(123), (231), (312)$
\item $\varepsilon_{ijk} = -1$ for odd permutations: $(213), (132), (321)$
\item $\varepsilon_{ijk} = 0$ if any two indices are equal
\end{itemize}

These are the fundamental angular momentum commutation relations.

\textbf{Anticommutator:} (special property of spin-$\frac{1}{2}$)
\begin{equation}
\boxed{\{S_i, S_j\} = \frac{1}{2}\hbar^2\delta_{ij}}
\tag{1.114}
\end{equation}

where the anticommutator is defined as:
\begin{align}
[A, B] &\equiv AB - BA \tag{1.115a} \\
\{A, B\} &\equiv AB + BA \tag{1.115b}
\end{align}

Note: Anticommutation relation (1.114) is special to spin-$\frac{1}{2}$; does NOT hold for higher spins.

\subsubsection{Total Spin Operator}

\textbf{Definition:}
\begin{equation}
\mathbf{S}^2 \equiv S_x^2 + S_y^2 + S_z^2
\tag{1.116}
\end{equation}

Using the anticommutation relation (1.114), this becomes:
\begin{equation}
\boxed{\mathbf{S}^2 = \left(\frac{3}{4}\right)\hbar^2 = \text{constant} \times \mathbf{1}}
\tag{1.117}
\end{equation}

\textbf{Important property:}
\begin{equation}
[\mathbf{S}^2, S_i] = 0 \quad \text{for all } i
\tag{1.118}
\end{equation}

$\mathbf{S}^2$ commutes with all components of spin!

\textbf{Physical interpretation:}
\begin{itemize}
\item For spin-$\frac{1}{2}$: $\mathbf{S}^2 = \frac{3}{4}\hbar^2 = s(s+1)\hbar^2$ with $s = \frac{1}{2}$
\item This is the spin angular momentum magnitude squared
\item For higher spins, $\mathbf{S}^2$ is no longer a constant multiple of identity, but (1.118) still holds
\end{itemize}

\subsection{Compatible Observables}

\textbf{Definition:}
\begin{itemize}
\item Observables $A$ and $B$ are \textbf{compatible} when their operators commute:
\begin{equation}
[A, B] = 0
\tag{1.119}
\end{equation}

\item Observables are \textbf{incompatible} when:
\begin{equation}
[A, B] \neq 0
\tag{1.120}
\end{equation}
\end{itemize}

\textbf{Examples:}
\begin{itemize}
\item $\mathbf{S}^2$ and $S_z$ are compatible
\item $S_x$ and $S_z$ are incompatible
\end{itemize}

\subsubsection{Degeneracy}

\textbf{Degenerate eigenvalues:} When two or more linearly independent eigenkets of $A$ have the same eigenvalue $a'$, the eigenvalue is \textbf{degenerate}.

\textbf{Problem with notation $|a'\rangle$:} Labeling by eigenvalue alone doesn't give complete description when there's degeneracy.

\textbf{Solution:} Use eigenvalues of some other commuting observable $B$ to label the degenerate eigenkets.

\subsubsection{Simultaneous Eigenkets}

\textbf{Theorem 2:} Suppose $A$ and $B$ are compatible observables, and the eigenvalues of $A$ are nondegenerate. Then the matrix elements $\langle a''|B|a'\rangle$ are all diagonal.

\textbf{Proof:}
Using $[A, B] = 0$:
\begin{equation}
\langle a''|[A, B]|a'\rangle = (a'' - a')\langle a''|B|a'\rangle = 0
\tag{1.121}
\end{equation}

So $\langle a''|B|a'\rangle$ must vanish unless $a' = a''$. Therefore:
\begin{equation}
\langle a''|B|a'\rangle = \delta_{a''a'}\langle a'|B|a'\rangle
\tag{1.122}
\end{equation}

Both $A$ and $B$ can be represented by diagonal matrices with the \textit{same} set of base kets.

\textbf{Simultaneous eigenket:}
Using (1.71) and (1.122), we can write $B$ as:
\begin{equation}
B = \sum_{a''}|a''\rangle\langle a''|B|a''\rangle\langle a''|
\tag{1.123}
\end{equation}

When this acts on an eigenket of $A$:
\begin{equation}
B|a'\rangle = \sum_{a''}|a''\rangle\langle a''|B|a''\rangle\langle a''|a'\rangle = (\langle a'|B|a'\rangle)|a'\rangle
\tag{1.124}
\end{equation}

This is the eigenvalue equation for $B$ with eigenvalue:
\begin{equation}
b' \equiv \langle a'|B|a'\rangle
\tag{1.125}
\end{equation}

The ket $|a'\rangle$ is therefore a \textbf{simultaneous eigenket} of $A$ and $B$.

\textbf{Notation:} Use $|a', b'\rangle$ to characterize the simultaneous eigenket:
\begin{align}
A|a', b'\rangle &= a'|a', b'\rangle \tag{1.127a} \\
B|a', b'\rangle &= b'|a', b'\rangle \tag{1.127b}
\end{align}

\subsubsection{Degenerate Case}

When there is $n$-fold degeneracy for eigenvalue $a'$:
\begin{equation}
A|a'^{(i)}\rangle = a'|a'^{(i)}\rangle \quad \text{for } i = 1, 2, \ldots, n
\tag{1.126}
\end{equation}

All $n$ eigenkets $|a'^{(i)}\rangle$ have the same eigenvalue $a'$ but are mutually orthonormal. We can construct linear combinations of $|a'^{(i)}\rangle$ that diagonalize $B$ (see Section 1.5).

\textbf{Example - Orbital angular momentum:}
\begin{itemize}
\item Eigenvalues of $\mathbf{L}^2$: $\hbar^2 l(l+1)$
\item Eigenvalues of $L_z$: $m_l\hbar$ with $m_l = -l, -l+1, \ldots, +l$
\item Must specify \textit{both} $l$ and $m_l$ to uniquely characterize the state
\item Notation: $|l, m_l\rangle$ or $|K'\rangle = |a', b'\rangle$
\end{itemize}

\subsubsection{Maximal Set of Commuting Observables}

For multiple mutually compatible observables:
\begin{equation}
[A, B] = [B, C] = [A, C] = \cdots = 0
\tag{1.129}
\end{equation}

\textbf{Maximal set:} Cannot add any more observables without violating (1.129).

\textbf{Collective index notation:}
\begin{equation}
|K'\rangle = |a', b', c', \ldots\rangle
\tag{1.130}
\end{equation}

\textbf{Orthonormality:}
\begin{equation}
\langle K''|K'\rangle = \delta_{K'K''} = \delta_{aa'}\delta_{bb'}\delta_{cc'}\ldots
\tag{1.131}
\end{equation}

\textbf{Completeness relation:}
\begin{equation}
\sum_{K'}|K'\rangle\langle K'| = \sum_{a'}\sum_{b'}\sum_{c'}\cdots |a', b', c', \ldots\rangle\langle a', b', c', \ldots| = 1
\tag{1.132}
\end{equation}

\subsubsection{Sequential Measurements}

\textbf{Compatible observables:} When $A$ and $B$ are compatible, measurements do not interfere.

\textbf{Non-degenerate case:}
\begin{equation}
|\alpha\rangle \xrightarrow{A \text{ measurement}} |a', b'\rangle \xrightarrow{B \text{ measurement}} |a', b'\rangle \xrightarrow{A \text{ measurement}} |a', b'\rangle
\tag{1.133}
\end{equation}

The third measurement always gives $a'$ with certainty. The second ($B$) measurement does not destroy the information from the first ($A$) measurement.

\textbf{Degenerate case:}
After first ($A$) measurement yields $a'$, system is in some linear combination:
\begin{equation}
\sum_{i=1}^n c_{a'}^{(i)}|a', b^{(i)}\rangle
\tag{1.134}
\end{equation}

where all kets have the same eigenvalue $a'$ for operator $A$. The second ($B$) measurement selects one term, say $|a', b^{(j)}\rangle$, but the third ($A$) measurement still yields $a'$.

\textbf{Key insight:} $A$ measurements and $B$ measurements do not interfere. The term "compatible" is appropriate.

\textbf{Contrast with incompatible observables:}
If $[A, B] \neq 0$, measuring $B$ will generally change the eigenstate of $A$, and a subsequent measurement of $A$ will give different results. Sequential measurements interfere with each other.

\subsection{Incompatible Observables}

\textbf{Key difference:} Incompatible observables do NOT have a complete set of simultaneous eigenkets.

\subsubsection{Proof: No Simultaneous Eigenkets}

Suppose (by contradiction) that simultaneous eigenkets $|a', b'\rangle$ exist for incompatible observables $A$ and $B$. Then:
\begin{align}
AB|a', b'\rangle &= Ab'|a', b'\rangle = a'b'|a', b'\rangle \tag{1.135} \\
BA|a', b'\rangle &= Ba'|a', b'\rangle = a'b'|a', b'\rangle \tag{1.136}
\end{align}

Hence:
\begin{equation}
AB|a', b'\rangle = BA|a', b'\rangle
\tag{1.137}
\end{equation}

This implies $[A, B] = 0$, contradicting the assumption that $A$ and $B$ are incompatible.

\textbf{Conclusion:} In general, $|a', b'\rangle$ does not make sense for incompatible observables.

\textbf{Exception:} There may exist a subspace where (1.137) holds even for incompatible operators.
\begin{itemize}
\item \textbf{Example:} $l = 0$ state (s-state) in orbital angular momentum
\item Even though $L_x$ and $L_z$ do NOT commute, the $l = 0$ state IS a simultaneous eigenstate of both (with eigenvalue zero)
\item This subspace is one-dimensional
\end{itemize}

\subsubsection{Sequential Measurements: The Heart of Quantum Mechanics}

Consider sequential selective measurements: $A$ filter → $B$ filter → $C$ filter (Figure 1.8a).

\textbf{Setup:}
\begin{itemize}
\item First filter selects $|a'\rangle$
\item Second filter selects $|b'\rangle$
\item Third filter selects $|c'\rangle$
\end{itemize}

\textbf{Case 1: With $B$ measurement (Figure 1.8a)}

Probability for $|c'\rangle$ when beam from first filter is normalized:
\begin{equation}
|\langle c'|b'\rangle|^2|\langle b'|a'\rangle|^2
\tag{1.138}
\end{equation}

Sum over all possible $b'$ routes (measure each $b'$ separately):
\begin{equation}
\sum_{b'}|\langle c'|b'\rangle|^2|\langle b'|a'\rangle|^2 = \sum_{b'}\langle c'|b'\rangle\langle b'|a'\rangle\langle a'|b'\rangle\langle b'|c'\rangle
\tag{1.139}
\end{equation}

\textbf{Case 2: Without $B$ measurement (Figure 1.8b)}

When $B$ filter is absent, the pure $|a'\rangle$ beam can be expanded:
\begin{equation}
|a'\rangle = \sum_{b'}|b'\rangle\langle b'|a'\rangle
\tag{1.141}
\end{equation}

Probability for $|c'\rangle$:
\begin{equation}
|\langle c'|a'\rangle|^2 = |\sum_{b'}\langle c'|b'\rangle\langle b'|a'\rangle|^2 = \sum_{b'}\sum_{b''}\langle c'|b'\rangle\langle b'|a'\rangle\langle a'|b''\rangle\langle b''|c'\rangle
\tag{1.140}
\end{equation}

\textbf{KEY OBSERVATION:} Expressions (1.139) and (1.140) are DIFFERENT!

\begin{equation}
\boxed{\sum_{b'}|\langle c'|b'\rangle|^2|\langle b'|a'\rangle|^2 \neq |\sum_{b'}\langle c'|b'\rangle\langle b'|a'\rangle|^2}
\end{equation}

\textbf{Physical interpretation:}
\begin{itemize}
\item Both cases: pure $|a'\rangle$ beam = superposition of $B$ eigenkets (1.141)
\item \textbf{With $B$ measurement:} We experimentally ascertain which $b'$ is realized
\item \textbf{Without $B$ measurement:} We merely imagine $|a'\rangle$ is built from various $|b'\rangle$
\item \textbf{The heart of QM:} Actually recording/measuring probabilities makes all the difference, even though we sum over $b'$ afterwards
\end{itemize}

\textbf{When are the two expressions equal?}

The peculiarity disappears when:
\begin{equation}
[A, B] = 0 \quad \text{or} \quad [B, C] = 0
\tag{1.142}
\end{equation}

In other words, this phenomenon is characteristic of \textbf{incompatible observables}.

\textbf{Mathematical distinction:}
\begin{itemize}
\item Sum of squared amplitudes: $\sum_{b'}|\langle c'|b'\rangle\langle b'|a'\rangle|^2$ (probabilities add)
\item Squared sum of amplitudes: $|\sum_{b'}\langle c'|b'\rangle\langle b'|a'\rangle|^2$ (amplitudes add, then square)
\item These are equal only when there's no quantum interference
\end{itemize}

\begin{center}
\textit{[Diagram: Figure 1.8 shows (a) A→B→C sequential filters with intermediate $|b'\rangle$ states, and (b) A→C filters with expansion $|a'\rangle = \sum_{b'}|b'\rangle\langle b'|a'\rangle$ shown between them]}
\end{center}

\textbf{Bottom line:} For incompatible observables, whether or not you actually perform an intermediate measurement fundamentally changes the outcome. This is not just about "disturbing" the system—it's about the quantum nature of measurement itself.

\subsection{The Uncertainty Relation}

\subsubsection{Dispersion and Uncertainty}

\textbf{Operator deviation from mean:}
\begin{equation}
\Delta A \equiv A - \langle A \rangle
\tag{1.143}
\end{equation}

where the expectation value is taken for a certain physical state.

\textbf{Dispersion (variance, mean square deviation):}
\begin{equation}
\langle(\Delta A)^2\rangle = \langle(A^2 - 2A\langle A \rangle + \langle A \rangle^2)\rangle = \langle A^2 \rangle - \langle A \rangle^2
\tag{1.144}
\end{equation}

\textbf{Physical interpretation:}
\begin{itemize}
\item Dispersion vanishes when the state is an eigenstate of $A$
\item Dispersion characterizes "fuzziness" of measurements
\item Example: For $S_z+$ state of spin-$\frac{1}{2}$:
\begin{equation}
\langle S_x^2 \rangle - \langle S_x \rangle^2 = \hbar^2/4
\tag{1.145}
\end{equation}
\item For $S_z+$ state: $S_z$ is "sharp" (zero dispersion), $S_x$ is "fuzzy" (non-zero dispersion)
\end{itemize}

\subsubsection{The Generalized Uncertainty Relation}

\textbf{Theorem:} For any observables $A$ and $B$, and any state:
\begin{equation}
\boxed{\langle(\Delta A)^2\rangle\langle(\Delta B)^2\rangle \geq \frac{1}{4}|\langle[A,B]\rangle|^2}
\tag{1.146}
\end{equation}

This is the generalization of the Heisenberg $x$-$p$ uncertainty relation.

\subsubsection{Proof: Three Lemmas}

\textbf{Lemma 1 - Schwarz inequality:}
\begin{equation}
\langle\alpha|\alpha\rangle\langle\beta|\beta\rangle \geq |\langle\alpha|\beta\rangle|^2
\tag{1.147}
\end{equation}

Analogous to $|\mathbf{a}|^2|\mathbf{b}|^2 \geq |\mathbf{a} \cdot \mathbf{b}|^2$ in Euclidean space.

\textbf{Proof of Lemma 1:}
Consider:
\begin{equation}
(\langle\alpha| + \lambda^*\langle\beta|) \cdot (|\alpha\rangle + \lambda|\beta\rangle) \geq 0
\tag{1.149}
\end{equation}

Setting $\lambda = -\langle\beta|\alpha\rangle/\langle\beta|\beta\rangle$:
\begin{equation}
\langle\alpha|\alpha\rangle\langle\beta|\beta\rangle - |\langle\alpha|\beta\rangle|^2 \geq 0
\tag{1.150}
\end{equation}

\textbf{Lemma 2:} The expectation value of a Hermitian operator is purely real.

\textbf{Proof:} Use (1.75): $\langle a''|B|a'\rangle = \langle a'|B|a''\rangle^*$. $\square$

\textbf{Lemma 3:} The expectation value of an anti-Hermitian operator, defined by $C = -C^\dagger$, is purely imaginary.

\textbf{Proof:} Trivial. $\square$

\subsubsection{Proof of Uncertainty Relation}

Apply Lemma 1 with:
\begin{equation}
|\alpha\rangle = \Delta A|~\rangle, \quad |\beta\rangle = \Delta B|~\rangle
\tag{1.151}
\end{equation}

where the blank ket $|~\rangle$ emphasizes this applies to any ket. We obtain:
\begin{equation}
\langle(\Delta A)^2\rangle\langle(\Delta B)^2\rangle \geq |\langle\Delta A\Delta B\rangle|^2
\tag{1.152}
\end{equation}

\textbf{Decompose into commutator and anticommutator:}
\begin{equation}
\Delta A\Delta B = \frac{1}{2}[\Delta A, \Delta B] + \frac{1}{2}\{\Delta A, \Delta B\}
\tag{1.153}
\end{equation}

\textbf{Properties:}
\begin{itemize}
\item Commutator $[\Delta A, \Delta B] = [A, B]$ is \textbf{anti-Hermitian}:
\begin{equation}
([A,B])^\dagger = (AB - BA)^\dagger = BA - AB = -[A,B]
\tag{1.154}
\end{equation}

\item Anticommutator $\{\Delta A, \Delta B\}$ is \textbf{Hermitian}
\end{itemize}

Therefore:
\begin{equation}
\langle\Delta A\Delta B\rangle = \frac{1}{2}\underbrace{\langle[A,B]\rangle}_{\text{purely imaginary}} + \frac{1}{2}\underbrace{\langle\{\Delta A, \Delta B\}\rangle}_{\text{purely real}}
\tag{1.155}
\end{equation}

The right-hand side of (1.152) becomes:
\begin{equation}
|\langle\Delta A\Delta B\rangle|^2 = \frac{1}{4}|\langle[A,B]\rangle|^2 + \frac{1}{4}|\langle\{\Delta A, \Delta B\}\rangle|^2
\tag{1.156}
\end{equation}

\textbf{Final step:}
Omitting the second (anticommutator) term only makes the inequality stronger. Therefore:
\begin{equation}
\boxed{\langle(\Delta A)^2\rangle\langle(\Delta B)^2\rangle \geq \frac{1}{4}|\langle[A,B]\rangle|^2}
\end{equation}

This completes the proof. $\square$

\subsubsection{Physical Interpretation}

\textbf{Key insights:}
\begin{itemize}
\item If $[A, B] = 0$ (compatible), then uncertainty relation gives no constraint
\item If $[A, B] \neq 0$ (incompatible), then $\Delta A$ and $\Delta B$ cannot both be zero
\item The product of uncertainties has a lower bound proportional to the expectation value of the commutator
\item Example: For position and momentum, $[x, p] = i\hbar$ leads to the Heisenberg uncertainty principle (Section 1.6)
\end{itemize}

\textbf{Important note:} Applications to spin-$\frac{1}{2}$ systems and the fundamental $x$-$p$ commutation relation (Heisenberg uncertainty principle) will be discussed in Section 1.6.
% ============================================
% ============================================
\newpage
\section{Change of Basis}

\subsection{Transformation Operator}

\textbf{Motivation:} Two incompatible observables $A$ and $B$ have different sets of eigenbases. The ket space can be spanned by either $\{|a'\rangle\}$ or $\{|b'\rangle\}$. How are these related?

\textbf{Terminology:}
\begin{itemize}
\item \textbf{Change of basis} = \textbf{change of representation}
\item $A$ representation: basis given by eigenkets $\{|a'\rangle\}$ of $A$
\item $A$ diagonal representation: matrix of $A$ is diagonal in this basis
\end{itemize}

\textbf{Example:} For spin-$\frac{1}{2}$: can use $\{|S_x; \pm\rangle\}$ or $\{|S_z; \pm\rangle\}$ as base kets.

\subsubsection{Unitary Operator}

\textbf{Theorem 3:} Given two sets of base kets, both satisfying orthonormality and completeness, there exists a \textbf{unitary operator} $U$ such that:
\begin{equation}
|b^{(1)}\rangle = U|a^{(1)}\rangle, |b^{(2)}\rangle = U|a^{(2)}\rangle, \ldots, |b^{(N)}\rangle = U|a^{(N)}\rangle
\tag{1.157}
\end{equation}

\textbf{Definition - Unitary operator:}
\begin{align}
U^\dagger U &= 1 \tag{1.158} \\
UU^\dagger &= 1 \tag{1.159}
\end{align}

\textbf{Proof by construction:}
Define:
\begin{equation}
U = \sum_k |b^{(k)}\rangle\langle a^{(k)}|
\tag{1.160}
\end{equation}

Applying this to $|a^{(l)}\rangle$:
\begin{equation}
U|a^{(l)}\rangle = |b^{(l)}\rangle
\tag{1.161}
\end{equation}
guaranteed by orthonormality of $\{|a'\rangle\}$.

\textbf{Verify unitarity:}
\begin{equation}
U^\dagger U = \sum_k\sum_l |a^{(l)}\rangle\langle b^{(l)}|b^{(k)}\rangle\langle a^{(k)}| = \sum_k |a^{(k)}\rangle\langle a^{(k)}| = 1
\tag{1.162}
\end{equation}
using orthonormality of $\{|b'\rangle\}$ and completeness of $\{|a'\rangle\}$. Similarly for $UU^\dagger = 1$. $\square$

\subsection{Transformation Matrix}

\textbf{Matrix elements of $U$:}
\begin{equation}
\langle a^{(k)}|U|a^{(l)}\rangle = \langle a^{(k)}|b^{(l)}\rangle
\tag{1.163}
\end{equation}

Matrix elements are built from inner products of old base bras and new base kets.

\textbf{Analogy with 3D rotations:}
Rotation matrix changing unit vectors $(\hat{\mathbf{x}}, \hat{\mathbf{y}}, \hat{\mathbf{z}})$ to $(\hat{\mathbf{x}}', \hat{\mathbf{y}}', \hat{\mathbf{z}}')$:
\begin{equation}
R = \begin{pmatrix}
\hat{\mathbf{x}} \cdot \hat{\mathbf{x}}' & \hat{\mathbf{x}} \cdot \hat{\mathbf{y}}' & \hat{\mathbf{x}} \cdot \hat{\mathbf{z}}' \\
\hat{\mathbf{y}} \cdot \hat{\mathbf{x}}' & \hat{\mathbf{y}} \cdot \hat{\mathbf{y}}' & \hat{\mathbf{y}} \cdot \hat{\mathbf{z}}' \\
\hat{\mathbf{z}} \cdot \hat{\mathbf{x}}' & \hat{\mathbf{z}} \cdot \hat{\mathbf{y}}' & \hat{\mathbf{z}} \cdot \hat{\mathbf{z}}'
\end{pmatrix}
\tag{1.164}
\end{equation}

The square matrix $\langle a^{(k)}|U|a^{(l)}\rangle$ is the \textbf{transformation matrix} from $\{|a'\rangle\}$ basis to $\{|b'\rangle\}$ basis.

\subsubsection{Transforming Expansion Coefficients}

\textbf{Expansion in old basis:}
\begin{equation}
|\alpha\rangle = \sum_{a'}|a'\rangle\langle a'|\alpha\rangle
\tag{1.165}
\end{equation}

\textbf{Expansion coefficients in new basis:}
Multiply by $\langle b^{(k)}|$:
\begin{equation}
\langle b^{(k)}|\alpha\rangle = \sum_l \langle b^{(k)}|a^{(l)}\rangle\langle a^{(l)}|\alpha\rangle = \sum_l \langle a^{(k)}|U^\dagger|a^{(l)}\rangle\langle a^{(l)}|\alpha\rangle
\tag{1.166}
\end{equation}

\textbf{In matrix notation:}
\begin{equation}
\boxed{(\text{new}) = (U^\dagger)(\text{old})}
\tag{1.167}
\end{equation}

Column matrix in new basis = $U^\dagger$ times column matrix in old basis.

\subsubsection{Similarity Transformation}

\textbf{Matrix elements in new basis:}
\begin{align}
\langle b^{(k)}|X|b^{(l)}\rangle &= \sum_m\sum_n \langle b^{(k)}|a^{(m)}\rangle\langle a^{(m)}|X|a^{(n)}\rangle\langle a^{(n)}|b^{(l)}\rangle \notag \\
&= \sum_m\sum_n \langle a^{(k)}|U^\dagger|a^{(m)}\rangle\langle a^{(m)}|X|a^{(n)}\rangle\langle a^{(n)}|U|a^{(l)}\rangle
\tag{1.168}
\end{align}

This is the \textbf{similarity transformation} in matrix algebra:
\begin{equation}
\boxed{X' = U^\dagger XU}
\tag{1.169}
\end{equation}

\subsubsection{Trace Properties}

\textbf{Definition of trace:}
\begin{equation}
\text{tr}(X) = \sum_{a'}\langle a'|X|a'\rangle
\tag{1.170}
\end{equation}

\textbf{Trace is basis-independent:}
\begin{equation}
\sum_{a'}\langle a'|X|a'\rangle = \sum_{a'}\sum_{b'}\sum_{b''}\langle a'|b'\rangle\langle b'|X|b''\rangle\langle b''|a'\rangle = \sum_{b'}\langle b'|X|b'\rangle
\tag{1.171}
\end{equation}

\textbf{Useful trace properties:}
\begin{align}
\text{tr}(XY) &= \text{tr}(YX) \tag{1.172a} \\
\text{tr}(U^\dagger XU) &= \text{tr}(X) \tag{1.172b} \\
\text{tr}(|a'\rangle\langle a''|) &= \delta_{a'a''} \tag{1.172c} \\
\text{tr}(|b'\rangle\langle a'|) &= \langle a'|b'\rangle \tag{1.172d}
\end{align}

\subsection{Diagonalization}

\textbf{Problem:} Find eigenvalues and eigenkets of operator $B$ whose matrix elements in old $\{|a'\rangle\}$ basis are known.

\textbf{Solution:} This is equivalent to finding the unitary matrix that diagonalizes $B$.

We seek eigenvalue $b'$ and eigenket $|b'\rangle$ such that:
\begin{equation}
B|b'\rangle = b'|b'\rangle
\tag{1.173}
\end{equation}

\textbf{Rewrite using completeness:}
\begin{equation}
\sum_{a''}\langle a''|B|a'\rangle\langle a'|b'\rangle = b'\langle a''|b'\rangle
\tag{1.174}
\end{equation}

\textbf{In matrix form:}
\begin{equation}
\begin{pmatrix}
B_{11} & B_{12} & B_{13} & \cdots \\
B_{21} & B_{22} & B_{23} & \cdots \\
\vdots & \vdots & \vdots & \ddots
\end{pmatrix}
\begin{pmatrix}
C_1^{(l)} \\
C_2^{(l)} \\
\vdots
\end{pmatrix}
= b^{(l)}
\begin{pmatrix}
C_1^{(l)} \\
C_2^{(l)} \\
\vdots
\end{pmatrix}
\tag{1.175}
\end{equation}

where:
\begin{align}
B_{ij} &= \langle a^{(i)}|B|a^{(j)}\rangle \tag{1.176a} \\
C_k^{(l)} &= \langle a^{(k)}|b^{(l)}\rangle \tag{1.176b}
\end{align}

\textbf{Characteristic equation:}
Nontrivial solutions for $C_k^{(l)}$ exist only if:
\begin{equation}
\boxed{\det(B - \lambda 1) = 0}
\tag{1.177}
\end{equation}

This is an $N$th order algebraic equation for $\lambda$. The $N$ roots are the eigenvalues $b^{(l)}$.

\textbf{Procedure:}
\begin{enumerate}
\item Solve characteristic equation (1.177) to find eigenvalues $b^{(l)}$
\item For each $b^{(l)}$, solve (1.175) for coefficients $C_k^{(l)}$ (up to normalization)
\item The $C_k^{(l)}$ are the elements of the unitary matrix $U$ for change of basis: $\{|a'\rangle\} \to \{|b'\rangle\}$
\end{enumerate}

\textbf{Important note:} This procedure requires $B$ to be Hermitian! 

\textbf{Non-Hermitian operators:}
Example: $S_+$ ladder operator (1.92) has matrix in $S_z$ basis:
\begin{equation}
S_+ \doteq \hbar\begin{pmatrix}
0 & 1 \\
0 & 0
\end{pmatrix}
\tag{1.178}
\end{equation}

This CANNOT be diagonalized by any unitary matrix. Non-Hermitian operators (like coherent state eigenvectors in Chapter 2) do NOT form a complete orthonormal set, so the formalism developed here doesn't apply.

\subsection{Unitary Equivalent Observables}

\textbf{Theorem 4:} Consider two sets of orthonormal bases $\{|a'\rangle\}$ and $\{|b'\rangle\}$ connected by the $U$ operator (1.160). Knowing $U$, we may construct a \textbf{unitary transform} of $A$:
\begin{equation}
UAU^{-1}
\end{equation}

$A$ and $UAU^{-1}$ are said to be \textbf{unitary equivalent observables}.

\subsubsection{Key Properties}

\textbf{Starting from eigenvalue equation for $A$:}
\begin{equation}
A|a^{(l)}\rangle = a^{(l)}|a^{(l)}\rangle
\tag{1.179}
\end{equation}

\textbf{Multiply by $U$ on left and insert $1 = U^{-1}U$:}
\begin{equation}
UAU^{-1}U|a^{(l)}\rangle = a^{(l)}U|a^{(l)}\rangle
\tag{1.180}
\end{equation}

\textbf{Using $U|a^{(l)}\rangle = |b^{(l)}\rangle$:}
\begin{equation}
\boxed{(UAU^{-1})|b^{(l)}\rangle = a^{(l)}|b^{(l)}\rangle}
\tag{1.181}
\end{equation}

\textbf{Profound result:} The $|b'\rangle$ are eigenkets of $UAU^{-1}$ with \textit{exactly the same eigenvalues} as the $A$ eigenvalues.

\begin{center}
\textit{Unitary equivalent observables have identical spectra.}
\end{center}

\subsubsection{Relationship with Observable $B$}

By definition, eigenket $|b^{(l)}\rangle$ satisfies:
\begin{equation}
B|b^{(l)}\rangle = b^{(l)}|b^{(l)}\rangle
\tag{1.182}
\end{equation}

\textbf{Comparing (1.181) and (1.182):}
\begin{itemize}
\item $B$ and $UAU^{-1}$ are \textbf{simultaneously diagonalizable}
\item They share the same eigenkets $|b^{(l)}\rangle$
\item But eigenvalues may differ: $b^{(l)}$ for $B$ vs. $a^{(l)}$ for $UAU^{-1}$
\end{itemize}

\textbf{Special case:} Is $UAU^{-1}$ the same as $B$ itself?

The answer is often YES in cases of physical interest!

\subsubsection{Example: Spin Components}

\textbf{$S_x$ and $S_z$:}
\begin{itemize}
\item Related by a unitary operator: rotation operator around $y$-axis by angle $\pi/2$
\item $S_x$ itself is the unitary transform of $S_z$
\item Both have the same set of eigenvalues: $+\hbar/2$ and $-\hbar/2$
\item The theorem holds in this particular example
\end{itemize}

\textbf{Physical interpretation:}
Different spin components are related by rotations (unitary transformations). Since physical laws are rotationally invariant, all spin components must have the same spectrum.

\textbf{General principle:}
When two observables are related by a symmetry transformation (like rotation), they are unitary equivalent and share the same spectrum.

\newpage
% ============================================
% ============================================
\section*{Diagonalization: Detailed Explanation}

\subsection*{Setup: Two Different Bases}

\textbf{Original basis (arbitrary):}
$$\{|1\rangle, |2\rangle, \ldots, |N\rangle\}$$

Operator $A$ has matrix elements:
$$A_{mn} = \langle m|A|n\rangle$$

This matrix is generally \textbf{NOT diagonal}.

\textbf{Eigenbasis (special basis):}
$$\{|a^{(1)}\rangle, |a^{(2)}\rangle, \ldots, |a^{(N)}\rangle\}$$

These satisfy:
$$A|a^{(i)}\rangle = a^{(i)}|a^{(i)}\rangle$$

where $a^{(i)}$ are the eigenvalues. We also write $|a'\rangle, |a''\rangle, \ldots$ for these eigenstates.

\subsection*{What is $U$? The Transformation Matrix}

\textbf{Definition:}
$$\boxed{U = \sum_{i=1}^N |a^{(i)}\rangle\langle i|}$$

\textbf{In matrix form:}
$$U = \begin{pmatrix}
\langle 1|a^{(1)}\rangle & \langle 1|a^{(2)}\rangle & \cdots & \langle 1|a^{(N)}\rangle \\
\langle 2|a^{(1)}\rangle & \langle 2|a^{(2)}\rangle & \cdots & \langle 2|a^{(N)}\rangle \\
\vdots & \vdots & \ddots & \vdots \\
\langle N|a^{(1)}\rangle & \langle N|a^{(2)}\rangle & \cdots & \langle N|a^{(N)}\rangle
\end{pmatrix}$$

\textbf{Reading the matrix:}
\begin{itemize}
\item \textbf{Column $j$} = eigenvector $|a^{(j)}\rangle$ expressed in original basis $\{|n\rangle\}$
\item Entry $U_{ij} = \langle i|a^{(j)}\rangle$ = component of $j$-th eigenvector along $i$-th original basis vector
\end{itemize}

\subsection*{What Does $U$ Do?}

\textbf{$U$ transforms component columns FROM eigenbasis TO original basis:}

If $|\psi\rangle = \sum_i c_i^{(\text{eigen})} |a^{(i)}\rangle = \sum_n c_n |n\rangle$, then:
$$\sum_i U_{ni}\, c_i^{(\text{eigen})} = c_n, \quad \text{i.e.} \quad U c^{(\text{eigen})} = c^{(\text{orig})}$$
(As an operator on kets, $U|i\rangle = |a^{(i)}\rangle$.)

\textbf{$U^\dagger$ transforms vectors FROM original basis TO eigenbasis:}
$$\boxed{U^\dagger = \sum_{i=1}^N |i\rangle\langle a^{(i)}|}$$

If $|\psi\rangle = \sum_n c_n |n\rangle$, then:
$$U^\dagger c^{(\text{orig})} = c^{(\text{eigen})}, \quad c_i^{(\text{eigen})} = \langle a^{(i)}|\psi\rangle$$

\subsection*{Why Does $U^\dagger A U$ Diagonalize $A$?}

\textbf{The key calculation:}

Compute the matrix element in the original basis:
$$\langle m|U^\dagger A U|n\rangle$$

\textbf{Step 1:} Insert completeness relations:
$$= \sum_{a',a''} \langle m|U^\dagger|a'\rangle\langle a'|A|a''\rangle\langle a''|U|n\rangle$$

\textbf{Step 2:} Use $U = \sum_i |a^{(i)}\rangle\langle i|$, so:
$$\langle m|U^\dagger = \langle m|\sum_i |i\rangle\langle a^{(i)}| = \langle a^{(m)}|$$
$$U|n\rangle = \sum_i |a^{(i)}\rangle\langle i|n\rangle = |a^{(n)}\rangle$$

\textbf{Step 3:} The matrix element becomes:
$$= \langle a^{(m)}|A|a^{(n)}\rangle$$

\textbf{Step 4:} Use eigenvalue equation $A|a^{(n)}\rangle = a^{(n)}|a^{(n)}\rangle$:
$$= a^{(n)}\langle a^{(m)}|a^{(n)}\rangle = a^{(n)}\delta_{mn}$$

\textbf{Result:}
$$\boxed{\langle m|U^\dagger A U|n\rangle = a^{(n)}\delta_{mn}}$$

This is \textbf{diagonal}! Only diagonal entries survive.

\subsection*{The Diagonalization Formula}

$$\boxed{U^\dagger A U = \Lambda}$$

where
$$\Lambda = \begin{pmatrix}
a^{(1)} & 0 & 0 & \cdots \\
0 & a^{(2)} & 0 & \cdots \\
0 & 0 & a^{(3)} & \cdots \\
\vdots & \vdots & \vdots & \ddots
\end{pmatrix}$$

Multiplying by $U$ on left and $U^\dagger$ on right:
$$\boxed{A = U\Lambda U^\dagger}$$

Equivalently (spectral decomposition):
$$\boxed{A = \sum_{a'} a' |a'\rangle\langle a'|}$$

\subsection*{Physical Picture: Change of Perspective}

\textbf{Original basis perspective:}
\begin{itemize}
\item View operator $A$ from basis $\{|n\rangle\}$
\item $A$ mixes basis vectors: $A|n\rangle = \sum_m \langle m|A|n\rangle|m\rangle$
\item Matrix has off-diagonal elements
\end{itemize}

\textbf{Eigenbasis perspective:}
\begin{itemize}
\item View along "natural axes" of $A$ (the eigenvectors)
\item Simple action: $A|a'\rangle = a'|a'\rangle$
\item No mixing! Just multiplication
\item Matrix is diagonal
\end{itemize}

\textbf{$U$ is the rotation between perspectives:}
\begin{itemize}
\item $U$ rotates \textbf{vectors} from eigenbasis to original basis
\item $U^\dagger$ rotates \textbf{vectors} from original to eigenbasis
\item $U^\dagger A U$ rotates the \textbf{operator} into eigenbasis perspective
\end{itemize}

\subsection*{How to Compute $A|\psi\rangle$}

\textbf{Method 1: Spectral decomposition}

Since $A = \sum_{a'} a' |a'\rangle\langle a'|$:
$$A|\psi\rangle = \sum_{a'} a' |a'\rangle\langle a'|\psi\rangle$$

\textbf{Steps:}
\begin{enumerate}
\item Compute $\langle a'|\psi\rangle$ = component along eigenvector
\item Multiply each by eigenvalue $a'$
\item Sum: $A|\psi\rangle = \sum_{a'} a' \langle a'|\psi\rangle |a'\rangle$
\end{enumerate}

\textbf{Method 2: Using $U$ explicitly}

$$A|\psi\rangle = U\Lambda U^\dagger|\psi\rangle$$

\textbf{Steps:}
\begin{enumerate}
\item \textbf{Transform to eigenbasis:} $|\psi'\rangle = U^\dagger|\psi\rangle$
   
   Components: $c_{a'}' = \langle a'|\psi\rangle$

\item \textbf{Apply diagonal operator:} $|\phi'\rangle = \Lambda|\psi'\rangle$
   
   In eigenbasis: $\phi_{a'}' = a' c_{a'}'$

\item \textbf{Transform back:} $|\phi\rangle = U|\phi'\rangle$
   
   Result in original basis
\end{enumerate}

\subsection*{Why Does This Work? The Deep Reason}

\textbf{Eigenvectors are special:}

In eigenbasis, $A$ has simplest form:
$$A_{\text{eigenbasis}} = \begin{pmatrix}
a^{(1)} & 0 & \cdots \\
0 & a^{(2)} & \cdots \\
\vdots & \vdots & \ddots
\end{pmatrix}$$

\textbf{Why?} Eigenvectors are "natural coordinates" - directions $A$ doesn't rotate, only stretches.

\textbf{Similarity transformations preserve the operator:}

$$\langle\beta|A|\alpha\rangle = \langle\beta'|A'|\alpha'\rangle$$

where $|\alpha'\rangle = U^\dagger|\alpha\rangle$, $A' = U^\dagger A U$

The physics is the same, just different coordinates!

\subsection*{Example: $S_x$ in the $S_z$ Basis}

\textbf{Original basis:} $\{|1\rangle, |2\rangle\} = \{|+\rangle_z, |-\rangle_z\}$ (eigenstates of $S_z$)

In this basis:
$$S_x = \frac{\hbar}{2}\begin{pmatrix} 0 & 1 \\ 1 & 0 \end{pmatrix} \quad \text{(off-diagonal!)}$$

\textbf{Find eigenvalues:}
$$\det\begin{pmatrix} -\lambda & \hbar/2 \\ \hbar/2 & -\lambda \end{pmatrix} = \lambda^2 - \frac{\hbar^2}{4} = 0$$
$$\lambda = \pm\frac{\hbar}{2}$$

So: $a^{(1)} = +\hbar/2$, $a^{(2)} = -\hbar/2$

\textbf{Find eigenvectors:}

For $a^{(1)} = +\hbar/2$:
$$|a^{(1)}\rangle = |+\rangle_x = \frac{1}{\sqrt{2}}(|1\rangle + |2\rangle) = \frac{1}{\sqrt{2}}\begin{pmatrix} 1 \\ 1 \end{pmatrix}$$

For $a^{(2)} = -\hbar/2$:
$$|a^{(2)}\rangle = |-\rangle_x = \frac{1}{\sqrt{2}}(|1\rangle - |2\rangle) = \frac{1}{\sqrt{2}}\begin{pmatrix} 1 \\ -1 \end{pmatrix}$$

\textbf{Construct $U$:}

Columns are eigenvectors in original basis $\{|n\rangle\}$:
$$U = (|a^{(1)}\rangle, |a^{(2)}\rangle) = \frac{1}{\sqrt{2}}\begin{pmatrix} 1 & 1 \\ 1 & -1 \end{pmatrix}$$

where $U_{n,i} = \langle n|a^{(i)}\rangle$

\textbf{Verify diagonalization:}
$$U^\dagger S_x U = \frac{1}{\sqrt{2}}\begin{pmatrix} 1 & 1 \\ 1 & -1 \end{pmatrix} \frac{\hbar}{2}\begin{pmatrix} 0 & 1 \\ 1 & 0 \end{pmatrix} \frac{1}{\sqrt{2}}\begin{pmatrix} 1 & 1 \\ 1 & -1 \end{pmatrix}$$
$$= \begin{pmatrix} \hbar/2 & 0 \\ 0 & -\hbar/2 \end{pmatrix} = \Lambda \quad \checkmark$$

\textbf{What $U$ did:}
\begin{itemize}
\item \textbf{Original basis:} $\{|n\rangle\} = $ $S_z$ eigenstates where $S_x$ is off-diagonal
\item \textbf{$U$ transforms TO:} $\{|a'\rangle\} = $ $S_x$ eigenstates where $S_x$ is diagonal
\item \textbf{$U^\dagger$ transforms FROM:} $\{|n\rangle\}$ TO $\{|a'\rangle\}$
\end{itemize}

\subsection*{Summary Table}

\begin{center}
\begin{tabular}{|l|l|l|}
\hline
\textbf{Matrix} & \textbf{What It Does} & \textbf{Direction} \\
\hline
$U$ & Transforms vectors & Eigenbasis $\{|a'\rangle\}$ $\to$ Original $\{|n\rangle\}$ \\
$U^\dagger$ & Transforms vectors & Original $\{|n\rangle\}$ $\to$ Eigenbasis $\{|a'\rangle\}$ \\
$U^\dagger A U$ & Transforms operator & Makes $A$ diagonal \\
\hline
\end{tabular}
\end{center}

\textbf{Key formulas:}
$$A = U\Lambda U^\dagger = \sum_{a'} a' |a'\rangle\langle a'|$$

\textbf{Why it works:} In eigenbasis $\{|a'\rangle\}$, operator $A$ is naturally diagonal. $U$ and $U^\dagger$ change perspective between bases $\{|n\rangle\} \leftrightarrow \{|a'\rangle\}$ while preserving physics.

\textbf{Notation summary:}
\begin{itemize}
\item $A$ = arbitrary observable operator
\item $\{|n\rangle\} = \{|1\rangle, |2\rangle, \ldots, |N\rangle\}$ = original (arbitrary) basis
\item $\{|a'\rangle\} = \{|a^{(1)}\rangle, |a^{(2)}\rangle, \ldots, |a^{(N)}\rangle\}$ = eigenbasis of $A$
\item $a', a'', \ldots$ or $a^{(1)}, a^{(2)}, \ldots$ = eigenvalues of $A$
\item $U_{ni} = \langle n|a^{(i)}\rangle$ = transformation matrix
\end{itemize}

\newpage
\subsection*{Diagonalization: Derivation with Two Completeness Relations}

\textbf{Setup:}
\begin{itemize}
\item $A$ is an observable
\item $\{|n\rangle\} = \{|1\rangle, |2\rangle, \ldots\}$ is the standard basis
\item $\{|a_k\rangle\}$ are eigenstates with eigenvalues $a_k$: $A|a_k\rangle = a_k|a_k\rangle$
\item $|\psi\rangle$ is an arbitrary state
\end{itemize}

\subsection*{Starting Point: Insert Two Completeness Relations}

Begin with $A|\psi\rangle$ and insert completeness for both bases:

\textbf{Insert eigenbasis completeness on the right:}
$$A|\psi\rangle = A \left(\sum_k |a_k\rangle\langle a_k|\right) |\psi\rangle$$

\textbf{Insert standard basis completeness on the left:}
$$= \left(\sum_j |j\rangle\langle j|\right) A \left(\sum_k |a_k\rangle\langle a_k|\right) |\psi\rangle$$

\subsection*{Rearrange and Apply Eigenvalue Equation}

$$= \sum_j \sum_k |j\rangle\langle j|A|a_k\rangle\langle a_k|\psi\rangle$$

\textbf{Use the eigenvalue equation:} $A|a_k\rangle = a_k|a_k\rangle$

$$\langle j|A|a_k\rangle = \langle j|(a_k|a_k\rangle) = a_k\langle j|a_k\rangle$$

Therefore:
$$\boxed{A|\psi\rangle = \sum_j \sum_k |j\rangle \underbrace{\langle j|a_k\rangle}_{U_{jk}} \cdot \underbrace{a_k}_{\Lambda_{kk}} \cdot \underbrace{\langle a_k|\psi\rangle}_{(U^\dagger|\psi\rangle)_k}}$$

\subsection*{Identifying the Matrix Components}

\textbf{Transformation matrix $U$:}
$$U_{jk} = \langle j|a_k\rangle$$

\textbf{Meaning:} Component of eigenvector $|a_k\rangle$ along standard basis vector $|j\rangle$

Column $k$ of $U$ is eigenvector $|a_k\rangle$ expressed in standard basis:
$$U = \begin{pmatrix}
\langle 1|a_1\rangle & \langle 1|a_2\rangle & \cdots \\
\langle 2|a_1\rangle & \langle 2|a_2\rangle & \cdots \\
\vdots & \vdots & \ddots
\end{pmatrix} = (|a_1\rangle, |a_2\rangle, \ldots)$$

\textbf{Diagonal eigenvalue matrix $\Lambda$:}
$$\Lambda_{kk'} = a_k \delta_{kk'}$$

$$\Lambda = \begin{pmatrix}
a_1 & 0 & 0 & \cdots \\
0 & a_2 & 0 & \cdots \\
0 & 0 & a_3 & \cdots \\
\vdots & \vdots & \vdots & \ddots
\end{pmatrix}$$

\textbf{Transformed state $(U^\dagger|\psi\rangle)$:}
$$(U^\dagger|\psi\rangle)_k = \langle a_k|\psi\rangle$$

This is the $k$-th component of $|\psi\rangle$ in the eigenbasis.

\subsection*{Matrix Form}

The double sum can be written as matrix multiplication:

$$A|\psi\rangle = \sum_j |j\rangle \left[\sum_k U_{jk} \cdot \Lambda_{kk} \cdot (U^\dagger|\psi\rangle)_k\right]$$

\textbf{In matrix notation:}
$$\boxed{A|\psi\rangle = U \Lambda U^\dagger |\psi\rangle}$$

Since this holds for any state $|\psi\rangle$:
$$\boxed{A = U \Lambda U^\dagger}$$

\subsection*{Why This is Diagonalization}

\textbf{Step-by-step action:}

\textbf{Step 1:} $\langle a_k|\psi\rangle$ projects $|\psi\rangle$ onto the $k$-th eigenvector
$$\text{Component in eigenbasis: } c_k = \langle a_k|\psi\rangle$$

\textbf{Step 2:} $a_k \langle a_k|\psi\rangle$ scales by eigenvalue
$$\text{Scaled component: } a_k c_k$$

\textbf{Step 3:} $\sum_k \langle j|a_k\rangle [a_k \langle a_k|\psi\rangle]$ reconstructs in standard basis
$$\text{Result in standard basis: } \sum_k U_{jk} (a_k c_k)$$

\subsection*{The Diagonal Matrix Explicitly}

To see $\Lambda$ explicitly, write the matrix multiplication:

$$(U \Lambda U^\dagger)_{jm} = \sum_k \sum_{k'} U_{jk} \Lambda_{kk'} (U^\dagger)_{k'm}$$

Using $\Lambda_{kk'} = a_k \delta_{kk'}$:
$$= \sum_k U_{jk} a_k (U^\dagger)_{km} = \sum_k \langle j|a_k\rangle a_k \langle a_k|m\rangle$$

Therefore:
$$A = \sum_k |a_k\rangle a_k \langle a_k| = \sum_k a_k |a_k\rangle\langle a_k|$$

This is the \textbf{spectral decomposition}.

\subsection*{Verifying $U^\dagger$}

\textbf{Definition:} $U = \sum_k |a_k\rangle\langle k|$ (in bra-ket notation)

\textbf{Hermitian conjugate:}
$$U^\dagger = \left(\sum_k |a_k\rangle\langle k|\right)^\dagger = \sum_k |k\rangle\langle a_k|$$

\textbf{Matrix elements:}
$$(U^\dagger)_{km} = \langle k|U^\dagger|m\rangle = \sum_{k'} \langle k|k'\rangle\langle a_{k'}|m\rangle = \langle a_k|m\rangle$$

\textbf{Action on $|\psi\rangle$:}
$$(U^\dagger|\psi\rangle)_k = \sum_m (U^\dagger)_{km}\langle m|\psi\rangle = \sum_m \langle a_k|m\rangle\langle m|\psi\rangle = \langle a_k|\psi\rangle$$

Using completeness $\sum_m |m\rangle\langle m| = \mathbb{1}$.

\subsection*{Summary of Two-Index Form}

\textbf{Starting expression:}
$$A|\psi\rangle = \left(\sum_j |j\rangle\langle j|\right) A \left(\sum_k |a_k\rangle\langle a_k|\right) |\psi\rangle$$

\textbf{After using eigenvalue equation:}
$$\boxed{A|\psi\rangle = \sum_j \sum_k |j\rangle \langle j|a_k\rangle \cdot a_k \cdot \langle a_k|\psi\rangle}$$

\textbf{Matrix components:}
\begin{itemize}
\item $U_{jk} = \langle j|a_k\rangle$ (columns are eigenvectors)
\item $\Lambda_{kk} = a_k$ (diagonal eigenvalues)
\item $(U^\dagger|\psi\rangle)_k = \langle a_k|\psi\rangle$ (state in eigenbasis)
\end{itemize}

\textbf{Final form:}
$$\boxed{A = U\Lambda U^\dagger = \sum_j \sum_k |j\rangle\langle j|a_k\rangle \cdot a_k \cdot \langle a_k|}$$

\subsection*{Visual Representation}

$$\begin{array}{ccccc}
|\psi\rangle_{\text{std}} & \xrightarrow{U^\dagger} & \begin{pmatrix} \langle a_1|\psi\rangle \\ \langle a_2|\psi\rangle \\ \vdots \end{pmatrix}_{\text{eigen}} & \xrightarrow{\Lambda} & \begin{pmatrix} a_1\langle a_1|\psi\rangle \\ a_2\langle a_2|\psi\rangle \\ \vdots \end{pmatrix}_{\text{eigen}} & \xrightarrow{U} & A|\psi\rangle_{\text{std}}
\end{array}$$

\textbf{Key insight:} Only \textbf{two completeness relations} needed:
\begin{enumerate}
\item $\sum_j |j\rangle\langle j| = \mathbb{1}$ (standard basis)
\item $\sum_k |a_k\rangle\langle a_k| = \mathbb{1}$ (eigenbasis)
\end{enumerate}

This naturally gives the two-index form with $j$ (standard basis) and $k$ (eigenbasis).

\subsection*{Why Diagonal in Eigenbasis?}

In the eigenbasis itself, the matrix of $A$ is:
$$\langle a_{k'}|A|a_k\rangle = a_k \langle a_{k'}|a_k\rangle = a_k \delta_{k'k}$$

This is diagonal! The transformation $U$ rotates from standard basis (where $A$ may not be diagonal) to eigenbasis (where $A$ is diagonal).
% ============================================
\newpage
\section{Position, Momentum, and Translation}

\subsection{Continuous Spectra}

\textbf{Continuous eigenvalues:} Some observables (like position $x$ or momentum $p_z$) have continuous eigenvalues ranging from $-\infty$ to $\infty$.

\textbf{Eigenvalue equation for continuous spectrum:}
\begin{equation}
\xi|\xi'\rangle = \xi'|\xi'\rangle
\tag{1.183}
\end{equation}

where $\xi$ is an operator and $\xi'$ is simply a number (the eigenvalue).

\textbf{Key modifications for continuous spectra:}

Replace Kronecker delta by Dirac delta function, discrete sums by integrals:

\begin{align}
\langle a'|a''\rangle = \delta_{a'a''} &\to \langle\xi'|\xi''\rangle = \delta(\xi' - \xi'') \tag{1.184a} \\
\sum_{a'}|a'\rangle\langle a'| = 1 &\to \int d\xi'|\xi'\rangle\langle\xi'| = 1 \tag{1.184b} \\
|\alpha\rangle = \sum_{a'}|a'\rangle\langle a'|\alpha\rangle &\to |\alpha\rangle = \int d\xi'|\xi'\rangle\langle\xi'|\alpha\rangle \tag{1.184c} \\
\sum_{a'}|\langle a'|\alpha\rangle|^2 = 1 &\to \int d\xi'|\langle\xi'|\alpha\rangle|^2 = 1 \tag{1.184d} \\
\langle\beta|\alpha\rangle = \sum_{a'}\langle\beta|a'\rangle\langle a'|\alpha\rangle &\to \langle\beta|\alpha\rangle = \int d\xi'\langle\beta|\xi'\rangle\langle\xi'|\alpha\rangle \tag{1.184e} \\
\langle a''|A|a'\rangle = a'\delta_{a''a'} &\to \langle\xi''|\xi|\xi'\rangle = \xi'\delta(\xi'' - \xi') \tag{1.184f}
\end{align}

\subsection{Position Eigenkets and Position Measurements}

\textbf{Eigenkets of position operator:}
\begin{equation}
x|x'\rangle = x'|x'\rangle
\tag{1.185}
\end{equation}

where $x'$ is a number (coordinate value), $x$ is an operator.

\textbf{Expansion in position basis (1D):}
\begin{equation}
|\alpha\rangle = \int_{-\infty}^{\infty} dx'|x'\rangle\langle x'|\alpha\rangle
\tag{1.186}
\end{equation}

\textbf{Position measurement:}
\begin{itemize}
\item Idealized detector: clicks only when particle is at $x'$
\item After click: state "jumps into" $|x'\rangle$
\item Realistic detector: localizes particle within narrow interval $(x' - \Delta/2, x' + \Delta/2)$
\end{itemize}

\textbf{State change upon measurement:}
\begin{equation}
|\alpha\rangle = \int_{-\infty}^{\infty} dx''|x''\rangle\langle x''|\alpha\rangle \xrightarrow{\text{measurement}} \int_{x'-\Delta/2}^{x'+\Delta/2} dx''|x''\rangle\langle x''|\alpha\rangle
\tag{1.187}
\end{equation}

\textbf{Probability for detection in interval:}
\begin{equation}
|\langle x'|\alpha\rangle|^2 dx'
\tag{1.188}
\end{equation}

where $dx'$ is written for $\Delta$ (narrow interval width).

\textbf{Normalization:}
\begin{equation}
\langle\alpha|\alpha\rangle = 1 \Rightarrow \int_{-\infty}^{\infty} dx'\langle\alpha|x'\rangle\langle x'|\alpha\rangle = 1
\tag{1.190}
\end{equation}

\textbf{Wave function:} $\langle x'|\alpha\rangle$ is the wave function for state $|\alpha\rangle$ in the $x$-representation.

\subsection*{Dirac Delta in Position Measurements}

\textbf{Position eigenstate orthonormality:}
$$\boxed{\langle x''|x'\rangle = \delta(x'' - x')}$$

\textbf{Completeness relation:}
$$\boxed{\int_{-\infty}^{\infty} dx' |x'\rangle\langle x'| = \mathbb{1}}$$

\textbf{Before measurement:}
$$|\alpha\rangle = \int_{-\infty}^{\infty} dx'' |x''\rangle\langle x''|\alpha\rangle$$

\textbf{After measurement in interval $[x' - \Delta/2, x' + \Delta/2]$:}
$$|\alpha\rangle_{\text{after}} = \int_{x'-\Delta/2}^{x'+\Delta/2} dx'' |x''\rangle\langle x''|\alpha\rangle$$

\textbf{Idealized point measurement ($\Delta \to 0$):}
$$\lim_{\Delta \to 0} \frac{1}{\Delta}|\alpha\rangle_{\text{after}} = |x'\rangle\langle x'|\alpha\rangle = |x'\rangle\psi_\alpha(x')$$

\textbf{Projection-type operator} (not a true projection: $\Lambda_{x'}^2 = \delta(0)\Lambda_{x'}$):
$$\Lambda_{x'} = |x'\rangle\langle x'|$$

Matrix elements:
$$\langle x''|\Lambda_{x'}|x'''\rangle = \delta(x'' - x')\delta(x' - x''')$$

\textbf{Key insight:} The Dirac delta $\delta(x'' - x')$ replaces the Kronecker delta $\delta_{a''a'}$ for continuous spectra.

\subsubsection{Three Dimensions}

\textbf{3D position eigenket:}
\begin{equation}
|\alpha\rangle = \int d^3\mathbf{x}'|\mathbf{x}'\rangle\langle\mathbf{x}'|\alpha\rangle
\tag{1.191}
\end{equation}

where $|\mathbf{x}'\rangle \equiv |x', y', z'\rangle$ is a simultaneous eigenket:

\begin{align}
|\mathbf{x}'\rangle &\equiv |x', y', z'\rangle \tag{1.192a}
\end{align}

\begin{align}
x|\mathbf{x}'\rangle = x'|\mathbf{x}'\rangle, \quad y|\mathbf{x}'\rangle = y'|\mathbf{x}'\rangle, \quad z|\mathbf{x}'\rangle = z'|\mathbf{x}'\rangle \tag{1.192b}
\end{align}

\textbf{Position components commute:}
\begin{equation}
[x_i, x_j] = 0
\tag{1.193}
\end{equation}

where $x_1, x_2, x_3$ stand for $x, y, z$ respectively. All three position components can be measured simultaneously to arbitrary accuracy.

\subsection{Translation}

\textbf{Infinitesimal translation operator:}

An operation that changes a well-localized state at $\mathbf{x}'$ to another at $\mathbf{x}' + d\mathbf{x}'$ is an \textbf{infinitesimal translation} by $d\mathbf{x}'$:

\begin{equation}
\mathcal{T}(d\mathbf{x}')|\mathbf{x}'\rangle = |\mathbf{x}' + d\mathbf{x}'\rangle
\tag{1.194}
\end{equation}

Note: $|\mathbf{x}'\rangle$ is NOT an eigenket of the infinitesimal translation operator.

\textbf{Effect on arbitrary state:}
\begin{equation}
|\alpha\rangle \to \mathcal{T}(d\mathbf{x}')|\alpha\rangle = \mathcal{T}(d\mathbf{x}') \int d^3\mathbf{x}'|\mathbf{x}'\rangle\langle\mathbf{x}'|\alpha\rangle = \int d^3\mathbf{x}'|\mathbf{x}' + d\mathbf{x}'\rangle\langle\mathbf{x}'|\alpha\rangle
\tag{1.195}
\end{equation}

Relabeling integration variable:
\begin{equation}
\int d^3\mathbf{x}'|\mathbf{x}' + d\mathbf{x}'\rangle\langle\mathbf{x}'|\alpha\rangle = \int d^3\mathbf{x}'|\mathbf{x}'\rangle\langle\mathbf{x}' - d\mathbf{x}'|\alpha\rangle
\tag{1.196}
\end{equation}

\textbf{Key result:} The wave function of the translated state is obtained by substituting $\mathbf{x}' - d\mathbf{x}'$ for $\mathbf{x}'$ in $\langle\mathbf{x}'|\alpha\rangle$.

\subsubsection{Properties of Translation Operator}

\textbf{Four fundamental properties:}

\textbf{1. Unitarity:} (probability conservation)
\begin{equation}
\mathcal{T}^\dagger(d\mathbf{x}')\mathcal{T}(d\mathbf{x}') = 1
\tag{1.198}
\end{equation}

\textbf{2. Composition:} (successive translations add)
\begin{equation}
\mathcal{T}(d\mathbf{x}'')\mathcal{T}(d\mathbf{x}') = \mathcal{T}(d\mathbf{x}' + d\mathbf{x}'')
\tag{1.199}
\end{equation}

\textbf{3. Inverse:} (opposite direction)
\begin{equation}
\mathcal{T}(-d\mathbf{x}') = \mathcal{T}^{-1}(d\mathbf{x}')
\tag{1.200}
\end{equation}

\textbf{4. Identity limit:}
\begin{equation}
\lim_{d\mathbf{x}' \to 0}\mathcal{T}(d\mathbf{x}') = 1
\tag{1.201}
\end{equation}

\subsubsection{Generator of Translation}

\textbf{Infinitesimal translation operator form:}
\begin{equation}
\boxed{\mathcal{T}(d\mathbf{x}') = 1 - i\mathbf{K} \cdot d\mathbf{x}'}
\tag{1.202}
\end{equation}

where $\mathbf{K}$ components are \textbf{Hermitian operators}.

\textbf{Verify unitarity:}
\begin{align}
\mathcal{T}^\dagger(d\mathbf{x}')\mathcal{T}(d\mathbf{x}') &= (1 + i\mathbf{K}^\dagger \cdot d\mathbf{x}')(1 - i\mathbf{K} \cdot d\mathbf{x}') \notag \\
&= 1 - i(\mathbf{K} - \mathbf{K}^\dagger) \cdot d\mathbf{x}' + O[(d\mathbf{x}')^2] \notag \\
&\simeq 1
\tag{1.203}
\end{align}

\textbf{Verify composition:}
\begin{align}
\mathcal{T}(d\mathbf{x}'')\mathcal{T}(d\mathbf{x}') &= (1 - i\mathbf{K} \cdot d\mathbf{x}'')(1 - i\mathbf{K} \cdot d\mathbf{x}') \notag \\
&\simeq 1 - i\mathbf{K} \cdot (d\mathbf{x}' + d\mathbf{x}'') \notag \\
&= \mathcal{T}(d\mathbf{x}' + d\mathbf{x}'')
\tag{1.204}
\end{align}

Properties 3 and 4 are obviously satisfied.

\subsection{Momentum as a Generator of Translation}

\textbf{Fundamental commutation relation:}

Computing $[\mathbf{x}, \mathcal{T}(d\mathbf{x}')]$:
\begin{align}
\mathbf{x}\mathcal{T}(d\mathbf{x}')|\mathbf{x}'\rangle &= \mathbf{x}|\mathbf{x}' + d\mathbf{x}'\rangle = (\mathbf{x}' + d\mathbf{x}')|\mathbf{x}' + d\mathbf{x}'\rangle \tag{1.205a} \\
\mathcal{T}(d\mathbf{x}')\mathbf{x}|\mathbf{x}'\rangle &= \mathbf{x}'\mathcal{T}(d\mathbf{x}')|\mathbf{x}'\rangle = \mathbf{x}'|\mathbf{x}' + d\mathbf{x}'\rangle \tag{1.205b}
\end{align}

Therefore:
\begin{equation}
[\mathbf{x}, \mathcal{T}(d\mathbf{x}')] |\mathbf{x}'\rangle = d\mathbf{x}'|\mathbf{x}' + d\mathbf{x}'\rangle \simeq d\mathbf{x}'|\mathbf{x}'\rangle
\tag{1.206}
\end{equation}

Since $|\mathbf{x}'\rangle$ is ANY position eigenket, this must be an \textbf{operator identity}:
\begin{equation}
[\mathbf{x}, \mathcal{T}(d\mathbf{x}')] = d\mathbf{x}'
\tag{1.207}
\end{equation}

or
\begin{equation}
-ix\mathbf{K} \cdot d\mathbf{x}' + i\mathbf{K} \cdot d\mathbf{x}'\mathbf{x} = d\mathbf{x}'
\tag{1.208}
\end{equation}

Choosing $d\mathbf{x}'$ in direction of $\hat{\mathbf{x}}_j$ and forming scalar product with $\hat{\mathbf{x}}_i$:
\begin{equation}
\boxed{[x_i, K_j] = i\delta_{ij}}
\tag{1.209}
\end{equation}

where $\delta_{ij}$ is multiplied by the identity operator.

\textbf{Physical significance:} This is the fundamental commutation relation between position operators and the $K$ operators.

\textbf{Connection to momentum:}

From classical mechanics, momentum is the generator of infinitesimal translation:
\begin{equation}
\mathbf{x}_{\text{new}} = \mathbf{x} + d\mathbf{x}, \quad \mathbf{p}_{\text{new}} = \mathbf{p}
\tag{1.210}
\end{equation}

Generating function:
\begin{equation}
F(\mathbf{x}, \mathbf{P}) = \mathbf{x} \cdot \mathbf{P} + \mathbf{p} \cdot d\mathbf{x}
\tag{1.211}
\end{equation}

This has striking similarity to (1.202), suggesting $\mathbf{K}$ is related to momentum.

\textbf{Identification:}

Cannot directly identify $\mathbf{K}$ with momentum (wrong dimensions!). Instead:
\begin{equation}
\boxed{\mathbf{K} = \frac{\mathbf{p}}{\text{universal constant with dimension of action}}}
\tag{1.212}
\end{equation}

This universal constant is $\hbar$ (Planck's constant divided by $2\pi$), from de Broglie relation:
\begin{equation}
\frac{2\pi}{\lambda} = \frac{p}{\hbar}
\tag{1.213}
\end{equation}

Therefore:
\begin{equation}
\boxed{\mathcal{T}(d\mathbf{x}') = 1 - i\mathbf{p} \cdot d\mathbf{x}'/\hbar}
\tag{1.214}
\end{equation}

\textbf{Canonical commutation relation:}
\begin{equation}
\boxed{[x_i, p_j] = i\hbar\delta_{ij}}
\tag{1.215}
\end{equation}

This is one of the most fundamental relations in quantum mechanics!

\subsubsection{Heisenberg Uncertainty Principle}

From (1.215), $x$ and $p_x$ are incompatible observables. Applying the uncertainty relation (1.146):
\begin{equation}
\boxed{\langle(\Delta x)^2\rangle\langle(\Delta p_x)^2\rangle \geq \hbar^2/4}
\tag{1.216}
\end{equation}

This is the \textbf{Heisenberg uncertainty principle}.

\subsubsection{Finite Translation}

By compounding $N$ infinitesimal translations, each of $\Delta x'/N$, and taking $N \to \infty$:
\begin{equation}
\mathcal{T}(\Delta x'\hat{\mathbf{x}}) = \lim_{N\to\infty}\left(1 - \frac{ip_x\Delta x'}{N\hbar}\right)^N = \exp\left(-\frac{ip_x\Delta x'}{\hbar}\right)
\tag{1.218}
\end{equation}

where $\exp(X) \equiv 1 + X + X^2/2! + \cdots$ for any operator $X$.

\subsubsection{Commutation of Translations}

\textbf{Key property:} Successive translations in different directions commute.

From Figure 1.9, going from $A$ to $B$ via $C$ or via $D$ gives the same result:
\begin{align}
\mathcal{T}(\Delta y'\hat{\mathbf{y}})\mathcal{T}(\Delta x'\hat{\mathbf{x}}) &= \mathcal{T}(\Delta x'\hat{\mathbf{x}} + \Delta y'\hat{\mathbf{y}}) \tag{1.220} \\
\mathcal{T}(\Delta x'\hat{\mathbf{x}})\mathcal{T}(\Delta y'\hat{\mathbf{y}}) &= \mathcal{T}(\Delta x'\hat{\mathbf{x}} + \Delta y'\hat{\mathbf{y}})
\end{align}

Expanding to second order in $\Delta x'$ and $\Delta y'$:
\begin{equation}
[\mathcal{T}(\Delta y'\hat{\mathbf{y}}), \mathcal{T}(\Delta x'\hat{\mathbf{x}})] \simeq -\frac{(\Delta x')(\Delta y')}{\hbar^2}[p_y, p_x]
\tag{1.221}
\end{equation}

Since translations commute, requirement (1.220) immediately leads to:
\begin{equation}
\boxed{[p_x, p_y] = 0}
\tag{1.223}
\end{equation}

or more generally:
\begin{equation}
\boxed{[p_i, p_j] = 0}
\tag{1.224}
\end{equation}

\textbf{Physical interpretation:} All momentum components are mutually compatible observables. This is a direct consequence of the fact that translations in different directions commute.

\textbf{Abelian group:} When generators of transformations commute, the group is called \textbf{Abelian}. The translation group in three dimensions is Abelian.

\subsubsection{Momentum Eigenkets}

Since $p_x, p_y, p_z$ all commute, we can conceive of a simultaneous eigenket:
\begin{align}
|\mathbf{p}'\rangle &\equiv |p_x', p_y', p_z'\rangle \tag{1.225a} \\
p_x|\mathbf{p}'\rangle = p_x'|\mathbf{p}'\rangle, \quad p_y|\mathbf{p}'\rangle = p_y'|\mathbf{p}'\rangle, \quad p_z|\mathbf{p}'\rangle = p_z'|\mathbf{p}'\rangle \tag{1.225b}
\end{align}

\textbf{Effect of translation on momentum eigenket:}
\begin{equation}
\mathcal{T}(d\mathbf{x}')|\mathbf{p}'\rangle = \left(1 - \frac{i\mathbf{p} \cdot d\mathbf{x}'}{\hbar}\right)|\mathbf{p}'\rangle = \left(1 - \frac{i\mathbf{p}' \cdot d\mathbf{x}'}{\hbar}\right)|\mathbf{p}'\rangle
\tag{1.226}
\end{equation}

The momentum eigenket remains the same (up to phase). Unlike $|\mathbf{x}'\rangle$, $|\mathbf{p}'\rangle$ IS an eigenket of $\mathcal{T}(d\mathbf{x}')$:
\begin{equation}
[\mathbf{p}, \mathcal{T}(d\mathbf{x}')] = 0
\tag{1.227}
\end{equation}

However, the eigenvalue is complex; $\mathcal{T}(d\mathbf{x}')$, though unitary, is not Hermitian.

\subsection{The Canonical Commutation Relations}

\textbf{Summary of fundamental commutation relations:}
\begin{equation}
\boxed{[x_i, x_j] = 0, \quad [p_i, p_j] = 0, \quad [x_i, p_j] = i\hbar\delta_{ij}}
\tag{1.228}
\end{equation}

These are the \textbf{canonical commutation relations} or \textbf{fundamental commutation relations} - the cornerstone of quantum mechanics (Dirac called them the "fundamental quantum conditions").

\subsubsection{Historical Context}

\textbf{Heisenberg (1925):} Showed that atomic transition frequencies could be understood by associating arrays of numbers (matrices) obeying certain multiplication rules. Born and Jordan recognized these as matrix algebra rules, leading to \textbf{matrix mechanics}.

\textbf{Dirac (1925):} Observed that quantum-mechanical relations can be obtained from classical relations by replacing classical Poisson brackets by commutators:
\begin{equation}
\boxed{[\, , \,]_{\text{classical}} \to \frac{[\, , \,]}{i\hbar}}
\tag{1.229}
\end{equation}

\textbf{Classical Poisson brackets:} For functions of $q$ and $p$:
\begin{equation}
[A(q,p), B(q,p)]_{\text{classical}} \equiv \sum_s \left(\frac{\partial A}{\partial q_s}\frac{\partial B}{\partial p_s} - \frac{\partial A}{\partial p_s}\frac{\partial B}{\partial q_s}\right)
\tag{1.230}
\end{equation}

\textbf{Example:}
\begin{equation}
[x_i, p_j]_{\text{classical}} = \delta_{ij}
\tag{1.231}
\end{equation}

which becomes $[x_i, p_j] = i\hbar\delta_{ij}$ in quantum mechanics (1.215).

\subsubsection{Algebraic Properties}

The following relations hold for both classical Poisson brackets and quantum commutators:

\begin{align}
[A, A] &= 0 \tag{1.232a} \\
[A, B] &= -[B, A] \tag{1.232b} \\
[A, c] &= 0 \quad (c \text{ is just a number}) \tag{1.232c} \\
[A + B, C] &= [A, C] + [B, C] \tag{1.232d} \\
[A, BC] &= [A, B]C + B[A, C] \tag{1.232e} \\
[A, [B, C]] + [B, [C, A]] + [C, [A, B]] &= 0 \tag{1.232f}
\end{align}

The last relation is known as the \textbf{Jacobi identity}.

\subsubsection{Differences: Classical vs Quantum}

\textbf{Important differences:}
\begin{enumerate}
\item \textbf{Dimensions:} Classical Poisson bracket has different dimensions than quantum commutator (due to differentiations in (1.230)). The factor $i\hbar$ in (1.229) corrects this.

\item \textbf{Hermiticity:} 
\begin{itemize}
\item Classical Poisson bracket of real functions is purely real
\item Quantum commutator of two Hermitian operators is anti-Hermitian (Lemma 3, Section 1.4)
\item The factor $i\hbar$ makes $[A,B]/(i\hbar)$ Hermitian when $A$ and $B$ are Hermitian
\end{itemize}
\end{enumerate}

\subsubsection{Our Approach vs Dirac's Approach}

\textbf{Our derivation:} Based on:
\begin{enumerate}
\item Properties of translations
\item Identification of generator of translation with momentum (modulo $\hbar$)
\end{enumerate}

\textbf{Advantages of our approach:}
\begin{itemize}
\item More powerful - can be generalized to observables with no classical analogues
\item Example: Spin angular momentum (Section 1.4) has nothing to do with classical $p$ and $q$
\item Yet spin commutation relations can be derived from properties of rotations (Chapter 3), just as canonical commutation relations were derived from properties of translations
\end{itemize}

\subsubsection{Connection to Classical Mechanics}

\textbf{Dirac's correspondence principle (1.229):} Provides a bridge between classical and quantum mechanics, but is limited to observables with classical analogues.

\textbf{M. Born's gravestone (Göttingen):} 
$$pq - qp = \frac{h}{2\pi i}$$

This encapsulates the birth of matrix mechanics and the fundamental non-commutativity of quantum observables.

\textbf{Key insight:} The canonical commutation relations (1.228) encode:
\begin{itemize}
\item Position and momentum are incompatible observables
\item Position components are mutually compatible
\item Momentum components are mutually compatible
\item The fundamental scale of quantum mechanics is set by $\hbar$
\end{itemize}

% ============================================
% ============================================
\newpage
\section{Wave Functions in Position and Momentum Space}

\subsection{Position-Space Wave Function}

\textbf{Base kets:} Position eigenkets satisfying
\begin{equation}
x|x'\rangle = x'|x'\rangle
\tag{1.233}
\end{equation}

with orthonormality:
\begin{equation}
\langle x''|x'\rangle = \delta(x'' - x')
\tag{1.234}
\end{equation}

\textbf{Expansion in position basis:}
\begin{equation}
|\alpha\rangle = \int dx'|x'\rangle\langle x'|\alpha\rangle
\tag{1.235}
\end{equation}

\textbf{Probability interpretation:}
\begin{equation}
|\langle x'|\alpha\rangle|^2 dx'
\tag{1.236}
\end{equation}
is the probability to find the particle in interval $dx'$ around $x'$.

\textbf{Wave function definition:}
\begin{equation}
\boxed{\langle x'|\alpha\rangle = \psi_\alpha(x')}
\tag{1.237}
\end{equation}

The inner product $\langle x'|\alpha\rangle$ is the \textbf{wave function} for state $|\alpha\rangle$ in position space.

\textbf{Key insight:} In Dirac's formalism, the probabilistic interpretation for expansion coefficient $c_{a'}$ and wave function $\psi_\alpha(x')$ are unified - both are expansion coefficients in different bases.

\subsubsection{Overlap Integral}

\textbf{Inner product as overlap:}
\begin{align}
\langle\beta|\alpha\rangle &= \int dx'\langle\beta|x'\rangle\langle x'|\alpha\rangle \notag \\
&= \int dx'\psi_\beta^*(x')\psi_\alpha(x')
\tag{1.238}
\end{align}

This characterizes the overlap between two wave functions. More generally (independent of representation), $\langle\beta|\alpha\rangle$ is the probability amplitude for state $|\alpha\rangle$ to be found in state $|\beta\rangle$.

\subsubsection{Eigenfunction Expansion}

For discrete eigenvalues:
\begin{equation}
|\alpha\rangle = \sum_{a'}|a'\rangle\langle a'|\alpha\rangle
\tag{1.239}
\end{equation}

Multiply both sides by $\langle x'|$:
\begin{equation}
\langle x'|\alpha\rangle = \sum_{a'}\langle x'|a'\rangle\langle a'|\alpha\rangle
\tag{1.240}
\end{equation}

In wave mechanics notation:
\begin{equation}
\psi_\alpha(x') = \sum_{a'}c_{a'}u_{a'}(x')
\end{equation}

where the \textbf{eigenfunction} is defined as:
\begin{equation}
u_{a'}(x') = \langle x'|a'\rangle
\tag{1.241}
\end{equation}

\subsubsection{Expectation Values}

\textbf{Matrix element in position basis:}
\begin{align}
\langle\beta|A|\alpha\rangle &= \int dx' \int dx''\langle\beta|x'\rangle\langle x'|A|x''\rangle\langle x''|\alpha\rangle \notag \\
&= \int dx' \int dx''\psi_\beta^*(x')\langle x'|A|x''\rangle\psi_\alpha(x'')
\tag{1.242}
\end{align}

\textbf{For position-dependent operators:}

If $A = f(x)$ is a function of position:
\begin{equation}
\langle x'|x^2|x''\rangle = (\langle x'|) \cdot (x''^2|x''\rangle) = x'^2\delta(x' - x'')
\tag{1.244}
\end{equation}

Then the double integral reduces to a single integral:
\begin{align}
\langle\beta|x^2|\alpha\rangle &= \int dx'\langle\beta|x'\rangle x'^2\langle x'|\alpha\rangle \notag \\
&= \int dx'\psi_\beta^*(x')x'^2\psi_\alpha(x')
\tag{1.245}
\end{align}

\textbf{General position function:}
\begin{equation}
\boxed{\langle\beta|f(x)|\alpha\rangle = \int dx'\psi_\beta^*(x')f(x')\psi_\alpha(x')}
\tag{1.246}
\end{equation}

Note: $f(x)$ on left is an operator, $f(x')$ on right is not an operator.

\subsection{Momentum Operator in the Position Basis}

\textbf{Starting from infinitesimal translation:}
\begin{equation}
\left(1 - \frac{ip\Delta x'}{\hbar}\right)|\alpha\rangle = \int dx'\mathcal{T}(\Delta x')|x'\rangle\langle x'|\alpha\rangle
\tag{1.247}
\end{equation}

Using $\mathcal{T}(\Delta x')|x'\rangle = |x' + \Delta x'\rangle$ and relabeling:
\begin{equation}
= \int dx'|x'\rangle\left(\langle x'|\alpha\rangle - \Delta x'\frac{\partial}{\partial x'}\langle x'|\alpha\rangle\right)
\end{equation}

Comparing both sides:
\begin{equation}
\boxed{p|\alpha\rangle = \int dx'|x'\rangle\left(-i\hbar\frac{\partial}{\partial x'}\langle x'|\alpha\rangle\right)}
\tag{1.248}
\end{equation}

or
\begin{equation}
\boxed{\langle x'|p|\alpha\rangle = -i\hbar\frac{\partial}{\partial x'}\langle x'|\alpha\rangle}
\tag{1.249}
\end{equation}

\textbf{Matrix element of momentum:}
\begin{equation}
\boxed{\langle x'|p|x''\rangle = -i\hbar\frac{\partial}{\partial x'}\delta(x' - x'')}
\tag{1.250}
\end{equation}

\textbf{Momentum expectation value:}
\begin{align}
\langle\beta|p|\alpha\rangle &= \int dx'\langle\beta|x'\rangle\left(-i\hbar\frac{\partial}{\partial x'}\langle x'|\alpha\rangle\right) \notag \\
&= \int dx'\psi_\beta^*(x')\left(-i\hbar\frac{\partial}{\partial x'}\right)\psi_\alpha(x')
\tag{1.251}
\end{align}

This is DERIVED, not postulated, from properties of momentum as generator of translation.

\textbf{Higher powers of momentum:}
\begin{align}
\langle x'|p^n|\alpha\rangle &= (-i\hbar)^n\frac{\partial^n}{\partial x'^n}\langle x'|\alpha\rangle \tag{1.252} \\
\langle\beta|p^n|\alpha\rangle &= \int dx'\psi_\beta^*(x')(-i\hbar)^n\frac{\partial^n}{\partial x'^n}\psi_\alpha(x') \tag{1.253}
\end{align}

\subsection{Momentum-Space Wave Function}

\textbf{Symmetry between $x$ and $p$:} Complete symmetry exists (apart from minus signs) due to canonical commutation relations.

\textbf{Momentum eigenkets:}
\begin{equation}
p|p'\rangle = p'|p'\rangle
\tag{1.254}
\end{equation}

with orthonormality:
\begin{equation}
\langle p'|p''\rangle = \delta(p' - p'')
\tag{1.255}
\end{equation}

\textbf{Expansion in momentum basis:}
\begin{equation}
|\alpha\rangle = \int dp'|p'\rangle\langle p'|\alpha\rangle
\tag{1.256}
\end{equation}

\textbf{Momentum-space wave function:}
\begin{equation}
\boxed{\langle p'|\alpha\rangle = \phi_\alpha(p')}
\tag{1.257}
\end{equation}

Probability for momentum in interval $dp'$ is $|\langle p'|\alpha\rangle|^2dp'$.

\textbf{Normalization:}
\begin{equation}
\int dp'\langle\alpha|p'\rangle\langle p'|\alpha\rangle = \int dp'|\phi_\alpha(p')|^2 = 1
\tag{1.258}
\end{equation}

\subsubsection{Transformation Function}

\textbf{Connection between $x$-representation and $p$-representation:}

The transformation function $\langle x'|p'\rangle$ connects the two representations. From (1.249) with $|\alpha\rangle = |p'\rangle$:
\begin{equation}
\langle x'|p|p'\rangle = -i\hbar\frac{\partial}{\partial x'}\langle x'|p'\rangle
\tag{1.259}
\end{equation}

or
\begin{equation}
p'\langle x'|p'\rangle = -i\hbar\frac{\partial}{\partial x'}\langle x'|p'\rangle
\tag{1.260}
\end{equation}

\textbf{Solution:}
\begin{equation}
\boxed{\langle x'|p'\rangle = \mathcal{N}\exp\left(\frac{ip'x'}{\hbar}\right)}
\tag{1.261}
\end{equation}

This is a plane wave! The wave function of a momentum eigenstate is a plane wave (obtained without solving Schrödinger equation).

\textbf{Normalization:}
\begin{equation}
\langle x'|x''\rangle = \int dp'\langle x'|p'\rangle\langle p'|x''\rangle
\tag{1.262}
\end{equation}

gives
\begin{equation}
\delta(x' - x'') = |\mathcal{N}|^2\int dp'\exp\left[\frac{ip'(x' - x'')}{\hbar}\right] = 2\pi\hbar|\mathcal{N}|^2\delta(x' - x'')
\tag{1.263}
\end{equation}

Therefore:
\begin{equation}
\boxed{\langle x'|p'\rangle = \frac{1}{\sqrt{2\pi\hbar}}\exp\left(\frac{ip'x'}{\hbar}\right)}
\tag{1.264}
\end{equation}

\textbf{Fourier transform relations:}
\begin{align}
\langle x'|\alpha\rangle &= \int dp'\langle x'|p'\rangle\langle p'|\alpha\rangle \tag{1.265a} \\
\langle p'|\alpha\rangle &= \int dx'\langle p'|x'\rangle\langle x'|\alpha\rangle \tag{1.265b}
\end{align}

Explicitly:
\begin{align}
\psi_\alpha(x') &= \left[\frac{1}{\sqrt{2\pi\hbar}}\right]\int dp'\exp\left(\frac{ip'x'}{\hbar}\right)\phi_\alpha(p') \tag{1.266a} \\
\phi_\alpha(p') &= \left[\frac{1}{\sqrt{2\pi\hbar}}\right]\int dx'\exp\left(\frac{-ip'x'}{\hbar}\right)\psi_\alpha(x') \tag{1.266b}
\end{align}

These are Fourier transform pairs! The mathematics "knows" Fourier's work on integral transforms.

\subsection{Gaussian Wave Packets}

\textbf{Example:} \textbf{Gaussian wave packet} in position space:
\begin{equation}
\langle x'|\alpha\rangle = \left[\frac{1}{\pi^{1/4}\sqrt{d}}\right]\exp\left[ikx' - \frac{x'^2}{2d^2}\right]
\tag{1.267}
\end{equation}

Plane wave with wave number $k$ modulated by Gaussian profile centered at origin.

\textbf{Properties:}
\begin{itemize}
\item Probability density $|\langle x'|\alpha\rangle|^2$ has Gaussian shape with width $d$
\item Vanishes rapidly for $|x'| > d$
\end{itemize}

\textbf{Expectation values:}
\begin{align}
\langle x \rangle &= 0 \quad \text{(by symmetry)} \tag{1.268} \\
\langle x^2 \rangle &= \frac{d^2}{2} \tag{1.269} \\
\langle(\Delta x)^2\rangle &= \langle x^2\rangle - \langle x\rangle^2 = \frac{d^2}{2} \tag{1.270}
\end{align}

\begin{align}
\langle p \rangle &= \hbar k \tag{1.271a} \\
\langle p^2 \rangle &= \frac{\hbar^2}{2d^2} + \hbar^2k^2 \tag{1.271b} \\
\langle(\Delta p)^2\rangle &= \frac{\hbar^2}{2d^2} \tag{1.272}
\end{align}

\textbf{Uncertainty product:}
\begin{equation}
\boxed{\langle(\Delta x)^2\rangle\langle(\Delta p)^2\rangle = \frac{\hbar^2}{4}}
\tag{1.273}
\end{equation}

This saturates the Heisenberg uncertainty relation! Gaussian wave packet is a \textbf{minimum uncertainty wave packet}.

\textbf{Momentum-space wave function:}

By Fourier transform (completing the square):
\begin{equation}
\langle p'|\alpha\rangle = \sqrt{\frac{d}{\hbar\sqrt{\pi}}}\exp\left[\frac{-(p' - \hbar k)^2d^2}{2\hbar^2}\right]
\tag{1.274}
\end{equation}

Also Gaussian, centered at $\hbar k$ in momentum space!

\textbf{Key insight:} The widths of the two Gaussians are inversely proportional - expressing the constancy of $\langle(\Delta x)^2\rangle\langle(\Delta p)^2\rangle = \hbar^2/4$. Wider spread in $p$-space means narrower spread in $x$-space, and vice versa.

\textbf{Extreme cases:}
\begin{itemize}
\item $d \to \infty$: position-space wave function becomes plane wave (constant everywhere); momentum-space becomes $\delta$-function at $\hbar k$
\item $d \to 0$: position-space becomes $\delta$-function; momentum-space becomes constant
\end{itemize}

Well-localized state (in $x$-space) = superposition of momentum eigenstates with all possible momenta.

\subsection{Generalization to Three Dimensions}

\textbf{Position eigenkets:}
\begin{equation}
\mathbf{x}|\mathbf{x}'\rangle = \mathbf{x}'|\mathbf{x}'\rangle
\tag{1.275}
\end{equation}

\textbf{Momentum eigenkets:}
\begin{equation}
\mathbf{p}|\mathbf{p}'\rangle = \mathbf{p}'|\mathbf{p}'\rangle
\tag{1.276}
\end{equation}

\textbf{Orthonormality:}
\begin{align}
\langle\mathbf{x}'|\mathbf{x}''\rangle &= \delta^3(\mathbf{x}' - \mathbf{x}'') \tag{1.277a} \\
\langle\mathbf{p}'|\mathbf{p}''\rangle &= \delta^3(\mathbf{p}' - \mathbf{p}'') \tag{1.277b}
\end{align}

where $\delta^3(\mathbf{x}' - \mathbf{x}'') = \delta(x' - x'')\delta(y' - y'')\delta(z' - z'')$ (1.278).

\textbf{Completeness:}
\begin{align}
\int d^3\mathbf{x}'|\mathbf{x}'\rangle\langle\mathbf{x}'| &= 1 \tag{1.279a} \\
\int d^3p'|\mathbf{p}'\rangle\langle\mathbf{p}'| &= 1 \tag{1.279b}
\end{align}

\textbf{Expansions:}
\begin{align}
|\alpha\rangle &= \int d^3\mathbf{x}'|\mathbf{x}'\rangle\langle\mathbf{x}'|\alpha\rangle \tag{1.280a} \\
|\alpha\rangle &= \int d^3p'|\mathbf{p}'\rangle\langle\mathbf{p}'|\alpha\rangle \tag{1.280b}
\end{align}

Wave functions: $\psi_\alpha(\mathbf{x}') = \langle\mathbf{x}'|\alpha\rangle$ and $\phi_\alpha(\mathbf{p}') = \langle\mathbf{p}'|\alpha\rangle$.

\textbf{Momentum operator:}
\begin{equation}
\langle\beta|\mathbf{p}|\alpha\rangle = \int d^3\mathbf{x}'\psi_\beta^*(\mathbf{x}')(-i\hbar\nabla')\psi_\alpha(\mathbf{x}')
\tag{1.281}
\end{equation}

\textbf{3D transformation function:}
\begin{equation}
\boxed{\langle\mathbf{x}'|\mathbf{p}'\rangle = \left[\frac{1}{(2\pi\hbar)^{3/2}}\right]\exp\left(\frac{i\mathbf{p}' \cdot \mathbf{x}'}{\hbar}\right)}
\tag{1.282}
\end{equation}

\textbf{3D Fourier transforms:}
\begin{align}
\psi_\alpha(\mathbf{x}') &= \left[\frac{1}{(2\pi\hbar)^{3/2}}\right]\int d^3p'\exp\left(\frac{i\mathbf{p}' \cdot \mathbf{x}'}{\hbar}\right)\phi_\alpha(\mathbf{p}') \tag{1.283a} \\
\phi_\alpha(\mathbf{p}') &= \left[\frac{1}{(2\pi\hbar)^{3/2}}\right]\int d^3\mathbf{x}'\exp\left(\frac{-i\mathbf{p}' \cdot \mathbf{x}'}{\hbar}\right)\psi_\alpha(\mathbf{x}') \tag{1.283b}
\end{align}

\textbf{Dimensional analysis:}
\begin{itemize}
\item 1D: $|\langle x'|\alpha\rangle|^2$ has dimension (length)$^{-1}$, so wave function has dimension (length)$^{-1/2}$
\item 3D: $|\langle\mathbf{x}'|\alpha\rangle|^2$ integrated over volume must be unity, so wave function has dimension (length)$^{-3/2}$
\end{itemize}

% ============================================
% ============================================
\newpage
\section{Time Evolution and the Schrödinger Equation}

\subsection{Time Evolution Operator}

\textbf{Critical distinction:} Time is a \textbf{parameter}, NOT an operator in quantum mechanics. There is no "time operator" analogous to the position operator.

\textbf{Historical note:} de Broglie and Schrödinger were guided by covariant analogy between energy-time and momentum-position. In finished QM, there is NO symmetrical treatment between time and space.

\subsubsection{State Evolution}

\textbf{Notation for time-dependent states:}

State at time $t_0$: $|\alpha\rangle$

State at later time $t$:
\begin{equation}
|\alpha, t_0; t\rangle \quad (t > t_0)
\tag{2.1}
\end{equation}

Reminds us the system \textit{used to be} in state $|\alpha\rangle$ at reference time $t_0$.

\textbf{Continuity condition:}
\begin{equation}
\lim_{t \to t_0}|\alpha, t_0; t\rangle = |\alpha\rangle
\tag{2.2}
\end{equation}

\textbf{Shorthand notation:}
\begin{equation}
|\alpha, t_0; t_0\rangle = |\alpha, t_0\rangle
\tag{2.3}
\end{equation}

\textbf{Time evolution:}
\begin{equation}
|\alpha, t_0\rangle = |\alpha\rangle \xrightarrow{\text{time evolution}} |\alpha, t_0; t\rangle
\tag{2.4}
\end{equation}

\subsubsection{Time-Evolution Operator}

\textbf{Definition:} Operator $\mathscr{U}(t, t_0)$ relates states at different times:
\begin{equation}
\boxed{|\alpha, t_0; t\rangle = \mathscr{U}(t, t_0)|\alpha, t_0\rangle}
\tag{2.5}
\end{equation}

\textbf{Properties of $\mathscr{U}(t, t_0)$:}

\textbf{1. Unitarity:} (probability conservation)

At $t_0$, state expanded in eigenkets of observable $A$:
\begin{equation}
|\alpha, t_0\rangle = \sum_{a'}c_{a'}(t_0)|a'\rangle
\tag{2.6}
\end{equation}

At later time $t$:
\begin{equation}
|\alpha, t_0; t\rangle = \sum_{a'}c_{a'}(t)|a'\rangle
\tag{2.7}
\end{equation}

Individual coefficients change:
\begin{equation}
|c_{a'}(t)| \neq |c_{a'}(t_0)|
\tag{2.8}
\end{equation}

But sum of probabilities conserved:
\begin{equation}
\sum_{a'}|c_{a'}(t_0)|^2 = \sum_{a'}|c_{a'}(t)|^2
\tag{2.9}
\end{equation}

Normalization preserved:
\begin{equation}
\langle\alpha, t_0|\alpha, t_0\rangle = 1 \Rightarrow \langle\alpha, t_0; t|\alpha, t_0; t\rangle = 1
\tag{2.10}
\end{equation}

This requires:
\begin{equation}
\boxed{\mathscr{U}^\dagger(t, t_0)\mathscr{U}(t, t_0) = 1}
\tag{2.11}
\end{equation}

Unitarity is synonymous with probability conservation.

\textbf{2. Composition property:}
\begin{equation}
\mathscr{U}(t_2, t_0) = \mathscr{U}(t_2, t_1)\mathscr{U}(t_1, t_0) \quad (t_2 > t_1 > t_0)
\tag{2.12}
\end{equation}

Time evolution from $t_0$ to $t_2$ = evolution from $t_0$ to $t_1$, then $t_1$ to $t_2$.

Note: Read equation (2.12) from right to left!

\subsubsection{Infinitesimal Time-Evolution Operator}

For infinitesimal time displacement $dt$:
\begin{equation}
|\alpha, t_0; t_0 + dt\rangle = \mathscr{U}(t_0 + dt, t_0)|\alpha, t_0\rangle
\tag{2.13}
\end{equation}

\textbf{3. Continuity:}
\begin{equation}
\lim_{dt \to 0}\mathscr{U}(t_0 + dt, t_0) = 1
\tag{2.14}
\end{equation}

Difference from identity should be first order in $dt$.

\textbf{Infinitesimal form:}
\begin{equation}
\boxed{\mathscr{U}(t_0 + dt, t_0) = 1 - i\Omega dt}
\tag{2.15}
\end{equation}

where $\Omega$ is a \textbf{Hermitian operator}:
\begin{equation}
\Omega^\dagger = \Omega
\tag{2.16}
\end{equation}

\textbf{Verify composition:}
\begin{equation}
\mathscr{U}(t_0 + dt_1 + dt_2, t_0) = \mathscr{U}(t_0 + dt_1 + dt_2, t_0 + dt_1)\mathscr{U}(t_0 + dt_1, t_0)
\tag{2.17}
\end{equation}

Differs from identity by term of order $dt$.

\textbf{Verify unitarity:}
\begin{equation}
\mathscr{U}^\dagger(t_0 + dt, t_0)\mathscr{U}(t_0 + dt, t_0) = (1 + i\Omega^\dagger dt)(1 - i\Omega dt) \simeq 1
\tag{2.18}
\end{equation}

to order $(dt)^2$ or higher.

\subsubsection{Connection to Hamiltonian}

\textbf{Dimension analysis:} $\Omega$ has dimension of frequency (inverse time).

What observable has dimension of frequency?

\textbf{Planck-Einstein relation:}
\begin{equation}
E = \hbar\omega
\tag{2.19}
\end{equation}

Angular frequency $\omega$ is related to energy $E$.

\textbf{From classical mechanics:} Hamiltonian is the generator of time evolution.

\textbf{Identification:}
\begin{equation}
\boxed{\Omega = \frac{H}{\hbar}}
\tag{2.20}
\end{equation}

where $H$ is the \textbf{Hamiltonian operator} (Hermitian).

\textbf{Infinitesimal time-evolution operator:}
\begin{equation}
\boxed{\mathscr{U}(t_0 + dt, t_0) = 1 - \frac{iH dt}{\hbar}}
\tag{2.21}
\end{equation}

\textbf{Important consistency check:} The $\hbar$ here is the SAME $\hbar$ as in the translation operator (1.214). This is necessary to obtain:
\begin{equation}
\frac{d\mathbf{x}}{dt} = \frac{\mathbf{p}}{m}
\tag{2.22}
\end{equation}

as the classical limit of the corresponding quantum-mechanical relation. The two $\hbar$ must be the same!

\textbf{Key physical interpretation:}
\begin{itemize}
\item $H$ is the quantum Hamiltonian (energy operator)
\item $H$ generates time evolution, just as $\mathbf{p}$ generates spatial translation
\item Time evolution is unitary → probability conservation
\item Infinitesimal time evolution characterized by $H/\hbar$
\end{itemize}
\subsection{The Schrödinger Equation}

\textbf{Deriving the differential equation for $\mathscr{U}(t, t_0)$:}

Using the composition property (2.12) with $t_1 \to t$, $t_2 \to t + dt$:
\begin{equation}
\mathscr{U}(t + dt, t_0) = \mathscr{U}(t + dt, t)\mathscr{U}(t, t_0) = \left(1 - \frac{iHdt}{\hbar}\right)\mathscr{U}(t, t_0)
\tag{2.23}
\end{equation}

where $t - t_0$ need not be infinitesimal. We have:
\begin{equation}
\mathscr{U}(t + dt, t_0) - \mathscr{U}(t, t_0) = -i\left(\frac{H}{\hbar}\right)dt\mathscr{U}(t, t_0)
\tag{2.24}
\end{equation}

In differential equation form:
\begin{equation}
\boxed{i\hbar\frac{\partial}{\partial t}\mathscr{U}(t, t_0) = H\mathscr{U}(t, t_0)}
\tag{2.25}
\end{equation}

This is the **Schrödinger equation for the time-evolution operator**. Everything about time development follows from this fundamental equation.

\subsubsection{Schrödinger Equation for State Kets}

Multiplying both sides of (2.25) by $|\alpha, t_0\rangle$ on the right:
\begin{equation}
i\hbar\frac{\partial}{\partial t}\mathscr{U}(t, t_0)|\alpha, t_0\rangle = H\mathscr{U}(t, t_0)|\alpha, t_0\rangle
\tag{2.26}
\end{equation}

Since $|\alpha, t_0\rangle$ doesn't depend on $t$, using (2.5):
\begin{equation}
\boxed{i\hbar\frac{\partial}{\partial t}|\alpha, t_0; t\rangle = H|\alpha, t_0; t\rangle}
\tag{2.27}
\end{equation}

This is the **Schrödinger equation for state kets**.

\textbf{Key insight:} If we know $\mathscr{U}(t, t_0)$ and how it acts on the initial state $|\alpha, t_0\rangle$, we don't need to solve (2.27). Just apply $\mathscr{U}(t, t_0)$ to $|\alpha, t_0\rangle$ to get the state at any time $t$.

\subsubsection{Solutions to the Schrödinger Equation}

Three cases to consider:

\textbf{Case 1: Time-independent Hamiltonian}

$H$ operator remains unchanged as the parameter $t$ is changed. Example: spin-magnetic moment in time-independent magnetic field.

\textbf{Solution:}
\begin{equation}
\boxed{\mathscr{U}(t, t_0) = \exp\left[\frac{-iH(t - t_0)}{\hbar}\right]}
\tag{2.28}
\end{equation}

\textbf{Proof:} Expand the exponential:
\begin{equation}
\exp\left[\frac{-iH(t - t_0)}{\hbar}\right] = 1 + \frac{-iH(t - t_0)}{\hbar} + \left[\frac{(-i)^2}{2}\right]\left[\frac{H(t - t_0)}{\hbar}\right]^2 + \cdots
\tag{2.29}
\end{equation}

Time derivative:
\begin{equation}
\frac{\partial}{\partial t}\exp\left[\frac{-iH(t - t_0)}{\hbar}\right] = \frac{-iH}{\hbar} + \left[\frac{(-i)^2}{2}\right]2\left(\frac{H}{\hbar}\right)^2(t - t_0) + \cdots
\tag{2.30}
\end{equation}

This satisfies (2.25). Boundary condition also satisfied: as $t \to t_0$, (2.28) reduces to identity.

\textbf{Alternative derivation:} Compound infinitesimal time-evolution operators (analogous to finite translation):
\begin{equation}
\lim_{N \to \infty}\left[1 - \frac{(iH/\hbar)(t - t_0)}{N}\right]^N = \exp\left[\frac{-iH(t - t_0)}{\hbar}\right]
\tag{2.31}
\end{equation}

\textbf{Case 2: Time-dependent $H$ with commuting $H$ at different times}

Example: Spin-magnetic moment in magnetic field with time-varying strength but constant direction.

\textbf{Solution:}
\begin{equation}
\mathscr{U}(t, t_0) = \exp\left[-\left(\frac{i}{\hbar}\right)\int_{t_0}^t dt'H(t')\right]
\tag{2.32}
\end{equation}

Proved similarly by replacing $H(t - t_0)$ in (2.29) and (2.30) by $\int_{t_0}^t dt'H(t')$.

\textbf{Case 3: Non-commuting $H$ at different times}

Example: Spin-magnetic moment where magnetic field direction also changes with time.

Since $S_x$ and $S_y$ don't commute, $H(t_1)$ and $H(t_2)$ don't commute either.

\textbf{Solution (Dyson series):}
\begin{equation}
\mathscr{U}(t, t_0) = 1 + \sum_{n=1}^{\infty}\left(\frac{-i}{\hbar}\right)^n \int_{t_0}^t dt_1 \int_{t_0}^{t_1} dt_2 \cdots \int_{t_0}^{t_{n-1}} dt_n H(t_1)H(t_2)\cdots H(t_n)
\tag{2.33}
\end{equation}

This is the **Dyson series** (F. J. Dyson). Proof similar to Chapter 5 (interaction picture).

\textbf{Note:} In elementary applications, only Case 1 is of practical interest. We assume time-independent $H$ for the rest of this chapter. Time-dependent Hamiltonians appear in Chapter 5.

\subsection{Energy Eigenkets}

\textbf{Motivation:} To evaluate $\mathscr{U}(t, t_0)$ on a general initial ket $|\alpha\rangle$, we must know how it acts on the base kets.

\textbf{Key simplification:} Use base kets that are eigenkets of an observable $A$ such that:
\begin{equation}
\boxed{[A, H] = 0}
\tag{2.34}
\end{equation}

Then eigenkets of $A$ are also eigenkets of $H$, called **energy eigenkets**, with eigenvalues $E_{a'}$:
\begin{equation}
H|a'\rangle = E_{a'}|a'\rangle
\tag{2.35}
\end{equation}

\subsubsection{Time Evolution in Energy Basis}

Expand time-evolution operator (taking $t_0 = 0$ for simplicity):
\begin{align}
\exp\left(\frac{-iHt}{\hbar}\right) &= \sum_{a'}\sum_{a''}|a''\rangle\langle a''|\exp\left(\frac{-iHt}{\hbar}\right)|a'\rangle\langle a'| \notag \\
&= \sum_{a'}|a'\rangle\exp\left(\frac{-iE_{a'}t}{\hbar}\right)\langle a'|
\tag{2.36}
\end{align}

This enables solving any initial-value problem once we know the expansion of initial ket in terms of $\{|a'\rangle\}$.

\textbf{Example:} Initial ket expansion:
\begin{equation}
|\alpha, t_0 = 0\rangle = \sum_{a'}|a'\rangle\langle a'|\alpha\rangle = \sum_{a'}c_{a'}|a'\rangle
\tag{2.37}
\end{equation}

Then:
\begin{equation}
|\alpha, t_0 = 0; t\rangle = \exp\left(\frac{-iHt}{\hbar}\right)|\alpha, t_0 = 0\rangle = \sum_{a'}|a'\rangle\langle a'|\alpha\rangle\exp\left(\frac{-iE_{a'}t}{\hbar}\right)
\tag{2.38}
\end{equation}

\textbf{Expansion coefficient evolution:}
\begin{equation}
\boxed{c_{a'}(t = 0) \to c_{a'}(t) = c_{a'}(t = 0)\exp\left(\frac{-iE_{a'}t}{\hbar}\right)}
\tag{2.39}
\end{equation}

**Modulus unchanged!** Only relative phases among components vary because oscillation frequencies are different.

\subsubsection{Energy Eigenstates as Stationary States}

Special case: Initial state is one of $\{|a'\rangle\}$:
\begin{equation}
|\alpha, t_0 = 0\rangle = |a'\rangle
\tag{2.40}
\end{equation}

At later time:
\begin{equation}
|\alpha, t_0 = 0; t\rangle = |a'\rangle\exp\left(\frac{-iE_{a'}t}{\hbar}\right)
\tag{2.41}
\end{equation}

If system is initially a simultaneous eigenstate of $A$ and $H$, it remains so at all times. Only phase modulation $\exp(-iE_{a'}t/\hbar)$ occurs.

\textbf{Constants of motion:} Observable compatible with $H$ [see (2.34)] is a \textit{constant of the motion}. More on this in Heisenberg equation of motion.

\subsubsection{General Strategy for Solving Time Evolution}

\textbf{Basic task in quantum dynamics:}
\begin{enumerate}
\item Find observable that commutes with $H$
\item Evaluate its eigenvalues
\item Expand initial ket in terms of eigenkets of that observable
\item Apply time-evolution operator - this merely changes phase of each expansion coefficient [see (2.39)]
\end{enumerate}

\textbf{Extension to multiple compatible observables:}

If several mutually compatible observables all commute with $H$:
\begin{equation}
[A,B] = [B,C] = [A,C] = \cdots = 0, \quad [A,H] = [B,H] = [C,H] = \cdots = 0
\tag{2.42}
\end{equation}

Using collective index notation [Section 1.4, (1.130)]:
\begin{equation}
\exp\left(\frac{-iHt}{\hbar}\right) = \sum_{\mathcal{K}'}|\mathcal{K}'\rangle\exp\left(\frac{-iE_{\mathcal{K}'}t}{\hbar}\right)\langle\mathcal{K}'|
\tag{2.43}
\end{equation}

where $E_{\mathcal{K}'}$ is uniquely specified once $a', b', c', \ldots$ are specified.

\textbf{Critical importance:} Find a \textit{complete set of mutually compatible observables that also commute with $H$}. Express initial ket as superposition of simultaneous eigenkets of $A, B, C, \ldots$ and $H$. Apply time-evolution operator (2.43).

This is the general procedure for solving initial-value problems with time-independent $H$.

\subsection{Time Dependence of Expectation Values}

\subsubsection{Stationary States}

Suppose at $t = 0$ the initial state is an eigenstate of observable $A$ that commutes with $H$, as in (2.40).

Consider expectation value of another observable $B$ (which need not commute with $A$ or $H$).

At later time:
\begin{equation}
|a', t_0 = 0; t\rangle = \mathscr{U}(t, 0)|a'\rangle
\tag{2.44}
\end{equation}

The expectation value $\langle B \rangle$ is:
\begin{align}
\langle B \rangle &= (\langle a'|\mathscr{U}^\dagger(t, 0)) \cdot B \cdot (\mathscr{U}(t, 0)|a'\rangle) \notag \\
&= \langle a'|\exp\left(\frac{iE_{a'}t}{\hbar}\right)B\exp\left(\frac{-iE_{a'}t}{\hbar}\right)|a'\rangle \notag \\
&= \langle a'|B|a'\rangle
\tag{2.45}
\end{align}

which is **independent of $t$**.

\textbf{Key result:} The expectation value of an observable taken with respect to an energy eigenstate does NOT change with time.

\textbf{Definition:} An energy eigenstate is called a **stationary state**.

\subsubsection{Nonstationary States}

More interesting: expectation value taken with respect to a **superposition** of energy eigenstates (a **nonstationary state**).

Suppose initially:
\begin{equation}
|\alpha, t_0 = 0\rangle = \sum_{a'}c_{a'}|a'\rangle
\tag{2.46}
\end{equation}

Expectation value of $B$:
\begin{align}
\langle B \rangle &= \left[\sum_{a'}c_{a'}^*\langle a'|\exp\left(\frac{iE_{a'}t}{\hbar}\right)\right] \cdot B \cdot \left[\sum_{a''}c_{a''}\exp\left(\frac{-iE_{a''}t}{\hbar}\right)|a''\rangle\right] \notag \\
&= \sum_{a'}\sum_{a''}c_{a'}^*c_{a''}\langle a'|B|a''\rangle\exp\left[\frac{-i(E_{a''} - E_{a'})t}{\hbar}\right]
\tag{2.47}
\end{align}

\textbf{Key result:} The expectation value consists of **oscillating terms** whose angular frequencies are determined by **Bohr's frequency condition**:
\begin{equation}
\boxed{\omega_{a''a'} = \frac{(E_{a''} - E_{a'})}{\hbar}}
\tag{2.48}
\end{equation}

\textbf{Physical interpretation:}
\begin{itemize}
\item \textbf{Stationary states}: Expectation values are time-independent (constant)
\item \textbf{Nonstationary states}: Expectation values oscillate at frequencies corresponding to energy differences between eigenstates
\item The oscillation frequencies are exactly the Bohr frequencies from old quantum theory
\item Off-diagonal matrix elements $\langle a'|B|a''\rangle$ (with $a' \neq a''$) contribute oscillating terms
\item Diagonal elements $\langle a'|B|a'\rangle$ contribute time-independent terms
\end{itemize}

\textbf{Connection to old quantum theory:} Bohr's frequency condition, which was postulated in the old quantum theory, emerges naturally from the time evolution of quantum states. The frequency of radiation emitted during a transition between states with energies $E_{a''}$ and $E_{a'}$ is:
$$\omega = \frac{E_{a''} - E_{a'}}{\hbar}$$

This is exactly what we get from the time dependence of expectation values in the superposition state!

\textbf{Summary:}
\begin{itemize}
\item Energy eigenstates → stationary states → time-independent expectation values
\item Superpositions of energy eigenstates → nonstationary states → oscillating expectation values at Bohr frequencies
\item The quantum mechanical framework naturally reproduces Bohr's frequency condition without additional postulates
\end{itemize}

\subsection{Spin Precession}

\textbf{Example:} Spin-$\frac{1}{2}$ system with magnetic moment $e\hbar/(2m_ec)$ in external magnetic field $\mathbf{B}$.

\textbf{Hamiltonian:}
\begin{equation}
H = -\left(\frac{e}{m_ec}\right)\mathbf{S} \cdot \mathbf{B}
\tag{2.49}
\end{equation}

where $e < 0$ for the electron.

Taking $\mathbf{B}$ as static, uniform field in $z$-direction:
\begin{equation}
H = -\left(\frac{eB}{m_ec}\right)S_z
\tag{2.50}
\end{equation}

Since $S_z$ and $H$ differ by a multiplicative constant, they commute. $S_z$ eigenstates are also energy eigenstates:
\begin{equation}
E_{\pm} = \mp \frac{e\hbar B}{2m_ec}, \quad \text{for } S_z\pm
\tag{2.51}
\end{equation}

\textbf{Define angular frequency:}
\begin{equation}
\omega \equiv \frac{|e|B}{m_ec}
\tag{2.52}
\end{equation}

Then:
\begin{equation}
H = \omega S_z
\tag{2.53}
\end{equation}

\textbf{Time-evolution operator:}
\begin{equation}
\mathscr{U}(t, 0) = \exp\left(\frac{-i\omega S_z t}{\hbar}\right)
\tag{2.54}
\end{equation}

\subsubsection{General Superposition State}

Initial state (general superposition):
\begin{equation}
|\alpha\rangle = c_+|+\rangle + c_-|-\rangle
\tag{2.55}
\end{equation}

State at later time:
\begin{equation}
|\alpha, t_0 = 0; t\rangle = c_+\exp\left(\frac{-i\omega t}{2}\right)|+\rangle + c_-\exp\left(\frac{+i\omega t}{2}\right)|-\rangle
\tag{2.56}
\end{equation}

where:
\begin{equation}
H|\pm\rangle = \left(\frac{\pm\hbar\omega}{2}\right)|\pm\rangle
\tag{2.57}
\end{equation}

\textbf{Case 1: Initially spin-up}

If $c_+ = 1$, $c_- = 0$:
\begin{equation}
|\alpha, t_0 = 0; t\rangle = \exp\left(\frac{-i\omega t}{2}\right)|+\rangle
\tag{2.58}
\end{equation}

Remains in spin-up state (stationary state).

\textbf{Case 2: Initially $S_x+$ state}

Comparing (1.110a) with (2.55):
\begin{equation}
c_+ = c_- = \frac{1}{\sqrt{2}}
\tag{2.59}
\end{equation}

\textbf{Probabilities for $S_x\pm$ at later time:}
\begin{align}
|\langle S_x; \pm|\alpha, t_0 = 0; t\rangle|^2 &= \left|\left[\left(\frac{1}{\sqrt{2}}\right)\langle+| \pm \left(\frac{1}{\sqrt{2}}\right)\langle-|\right] \cdot \left[\left(\frac{1}{\sqrt{2}}\right)\exp\left(\frac{-i\omega t}{2}\right)|+\rangle\right.\right. \notag \\
&\quad\left.\left. + \left(\frac{1}{\sqrt{2}}\right)\exp\left(\frac{+i\omega t}{2}\right)|-\rangle\right]\right|^2 \notag \\
&= \left|\frac{1}{2}\exp\left(\frac{-i\omega t}{2}\right) \pm \frac{1}{2}\exp\left(\frac{+i\omega t}{2}\right)\right|^2 \notag \\
&= \cos^2\frac{\omega t}{2} \quad \text{for } S_x+ \tag{2.60a} \\
&= \sin^2\frac{\omega t}{2} \quad \text{for } S_x- \tag{2.60b}
\end{align}

Even though spin is initially in positive $x$-direction, the magnetic field in $z$-direction causes it to rotate; finite probability for finding $S_x-$ at later time.

Sum of probabilities equals unity at all times (unitarity).

\subsubsection{Expectation Values}

Using (1.99), expectation value of $S_x$:
\begin{align}
\langle S_x \rangle &= \left(\frac{\hbar}{2}\right)\cos^2\left(\frac{\omega t}{2}\right) + \left(\frac{-\hbar}{2}\right)\sin^2\left(\frac{\omega t}{2}\right) \notag \\
&= \left(\frac{\hbar}{2}\right)\cos\omega t
\tag{2.61}
\end{align}

Oscillates with angular frequency $\omega = (E_+ - E_-)/\hbar$ (Bohr frequency condition).

Similarly:
\begin{equation}
\langle S_y \rangle = \left(\frac{\hbar}{2}\right)\sin\omega t
\tag{2.62a}
\end{equation}

\begin{equation}
\langle S_z \rangle = 0
\tag{2.62b}
\end{equation}

\textbf{Physical interpretation:} The spin **precesses** in the $xy$-plane. More on spin precession in Chapter 3 (rotation operators).

\textbf{Experimental verification:} Spin precession is well established. Used to measure gyromagnetic ratio $g$ for pointlike particles (electrons, muons) obeying Dirac equation, where $g = 2$. Quantum field theory predicts small deviations from $g = 2$, which have been precisely measured.

Recent experiment: Bennett et al., *Phys. Rev. D*, **73** (2006) 072003. Muons injected into storage ring; spins precess in lock step with momentum vector only if $g \equiv 2$. Observation of precession measures $g - 2$ directly with precision better than one part per million.

\subsection{Neutrino Oscillations}

\textbf{Introduction:} Quantum-mechanical dynamics leading to interference in a two-state system - based on current physics research.

\textbf{Neutrinos:} Elementary particles with:
- No charge
- Very small mass (much smaller than electron)
- Three distinct "flavors"

For this discussion, consider only two flavors identified by their interactions:
- $|\nu_e\rangle$: interacts with electrons
- $|\nu_\mu\rangle$: interacts with muons

These are eigenstates of interaction Hamiltonian.

\subsubsection{Mass Eigenstates vs Flavor Eigenstates}

Neutrinos may have other interactions with energy eigenvalues corresponding to well-defined masses. "Mass eigenstates" with eigenvalues $E_1$ and $E_2$, masses $m_1$ and $m_2$, denoted $|\nu_1\rangle$ and $|\nu_2\rangle$.

\textbf{Mixing:} Flavor eigenstates related to mass eigenstates by unitary transformation with mixing angle $\theta$:
\begin{align}
|\nu_e\rangle &= \cos\theta|\nu_1\rangle - \sin\theta|\nu_2\rangle \tag{2.63a} \\
|\nu_\mu\rangle &= \sin\theta|\nu_1\rangle + \cos\theta|\nu_2\rangle \tag{2.63b}
\end{align}

If mixing angle were zero, $|\nu_e\rangle$ and $|\nu_\mu\rangle$ would be same as $|\nu_1\rangle$ and $|\nu_2\rangle$. However, no theoretical reason why this should be the case. $\theta$ is a free parameter determined only by experiment.

\subsubsection{Neutrino Oscillation Phenomenon}

Suppose at $t = 0$, we prepare momentum eigenstate of one flavor, say $|\nu_e\rangle$. According to (2.63a), the two different mass eigenstate components evolve with different frequencies, developing relative phase difference.

If mass difference is small enough, phase difference builds up over macroscopic distance. By measuring interference as function of distance, we observe **oscillations** with:
- Period depending on mass difference
- Amplitude depending on mixing angle

\textbf{Relativistic energy:} Neutrinos have very low mass, highly relativistic. For fixed momentum $p$:
\begin{equation}
E = [p^2c^2 + m^2c^4]^{1/2} \approx pc\left(1 + \frac{m^2c^2}{2p^2}\right)
\tag{2.64}
\end{equation}

Allowing state $|\nu_e\rangle$ to evolve, probability it still appears as $|\nu_e\rangle$ (not $|\nu_\mu\rangle$):
\begin{equation}
\boxed{P(\nu_e \to \nu_e) = 1 - \sin^2 2\theta \sin^2\left(\Delta m^2 c^4 \frac{L}{4E\hbar c}\right)}
\tag{2.65}
\end{equation}

where:
- $\Delta m^2 \equiv m_1^2 - m_2^2$
- $L = ct$ is flight distance
- $E = pc$ is nominal neutrino energy

\textbf{Experimental verification:} Oscillations dramatically observed by KamLAND experiment (Figure 2.2). Neutrinos from nuclear reactors detected at distance $\sim 150$ km. Rate compared to expected from reactor power and properties. Curve not perfect sine wave because reactors not all at same distance from detector.

\textbf{Key result:} Neutrino oscillations provide direct evidence for:
- Non-zero neutrino masses
- Neutrino mixing (non-zero mixing angle $\theta$)
- Quantum mechanical interference over macroscopic distances

\subsection{Correlation Amplitude and the Energy-Time Uncertainty Relation}

\textbf{Question:} How are state kets at different times correlated?

Initial state at $t = 0$: $|\alpha\rangle$

State at later time: $|\alpha, t_0 = 0; t\rangle = \mathscr{U}(t, 0)|\alpha\rangle$

\textbf{Measure of similarity:} Construct inner product between state kets at different times.

\subsubsection{Correlation Amplitude}

\textbf{Definition:}
\begin{equation}
C(t) \equiv \langle\alpha|\alpha, t_0 = 0; t\rangle = \langle\alpha|\mathscr{U}(t, 0)|\alpha\rangle
\tag{2.66}
\end{equation}

The **correlation amplitude** $C(t)$ provides quantitative measure of "resemblance" between state kets at different times.

\textbf{Case 1: Energy eigenstate (stationary state)}

If initial ket $|\alpha\rangle = |a'\rangle$ is eigenket of $H$:
\begin{equation}
C(t) = \langle a'|a', t_0 = 0; t\rangle = \exp\left(\frac{-iE_{a'}t}{\hbar}\right)
\tag{2.67}
\end{equation}

Modulus of correlation amplitude is **unity at all times** (not surprising for stationary state).

\textbf{Case 2: Superposition of energy eigenstates}

For general superposition $|\alpha\rangle = \sum_{a'}c_{a'}|a'\rangle$ (as in 2.37):
\begin{align}
C(t) &= \left(\sum_{a'}c_{a'}^*\langle a'|\right)\left[\sum_{a''}c_{a''}\exp\left(\frac{-iE_{a''}t}{\hbar}\right)|a''\rangle\right] \notag \\
&= \sum_{a'}|c_{a'}|^2\exp\left(\frac{-iE_{a'}t}{\hbar}\right)
\tag{2.68}
\end{align}

\textbf{Key observation:} Summing over many terms with oscillating time dependence of different frequencies → strong cancellation possible for moderately large $t$. Correlation amplitude starts with unity at $t = 0$ and decreases in magnitude with time.

\subsubsection{Quasi-Continuous Energy Spectrum}

For state ket as superposition of many energy eigenkets with similar energies (essentially quasi-continuous spectrum), replace sum by integral:
\begin{equation}
\sum_{a'} \to \int dE\rho(E), \quad c_{a'} \to g(E)\Big|_{E \simeq E_{a'}}
\tag{2.69}
\end{equation}

where $\rho(E)$ is density of energy eigenstates.

Expression (2.68) becomes:
\begin{equation}
C(t) = \int dE|g(E)|^2\rho(E)\exp\left(\frac{-iEt}{\hbar}\right)
\tag{2.70}
\end{equation}

subject to normalization:
\begin{equation}
\int dE|g(E)|^2\rho(E) = 1
\tag{2.71}
\end{equation}

\subsubsection{Energy Distribution Peaked at $E_0$}

In realistic physical situation, $|g(E)|^2\rho(E)$ peaked around $E = E_0$ with width $\Delta E$.

Writing (2.70) as:
\begin{equation}
C(t) = \exp\left(\frac{-iE_0 t}{\hbar}\right)\int dE|g(E)|^2\rho(E)\exp\left[\frac{-i(E - E_0)t}{\hbar}\right]
\tag{2.72}
\end{equation}

\textbf{Analysis:} As $t$ becomes large, integrand oscillates very rapidly unless energy interval $|E - E_0|$ is small compared with $\hbar/t$. 

If interval for which $|E - E_0| \simeq \hbar/t$ is much narrower than $\Delta E$ (width of $|g(E)|^2\rho(E)$), we get essentially no contribution to $C(t)$ because of strong cancellations.

\textbf{Characteristic time scale:}

Modulus of correlation amplitude starts becoming appreciably different from 1 at:
\begin{equation}
\boxed{t \simeq \frac{\hbar}{\Delta E}}
\tag{2.73}
\end{equation}

\textbf{Applicability:} Though derived for quasi-continuous spectrum, this also makes sense for two-level system. In spin-precession problem, state ket initially $|S_x+\rangle$ starts losing identity after $\sim 1/\omega = \hbar/(E_+ - E_-)$, evident from (2.60).

\subsubsection{Energy-Time Uncertainty Relation}

\textbf{Summary:} As result of time evolution, state ket of physical system ceases to retain its original form after time interval of order $\hbar/\Delta E$.

This is often said to illustrate the **energy-time uncertainty relation**:
\begin{equation}
\boxed{\Delta t \Delta E \simeq \hbar}
\tag{2.74}
\end{equation}

\textbf{CRITICAL DISTINCTION:} This energy-time uncertainty relation is of **very different nature** from the uncertainty relation between two incompatible observables discussed in Section 1.4.

\textbf{Key differences:}
\begin{itemize}
\item Time is a parameter, NOT an operator
\item There is no "time operator" in quantum mechanics
\item $\Delta t$ here represents characteristic time scale for state to change appreciably
\item NOT analogous to position-momentum uncertainty relation
\item Different physical and mathematical origin
\end{itemize}

\textbf{Physical interpretation of $\Delta t$:}
- Time scale over which correlation amplitude decreases significantly
- Time for superposition state to "lose its identity"
- Related to lifetime of unstable states
- Connection to time-dependent perturbation theory (Chapter 5)

\textbf{Physical interpretation of $\Delta E$:}
- Energy spread in the superposition state
- Width of energy distribution $|g(E)|^2\rho(E)$
- NOT uncertainty in energy measurement at fixed time

We return to (2.74) in Chapter 5 in connection with time-dependent perturbation theory, where it relates to:
- Lifetime of unstable quantum states
- Energy width of resonances
- Transition rates in time-dependent processes

\textbf{Warning:} Do not confuse this with the Heisenberg uncertainty principle for position-momentum! They have fundamentally different origins and meanings.

\newpage
% ============================================
\section{The Schrödinger Versus the Heisenberg Picture}

\subsection{Unitary Operators}

\textbf{Two formulations of quantum dynamics:}
\begin{itemize}
\item \textbf{Schrödinger picture:} Time-evolution operator affects state kets (state kets vary with time)
\item \textbf{Heisenberg picture:} Observables vary with time (state kets remain fixed)
\end{itemize}

Before discussing differences, we review unitary operators in general.

\subsubsection{Uses of Unitary Operators}

\textbf{Section 1.5:} Unitary operators relate base kets in different representations. State kets themselves don't change; only expansion coefficients change.

\textbf{Sections 1.6 and 2.1:} Unitary operators that actually **change state kets**:
- Translation operator $\mathcal{T}(d\mathbf{x})$ (Section 1.6)
- Time-evolution operator $\mathscr{U}(t, t_0)$ (Section 2.1)

\textbf{General transformation:}
\begin{equation}
|\alpha\rangle \to U|\alpha\rangle
\tag{2.75}
\end{equation}

where $U$ may stand for $\mathcal{T}(d\mathbf{x})$ or $\mathscr{U}(t, t_0)$. Here $U|\alpha\rangle$ is the state ket for a physical system that has actually undergone translation or time evolution.

\subsubsection{Invariance Under Unitary Transformations}

\textbf{Key property:} Under unitary transformation that changes state kets, inner product of state bra and state ket remains unchanged:
\begin{equation}
\langle\beta|\alpha\rangle \to \langle\beta|U^\dagger U|\alpha\rangle = \langle\beta|\alpha\rangle
\tag{2.76}
\end{equation}

\textbf{Expectation values:} Transformations affect state kets but not operators. How must $\langle\beta|X|\alpha\rangle$ change?
\begin{equation}
\langle\beta|X|\alpha\rangle \to (\langle\beta|U^\dagger) \cdot X \cdot (U|\alpha\rangle) = \langle\beta|U^\dagger XU|\alpha\rangle
\tag{2.77}
\end{equation}

\textbf{Mathematical observation from associative property:}
\begin{equation}
(\langle\beta|U^\dagger) \cdot X \cdot (U|\alpha\rangle) = \langle\beta| \cdot (U^\dagger XU) \cdot |\alpha\rangle
\tag{2.78}
\end{equation}

Is there any physics in this? This mathematical identity suggests **two equivalent approaches** to unitary transformations.

\subsubsection{Two Approaches to Unitary Transformations}

\textbf{Approach 1: Transform state kets}
\begin{equation}
|\alpha\rangle \to U|\alpha\rangle, \quad \text{with operators unchanged}
\tag{2.79a}
\end{equation}

\textbf{Approach 2: Transform operators}
\begin{equation}
X \to U^\dagger XU, \quad \text{with state kets unchanged}
\tag{2.79b}
\end{equation}

\textbf{Connection to classical physics:} In classical physics we don't introduce state kets, yet we talk about translation, time evolution, etc. This is possible because these operations actually change quantities such as $\mathbf{x}$ and $\mathbf{L}$, which are observables of classical mechanics.

\textbf{Conjecture:} A closer connection with classical physics may be established if we follow **approach 2**.

\subsubsection{Example: Infinitesimal Translation}

Return to infinitesimal translation operator $\mathcal{T}(d\mathbf{x}')$. Section 1.6 formalism based on **approach 1**: $\mathcal{T}(d\mathbf{x}')$ affects state kets, not position operator:

\textbf{Approach 1:}
\begin{equation}
|\alpha\rangle \to \left(1 - \frac{i\mathbf{p} \cdot d\mathbf{x}'}{\hbar}\right)|\alpha\rangle, \quad \mathbf{x} \to \mathbf{x}
\tag{2.80}
\end{equation}

\textbf{Approach 2:}
If we follow approach 2 instead:
\begin{align}
|\alpha\rangle &\to |\alpha\rangle, \notag \\
\mathbf{x} &\to \left(1 + \frac{i\mathbf{p} \cdot d\mathbf{x}'}{\hbar}\right)\mathbf{x}\left(1 - \frac{i\mathbf{p} \cdot d\mathbf{x}'}{\hbar}\right) \notag \\
&= \mathbf{x} + \left(\frac{i}{\hbar}\right)[\mathbf{p} \cdot d\mathbf{x}', \mathbf{x}] \notag \\
&= \mathbf{x} + d\mathbf{x}'
\tag{2.81}
\end{align}

\textbf{Result:} Both approaches lead to same result for expectation value:
\begin{equation}
\langle\mathbf{x}\rangle \to \langle\mathbf{x}\rangle + \langle d\mathbf{x}'\rangle
\tag{2.82}
\end{equation}

(Exercise for reader to verify.)

\textbf{Physical interpretation:} In approach 2, the position operator itself shifts by $d\mathbf{x}'$ - directly analogous to classical mechanics where position changes under translation!

\subsection{State Kets and Observables in the Schrödinger and Heisenberg Pictures}

Return to time-evolution operator $\mathscr{U}(t, t_0)$. In previous section (2.1) we examined how state kets evolve with time - this was **approach 1**, known as the **Schrödinger picture** when applied to time evolution.

Alternatively we may follow **approach 2**, known as the **Heisenberg picture** when applied to time evolution.

\subsubsection{Schrödinger Picture}

\textbf{Operators:} Fixed in time
- Operators corresponding to observables like $x$, $p_y$, $S_z$ are fixed in time

\textbf{State kets:} Vary with time
- As indicated in previous section

\subsubsection{Heisenberg Picture}

\textbf{Operators:} Vary with time
- Operators corresponding to observables vary with time

\textbf{State kets:} Fixed (frozen at $t_0$)
- Frozen at what they were at $t_0$

\textbf{Simplified notation:} Set $t_0 = 0$ for simplicity, work with:
\begin{equation}
\mathscr{U}(t, t_0 = 0) \equiv \mathscr{U}(t) = \exp\left(\frac{-iHt}{\hbar}\right)
\tag{2.83}
\end{equation}

\subsubsection{Heisenberg Picture Observable}

Motivated by (2.79b) of approach 2, define Heisenberg picture observable:
\begin{equation}
\boxed{A^{(H)}(t) \equiv \mathscr{U}^\dagger(t)A^{(S)}\mathscr{U}(t)}
\tag{2.84}
\end{equation}

where superscripts $H$ and $S$ stand for Heisenberg and Schrödinger, respectively.

\textbf{At $t = 0$:} Heisenberg and Schrödinger observables coincide:
\begin{equation}
A^{(H)}(0) = A^{(S)}
\tag{2.85}
\end{equation}

\subsubsection{State Kets in Both Pictures}

State kets also coincide at $t = 0$; at later $t$ the Heisenberg picture state ket is frozen:

\textbf{Heisenberg picture:}
\begin{equation}
|\alpha, t_0 = 0; t\rangle_H = |\alpha, t_0 = 0\rangle, \quad \text{independent of } t
\tag{2.86}
\end{equation}

This is in **dramatic contrast** with the Schrödinger picture state ket:

\textbf{Schrödinger picture:}
\begin{equation}
|\alpha, t_0 = 0; t\rangle_S = \mathscr{U}(t)|\alpha, t_0 = 0\rangle
\tag{2.87}
\end{equation}

\subsubsection{Equivalence of Expectation Values}

Expectation value $\langle A \rangle$ is obviously the same in both pictures:
\begin{align}
&{}_S\langle\alpha, t_0 = 0; t|A^{(S)}|\alpha, t_0 = 0; t\rangle_S \notag \\
&\quad = \langle\alpha, t_0 = 0|\mathscr{U}^\dagger A^{(S)}\mathscr{U}|\alpha, t_0 = 0\rangle \notag \\
&\quad = {}_H\langle\alpha, t_0 = 0; t|A^{(H)}(t)|\alpha, t_0 = 0; t\rangle_H
\tag{2.88}
\end{align}

\textbf{Key insight:} Both pictures give identical physical predictions (expectation values). The choice is one of mathematical convenience.

\textbf{Summary of differences:}

\begin{center}
\begin{tabular}{|l|c|c|}
\hline
& \textbf{Schrödinger Picture} & \textbf{Heisenberg Picture} \\
\hline
State kets & Time-dependent & Fixed (frozen at $t_0$) \\
& $|\alpha; t\rangle_S = \mathscr{U}(t)|\alpha\rangle$ & $|\alpha; t\rangle_H = |\alpha\rangle$ \\
\hline
Operators & Time-independent & Time-dependent \\
& $A^{(S)}$ & $A^{(H)}(t) = \mathscr{U}^\dagger(t)A^{(S)}\mathscr{U}(t)$ \\
\hline
Expectation values & \multicolumn{2}{c|}{Same in both pictures} \\
\hline
\end{tabular}
\end{center}

\textbf{Physical interpretation:}
\begin{itemize}
\item \textbf{Schrödinger:} "Active" transformation - state evolves, observables fixed
\item \textbf{Heisenberg:} "Passive" transformation - state fixed, observables evolve
\item Closer to classical mechanics where observables change with time
\item Both are equivalent formulations - matter of convenience
\end{itemize}

\subsection{The Heisenberg Equation of Motion}

\textbf{Derivation:} Assuming $A^{(S)}$ does not depend explicitly on time (most physical situations), differentiate (2.84):

\begin{align}
\frac{dA^{(H)}}{dt} &= \frac{\partial\mathscr{U}^\dagger}{\partial t}A^{(S)}\mathscr{U} + \mathscr{U}^\dagger A^{(S)}\frac{\partial\mathscr{U}}{\partial t} \notag \\
&= -\frac{1}{i\hbar}\mathscr{U}^\dagger H\mathscr{U}\mathscr{U}^\dagger A^{(S)}\mathscr{U} + \frac{1}{i\hbar}\mathscr{U}^\dagger A^{(S)}\mathscr{U}\mathscr{U}^\dagger H\mathscr{U} \notag \\
&= \frac{1}{i\hbar}[A^{(H)}, \mathscr{U}^\dagger H\mathscr{U}]
\tag{2.89}
\end{align}

where we used (2.25):
\begin{align}
\frac{\partial\mathscr{U}}{\partial t} &= \frac{1}{i\hbar}H\mathscr{U} \tag{2.90a} \\
\frac{\partial\mathscr{U}^\dagger}{\partial t} &= -\frac{1}{i\hbar}\mathscr{U}^\dagger H \tag{2.90b}
\end{align}

\textbf{Heisenberg picture Hamiltonian:}

Since $H$ was introduced in Schrödinger picture, we might define:
\begin{equation}
H^{(H)} = \mathscr{U}^\dagger H\mathscr{U}
\tag{2.91}
\end{equation}

But in elementary applications where $\mathscr{U}$ is given by (2.83), $\mathscr{U}$ and $H$ commute:
\begin{equation}
\mathscr{U}^\dagger H\mathscr{U} = H
\tag{2.92}
\end{equation}

Therefore (2.89) becomes:
\begin{equation}
\boxed{\frac{dA^{(H)}}{dt} = \frac{1}{i\hbar}[A^{(H)}, H]}
\tag{2.93}
\end{equation}

This is the **Heisenberg equation of motion**. Derived using properties of time-evolution operator and defining equation for $A^{(H)}$.

\subsubsection{Comparison with Classical Mechanics}

Compare (2.93) with classical equation of motion in Poisson bracket form. In classical physics, for function $A$ of $q$ and $p$ not explicitly time-dependent (Goldstein et al. (2002), pp. 396-397):
\begin{equation}
\frac{dA}{dt} = [A, H]_{\text{classical}}
\tag{2.94}
\end{equation}

**Again, Dirac's quantization rule leads to correct quantum equation!**

Historically (2.93) was first written by P. A. M. Dirac, who called it the Heisenberg equation of motion.

\textbf{Important note:} (2.93) makes sense whether or not $A^{(H)}$ has a classical analogue. For example, spin operator in Heisenberg picture satisfies:
\begin{equation}
\frac{dS_i^{(H)}}{dt} = \frac{1}{i\hbar}[S_i^{(H)}, H]
\tag{2.95}
\end{equation}

This can be used to discuss spin precession, but has no classical counterpart because $S_z$ **cannot** be written as function of $q$ and $p$.

\textbf{General principle:} For quantities possessing classical counterparts, correct classical equation can be obtained from corresponding quantum-mechanical equation via:
\begin{equation}
\frac{[,]}{i\hbar} \to [,]_{\text{classical}}
\tag{2.96}
\end{equation}

Classical mechanics can be derived from quantum mechanics, but the opposite is not true.

\subsection{Free Particles: Ehrenfest's Theorem}

Whether in Schrödinger or Heisenberg picture, to use equations of motion we must first construct the appropriate Hamiltonian operator.

\textbf{For systems with classical analogues:} Assume Hamiltonian has same form as classical physics; replace classical $x_i$ and $p_i$ by corresponding quantum operators. This reproduces correct classical equations in classical limit.

\textbf{Ordering ambiguity:} When noncommuting observables appear, resolve by requiring $H$ to be Hermitian. For example, quantum analogue of classical product $xp$ is $\frac{1}{2}(xp + px)$.

\textbf{For systems without classical analogues:} Can only guess structure of Hamiltonian, then test against experiment.

\subsubsection{Useful Commutator Formulas}

For practical applications, evaluate commutator of $x_i$ (or $p_i$) with functions of $x_j$ and $p_j$:
\begin{equation}
[x_i, F(\mathbf{p})] = i\hbar\frac{\partial F}{\partial p_i}
\tag{2.97a}
\end{equation}

\begin{equation}
[p_i, G(\mathbf{x})] = -i\hbar\frac{\partial G}{\partial x_i}
\tag{2.97b}
\end{equation}

where $F$ and $G$ are functions that can be expanded in powers of $p_j$ and $x_j$, respectively. Easily proved by repeatedly applying (1.232e).

\subsubsection{Free Particle}

Apply Heisenberg equation of motion to free particle of mass $m$. Hamiltonian (same form as classical):
\begin{equation}
H = \frac{\mathbf{p}^2}{2m} = \frac{(p_x^2 + p_y^2 + p_z^2)}{2m}
\tag{2.98}
\end{equation}

Observables $p_i$ and $x_i$ understood to be momentum and position operators in Heisenberg picture (though we omit superscript $(H)$).

\textbf{Momentum evolution:}

Since $p_i$ commutes with any function of $p_j$:
\begin{equation}
\frac{dp_i}{dt} = \frac{1}{i\hbar}[p_i, H] = 0
\tag{2.99}
\end{equation}

For free particle, momentum operator is **constant of motion**: $p_i(t) = p_i(0)$ at all times.

Generally, from (2.93): whenever $A^{(H)}$ commutes with Hamiltonian, $A^{(H)}$ is constant of motion.

\textbf{Position evolution:}
\begin{align}
\frac{dx_i}{dt} &= \frac{1}{i\hbar}[x_i, H] = \frac{1}{i\hbar}\frac{1}{2m}i\hbar\frac{\partial}{\partial p_i}\left(\sum_{j=1}^3 p_j^2\right) \notag \\
&= \frac{p_i}{m} = \frac{p_i(0)}{m}
\tag{2.100}
\end{align}

where we used (2.97a). Solution:
\begin{equation}
x_i(t) = x_i(0) + \left(\frac{p_i(0)}{m}\right)t
\tag{2.101}
\end{equation}

Reminiscent of classical trajectory equation for uniform rectilinear motion!

\textbf{Important note:} Even though
\begin{equation}
[x_i(0), x_j(0)] = 0
\tag{2.102}
\end{equation}

at equal times, commutator of $x_i$ at **different** times does **not** vanish:
\begin{equation}
[x_i(t), x_i(0)] = \left[\frac{p_i(0)t}{m}, x_i(0)\right] = \frac{-i\hbar t}{m}
\tag{2.103}
\end{equation}

Applying uncertainty relation (1.146):
\begin{equation}
\langle(\Delta x_i)^2\rangle_t\langle(\Delta x_i)^2\rangle_{t=0} \geq \frac{\hbar^2 t^2}{4m^2}
\tag{2.104}
\end{equation}

Even if particle well localized at $t = 0$, position becomes more uncertain with time (wave packet spreading in wave mechanics).

\subsubsection{Free Particle with Potential}

Add potential $V(\mathbf{x})$ to free-particle Hamiltonian:
\begin{equation}
H = \frac{\mathbf{p}^2}{2m} + V(\mathbf{x})
\tag{2.105}
\end{equation}

Here $V(\mathbf{x})$ is function of $x$-, $y$-, $z$-**operators**. Using (2.97b):
\begin{equation}
\frac{dp_i}{dt} = \frac{1}{i\hbar}[p_i, V(\mathbf{x})] = -\frac{\partial}{\partial x_i}V(\mathbf{x})
\tag{2.106}
\end{equation}

On the other hand:
\begin{equation}
\frac{dx_i}{dt} = \frac{p_i}{m}
\tag{2.107}
\end{equation}

still holds because $x_i$ commutes with newly added term $V(\mathbf{x})$. Using Heisenberg equation again:
\begin{equation}
\frac{d^2x_i}{dt^2} = \frac{1}{i\hbar}\left[\frac{dx_i}{dt}, H\right] = \frac{1}{i\hbar}\left[\frac{p_i}{m}, H\right] = \frac{1}{m}\frac{dp_i}{dt}
\tag{2.108}
\end{equation}

Combining with (2.106):
\begin{equation}
\boxed{m\frac{d^2\mathbf{x}}{dt^2} = -\nabla V(\mathbf{x})}
\tag{2.109}
\end{equation}

This is the **quantum-mechanical analogue of Newton's second law**.

\textbf{Ehrenfest's Theorem:}

Taking expectation values of both sides with respect to Heisenberg state ket (does **not** move with time):
\begin{equation}
\boxed{m\frac{d^2}{dt^2}\langle\mathbf{x}\rangle = \frac{d\langle\mathbf{p}\rangle}{dt} = -\langle\nabla V(\mathbf{x})\rangle}
\tag{2.110}
\end{equation}

This is **Ehrenfest's theorem** (P. Ehrenfest, 1927) using wave mechanics formalism.

\textbf{Key points:}
- Written in expectation form, validity independent of picture choice
- Expectation values are same in both pictures
- Operator form (2.109) meaningful only if $\mathbf{x}$ and $\mathbf{p}$ are Heisenberg picture operators
- The $\hbar$ have completely disappeared!
- Not surprising that center of wave packet moves like **classical** particle subjected to $V(\mathbf{x})$

\subsection{Base Kets and Transition Amplitudes}

\textbf{Common misconception:} All kets move in Schrödinger picture and are stationary in Heisenberg picture. This is **not** the case!

\textbf{Important distinction:} Behavior of **state kets** vs **base kets**.

\subsubsection{Base Kets in Schrödinger Picture}

Started discussion in Section 1.2: eigenkets of observables are used as base kets. What happens to eigenvalue equation:
\begin{equation}
A|a'\rangle = a'|a'\rangle
\tag{2.111}
\end{equation}

with time?

In Schrödinger picture, $A$ does not change. Base kets obtained as solutions to this eigenvalue equation at $t = 0$ must remain unchanged. **Unlike state kets, base kets do NOT change in Schrödinger picture.**

\subsubsection{Base Kets in Heisenberg Picture}

Very different situation. Eigenvalue equation is for time-dependent operator:
\begin{equation}
A^{(H)}(t) = \mathscr{U}^\dagger A(0)\mathscr{U}
\tag{2.112}
\end{equation}

From (2.111) evaluated at $t = 0$ (when two pictures coincide):
\begin{equation}
\mathscr{U}^\dagger A(0)\mathscr{U}\mathscr{U}^\dagger|a'\rangle = a'\mathscr{U}^\dagger|a'\rangle
\tag{2.113}
\end{equation}

which implies eigenvalue equation for $A^{(H)}$:
\begin{equation}
A^{(H)}(\mathscr{U}^\dagger|a'\rangle) = a'(\mathscr{U}^\dagger|a'\rangle)
\tag{2.114}
\end{equation}

If eigenkets of observables form base kets, then $\{\mathscr{U}^\dagger|a'\rangle\}$ must be used as base kets in Heisenberg picture. As time goes on, Heisenberg picture base kets, denoted $|a', t\rangle_H$, move:
\begin{equation}
|a', t\rangle_H = \mathscr{U}^\dagger|a'\rangle
\tag{2.115}
\end{equation}

Because of $\mathscr{U}^\dagger$ rather than $\mathscr{U}$ in (2.115), Heisenberg picture base kets **rotate oppositely** when compared with Schrödinger picture state kets.

Specifically, $|a', t\rangle_H$ satisfies "wrong-sign Schrödinger equation":
\begin{equation}
i\hbar\frac{\partial}{\partial t}|a', t\rangle_H = -H|a', t\rangle_H
\tag{2.116}
\end{equation}

\textbf{Eigenvalues:} Unchanged with time from (2.114). Consistent with theorem on unitary equivalent observables (Section 1.5).

\textbf{Expansion of $A^{(H)}(t)$:}
\begin{align}
A^{(H)}(t) &= \sum_{a'}|a', t\rangle_H a'{}_{H}\langle a', t| \notag \\
&= \sum_{a'}\mathscr{U}^\dagger|a'\rangle a'\langle a'|\mathscr{U} \notag \\
&= \mathscr{U}^\dagger A^{(S)}\mathscr{U}
\tag{2.117}
\end{align}

Everything consistent provided Heisenberg base kets change as in (2.115).

\subsubsection{Expansion Coefficients}

Expansion coefficients of state ket in terms of base kets are **same in both pictures**:

\textbf{Schrödinger picture:}
\begin{equation}
c_{a'}(t) = \underbrace{\langle a'|}_{\text{base bra}} \cdot \underbrace{(\mathscr{U}|\alpha, t_0 = 0\rangle)}_{\text{state ket}}
\tag{2.118a}
\end{equation}

\textbf{Heisenberg picture:}
\begin{equation}
c_{a'}(t) = \underbrace{(\langle a'|\mathscr{U})}_{\text{base bra}} \cdot \underbrace{|\alpha, t_0 = 0\rangle}_{\text{state ket}}
\tag{2.118b}
\end{equation}

**Pictorially:** Cosine of angle between state ket and base ket is same whether we rotate state ket counterclockwise or base ket clockwise.

\textbf{Wave function interpretation:} These considerations apply to continuous spectrum. Wave function $\langle\mathbf{x}'|\alpha\rangle$ can be regarded either as:
1. Inner product of stationary position eigenbra with moving state ket (Schrödinger picture)
2. Inner product of moving position eigenbra with stationary state ket (Heisenberg picture)

Time dependence of wave function discussed in Section 2.4 (Schrödinger wave equation).

\subsubsection{Transition Amplitudes}

To illustrate further equivalence between pictures, study **transition amplitudes** (fundamental role in Section 2.6).

Suppose system prepared at $t = 0$ to be in eigenstate of observable $A$ with eigenvalue $a'$. At later time $t$: What is probability amplitude (transition amplitude) for system to be found in eigenstate of observable $B$ with eigenvalue $b'$? (Here $A$ and $B$ can be same or different.)

\textbf{Schrödinger picture:} State ket at $t$ given by $\mathscr{U}|a'\rangle$. Base kets $|a'\rangle$ and $|b'\rangle$ do not vary with time:
\begin{equation}
\underbrace{\langle b'|}_{\text{base bra}} \cdot \underbrace{(\mathscr{U}|a'\rangle)}_{\text{state ket}}
\tag{2.119}
\end{equation}

\textbf{Heisenberg picture:} State ket is stationary (remains $|a'\rangle$ at all times). Base kets evolve oppositely. Transition amplitude:
\begin{equation}
\underbrace{(\langle b'|\mathscr{U})}_{\text{base bra}} \cdot \underbrace{|a'\rangle}_{\text{state ket}}
\tag{2.120}
\end{equation}

Obviously (2.119) and (2.120) are the same. Both written as:
\begin{equation}
\langle b'|\mathscr{U}(t, 0)|a'\rangle
\tag{2.121}
\end{equation}

In some loose sense, this is transition amplitude for "going" from state $|a'\rangle$ to state $|b'\rangle$.

\subsubsection{Summary Table}

\begin{center}
\begin{tabular}{|l|c|c|}
\hline
\textbf{Table 2.1} & \multicolumn{2}{c|}{\textbf{The Schrödinger Picture Versus the Heisenberg Picture}} \\
\hline
& \textbf{Schrödinger picture} & \textbf{Heisenberg picture} \\
\hline
State ket & Moving: (2.5), (2.27) & Stationary \\
Observable & Stationary & Moving: (2.84), (2.93) \\
Base ket & Stationary & Moving oppositely: (2.115), (2.116) \\
\hline
\end{tabular}
\end{center}

\textbf{Key takeaway:} The pictures are completely equivalent formulations. Choice is matter of convenience for the problem at hand.

% ============================================
\newpage
\section{Simple Harmonic Oscillator}

\textbf{Importance:} The simple harmonic oscillator is one of the most important problems in quantum mechanics. It:
\begin{itemize}
\item Illustrates many basic concepts and methods of quantum mechanics
\item Has much practical value - essentially any potential well can be approximated by a simple harmonic oscillator
\item Describes phenomena from molecular vibrations to nuclear structure
\item Is an important starting point for quantum field theory (Hamiltonian is sum of squares of two canonically conjugate variables)
\end{itemize}

\subsection{Energy Eigenkets and Energy Eigenvalues}

We use Dirac's elegant operator method (based on work of M. Born and N. Wiener) to obtain energy eigenkets and eigenvalues.

\textbf{Hamiltonian:}
\begin{equation}
H = \frac{p^2}{2m} + \frac{m\omega^2x^2}{2}
\tag{2.122}
\end{equation}

where $\omega$ is the angular frequency of the classical oscillator related to spring constant $k$ via $\omega = \sqrt{k/m}$. The operators $x$ and $p$ are Hermitian.

\subsubsection{Creation and Annihilation Operators}

\textbf{Define two non-Hermitian operators:}
\begin{equation}
a = \sqrt{\frac{m\omega}{2\hbar}}\left(x + \frac{ip}{m\omega}\right), \quad a^\dagger = \sqrt{\frac{m\omega}{2\hbar}}\left(x - \frac{ip}{m\omega}\right)
\tag{2.123}
\end{equation}

known as the \textbf{annihilation operator} and \textbf{creation operator}, respectively.

\textbf{Commutation relation:}
Using the canonical commutation relations:
\begin{equation}
[a, a^\dagger] = \left(\frac{1}{2\hbar}\right)(-i[x,p] + i[p,x]) = 1
\tag{2.124}
\end{equation}

\subsubsection{Number Operator}

\textbf{Definition:}
\begin{equation}
N = a^\dagger a
\tag{2.125}
\end{equation}

which is obviously Hermitian.

\textbf{Relation to Hamiltonian:}
Straightforward calculation shows:
\begin{align}
a^\dagger a &= \left(\frac{m\omega}{2\hbar}\right)\left(x^2 + \frac{p^2}{m^2\omega^2}\right) + \left(\frac{i}{2\hbar}\right)[x,p] \notag \\
&= \frac{H}{\hbar\omega} - \frac{1}{2}
\tag{2.126}
\end{align}

Therefore:
\begin{equation}
\boxed{H = \hbar\omega\left(N + \frac{1}{2}\right)}
\tag{2.127}
\end{equation}

Since $H$ is a linear function of $N$, $N$ can be diagonalized simultaneously with $H$.

\subsubsection{Energy Eigenvalue Equation}

Denote an energy eigenket of $N$ by its eigenvalue $n$:
\begin{equation}
N|n\rangle = n|n\rangle
\tag{2.128}
\end{equation}

From (2.127):
\begin{equation}
H|n\rangle = \left(n + \frac{1}{2}\right)\hbar\omega|n\rangle
\tag{2.129}
\end{equation}

\textbf{Energy eigenvalues:}
\begin{equation}
\boxed{E_n = \left(n + \frac{1}{2}\right)\hbar\omega}
\tag{2.130}
\end{equation}

We will show that $n$ must be a nonnegative integer.

\subsubsection{Ladder Operators}

\textbf{Commutators with $N$:}
\begin{align}
[N, a] &= [a^\dagger a, a] = a^\dagger[a,a] + [a^\dagger, a]a = -a \tag{2.131} \\
[N, a^\dagger] &= a^\dagger \tag{2.132}
\end{align}

\textbf{Action on eigenkets:}
\begin{align}
Na^\dagger|n\rangle &= ([N,a^\dagger] + a^\dagger N)|n\rangle = (n+1)a^\dagger|n\rangle \tag{2.133a} \\
Na|n\rangle &= ([N,a] + aN)|n\rangle = (n-1)a|n\rangle \tag{2.133b}
\end{align}

$a^\dagger|n\rangle$ ($a|n\rangle$) is also an eigenket of $N$ with eigenvalue increased (decreased) by one. The increase (decrease) of $n$ by one corresponds to creation (annihilation) of one quantum unit of energy $\hbar\omega$.

\subsubsection{Normalization Constants}

From (2.133b), $a|n\rangle$ and $|n-1\rangle$ are the same up to a multiplicative constant:
\begin{equation}
a|n\rangle = c|n-1\rangle
\tag{2.134}
\end{equation}

\textbf{Determine $c$ from normalization:}
\begin{equation}
\langle n|a^\dagger a|n\rangle = |c|^2
\tag{2.135}
\end{equation}

Since $a^\dagger a$ is the number operator:
\begin{equation}
n = |c|^2
\tag{2.136}
\end{equation}

Taking $c$ real and positive by convention:
\begin{equation}
\boxed{a|n\rangle = \sqrt{n}|n-1\rangle}
\tag{2.137}
\end{equation}

Similarly:
\begin{equation}
\boxed{a^\dagger|n\rangle = \sqrt{n+1}|n+1\rangle}
\tag{2.138}
\end{equation}

\subsubsection{Quantization of $n$}

Repeatedly applying annihilation operator:
\begin{align}
a^2|n\rangle &= \sqrt{n(n-1)}|n-2\rangle \notag \\
a^3|n\rangle &= \sqrt{n(n-1)(n-2)}|n-3\rangle \notag \\
&\vdots
\tag{2.139}
\end{align}

We obtain eigenkets with smaller and smaller $n$ until the sequence terminates. This must happen for positive integer $n$.

\textbf{Positivity requirement:}
\begin{equation}
n = \langle n|N|n\rangle = \langle n|a^\dagger a|n\rangle = (\langle n|a^\dagger) \cdot (a|n\rangle) \geq 0
\tag{2.140}
\end{equation}

Therefore $n$ can never be negative! The sequence must terminate with $n = 0$ and the allowed values of $n$ are \textbf{nonnegative integers}.

\subsubsection{Ground State and Excited States}

\textbf{Ground state energy:}
\begin{equation}
\boxed{E_0 = \frac{1}{2}\hbar\omega}
\tag{2.141}
\end{equation}

\textbf{Constructing energy eigenstates:}
Apply creation operator $a^\dagger$ successively to ground state $|0\rangle$:
\begin{align}
|1\rangle &= a^\dagger|0\rangle \notag \\
|2\rangle &= \left(\frac{a^\dagger}{\sqrt{2}}\right)|1\rangle = \left[\frac{(a^\dagger)^2}{\sqrt{2}}\right]|0\rangle \notag \\
|3\rangle &= \left(\frac{a^\dagger}{\sqrt{3}}\right)|2\rangle = \left[\frac{(a^\dagger)^3}{\sqrt{3!}}\right]|0\rangle \notag \\
&\vdots \notag \\
|n\rangle &= \left[\frac{(a^\dagger)^n}{\sqrt{n!}}\right]|0\rangle
\tag{2.142}
\end{align}

\textbf{Energy eigenvalues:}
\begin{equation}
\boxed{E_n = \left(n + \frac{1}{2}\right)\hbar\omega \quad (n = 0,1,2,3,\ldots)}
\tag{2.143}
\end{equation}

\subsubsection{Matrix Elements}

From (2.137), (2.138), and orthonormality:
\begin{equation}
\langle n'|a|n\rangle = \sqrt{n}\delta_{n',n-1}, \quad \langle n'|a^\dagger|n\rangle = \sqrt{n+1}\delta_{n',n+1}
\tag{2.144}
\end{equation}

\textbf{Position and momentum operators:}
Using
\begin{equation}
x = \sqrt{\frac{\hbar}{2m\omega}}(a + a^\dagger), \quad p = i\sqrt{\frac{m\hbar\omega}{2}}(-a + a^\dagger)
\tag{2.145}
\end{equation}

we derive:
\begin{align}
\langle n'|x|n\rangle &= \sqrt{\frac{\hbar}{2m\omega}}(\sqrt{n}\delta_{n',n-1} + \sqrt{n+1}\delta_{n',n+1}) \tag{2.146a} \\
\langle n'|p|n\rangle &= i\sqrt{\frac{m\hbar\omega}{2}}(-\sqrt{n}\delta_{n',n-1} + \sqrt{n+1}\delta_{n',n+1}) \tag{2.146b}
\end{align}

Neither $x$ nor $p$ is diagonal in the $N$-representation (not surprising since $x$, $p$, like $a$ and $a^\dagger$, do not commute with $N$).

\subsection{Wave Functions}

\subsubsection{Ground State Wave Function}

Start with ground state defined by:
\begin{equation}
a|0\rangle = 0
\tag{2.147}
\end{equation}

In $x$-representation:
\begin{equation}
\langle x'|a|0\rangle = \sqrt{\frac{m\omega}{2\hbar}}\langle x'|\left(x + \frac{ip}{m\omega}\right)|0\rangle = 0
\tag{2.148}
\end{equation}

This is a differential equation for the ground-state wave function $\langle x'|0\rangle$:
\begin{equation}
\left(x' + x_0^2\frac{d}{dx'}\right)\langle x'|0\rangle = 0
\tag{2.149}
\end{equation}

where the \textbf{characteristic length scale} is:
\begin{equation}
\boxed{x_0 \equiv \sqrt{\frac{\hbar}{m\omega}}}
\tag{2.150}
\end{equation}

\textbf{Normalized solution:}
\begin{equation}
\boxed{\langle x'|0\rangle = \left(\frac{1}{\pi^{1/4}\sqrt{x_0}}\right)\exp\left[-\frac{1}{2}\left(\frac{x'}{x_0}\right)^2\right]}
\tag{2.151}
\end{equation}

\subsubsection{Excited State Wave Functions}

Obtain energy eigenfunctions for excited states by evaluating:
\begin{align}
\langle x'|1\rangle &= \langle x'|a^\dagger|0\rangle = \left(\frac{1}{\sqrt{2x_0}}\right)\left(x' - x_0^2\frac{d}{dx'}\right)\langle x'|0\rangle \notag \\
\langle x'|2\rangle &= \left(\frac{1}{\sqrt{2}}\right)\langle x'|(a^\dagger)^2|0\rangle = \left(\frac{1}{\sqrt{2!}}\right)\left(\frac{1}{\sqrt{2x_0}}\right)^2\left(x' - x_0^2\frac{d}{dx'}\right)^2\langle x'|0\rangle,\ldots
\tag{2.152}
\end{align}

\textbf{General formula:}
\begin{equation}
\boxed{\langle x'|n\rangle = \left(\frac{1}{\pi^{1/4}\sqrt{2^n n!}}\right)\left(\frac{1}{x_0^{n+1/2}}\right)\left(x' - x_0^2\frac{d}{dx'}\right)^n\exp\left[-\frac{1}{2}\left(\frac{x'}{x_0}\right)^2\right]}
\tag{2.153}
\end{equation}

\subsection{Expectation Values and Uncertainty Relations}

\textbf{Position squared operator:}
\begin{equation}
x^2 = \left(\frac{\hbar}{2m\omega}\right)(a^2 + a^{\dagger 2} + a^\dagger a + aa^\dagger)
\tag{2.154}
\end{equation}

\textbf{Ground state expectation values:}
For the ground state, only the last term contributes:
\begin{align}
\langle x^2\rangle &= \frac{\hbar}{2m\omega} = \frac{x_0^2}{2} \tag{2.155} \\
\langle p^2\rangle &= \frac{\hbar m\omega}{2} \tag{2.156}
\end{align}

\textbf{Energy equipartition (virial theorem):}
\begin{equation}
\left\langle\frac{p^2}{2m}\right\rangle = \frac{\hbar\omega}{4} = \frac{\langle H\rangle}{2} \quad \text{and} \quad \left\langle\frac{m\omega^2x^2}{2}\right\rangle = \frac{\hbar\omega}{4} = \frac{\langle H\rangle}{2}
\tag{2.157}
\end{equation}

From (2.146):
\begin{equation}
\langle x\rangle = \langle p\rangle = 0
\tag{2.158}
\end{equation}

which also holds for excited states.

\textbf{Uncertainties:}
\begin{equation}
\langle(\Delta x)^2\rangle = \langle x^2\rangle = \frac{\hbar}{2m\omega} \quad \text{and} \quad \langle(\Delta p)^2\rangle = \langle p^2\rangle = \frac{\hbar m\omega}{2}
\tag{2.159}
\end{equation}

\textbf{Minimum uncertainty product:}
\begin{equation}
\boxed{\langle(\Delta x)^2\rangle\langle(\Delta p)^2\rangle = \frac{\hbar^2}{4}}
\tag{2.160}
\end{equation}

The ground-state wave function has Gaussian shape, saturating the uncertainty relation.

\textbf{Excited states:}
Uncertainty products are larger:
\begin{equation}
\langle(\Delta x)^2\rangle\langle(\Delta p)^2\rangle = \left(n + \frac{1}{2}\right)^2\hbar^2
\tag{2.161}
\end{equation}

\subsection{Time Development of the Oscillator}

So far we have not discussed time evolution of oscillator state kets or observables. Everything holds at some instant, say $t = 0$; operators $x$, $p$, $a$, and $a^\dagger$ are either Schrödinger picture operators (at all $t$) or Heisenberg picture operators at $t = 0$.

In the remaining part, we work exclusively in the \textbf{Heisenberg picture}, meaning $x$, $p$, $a$, and $a^\dagger$ are all time dependent even though we don't explicitly write $x^{(H)}(t)$.

\subsubsection{Heisenberg Equations of Motion}

From (2.106) and (2.107):
\begin{align}
\frac{dp}{dt} &= -m\omega^2 x \tag{2.162a} \\
\frac{dx}{dt} &= \frac{p}{m} \tag{2.162b}
\end{align}

These coupled equations are equivalent to two uncoupled equations for $a$ and $a^\dagger$:
\begin{align}
\frac{da}{dt} &= \sqrt{\frac{m\omega}{2\hbar}}\left(\frac{p}{m} - i\omega x\right) = -i\omega a \tag{2.163a} \\
\frac{da^\dagger}{dt} &= i\omega a^\dagger \tag{2.163b}
\end{align}

\textbf{Solutions:}
\begin{equation}
\boxed{a(t) = a(0)\exp(-i\omega t) \quad \text{and} \quad a^\dagger(t) = a^\dagger(0)\exp(i\omega t)}
\tag{2.164}
\end{equation}

These relations explicitly show that $N$ and $H$ are \textit{time-independent} operators even in Heisenberg picture, as they must be.

\textbf{In terms of $x$ and $p$:}
\begin{align}
x(t) + \frac{ip(t)}{m\omega} &= x(0)\exp(-i\omega t) + i\left[\frac{p(0)}{m\omega}\right]\exp(-i\omega t) \notag \\
x(t) - \frac{ip(t)}{m\omega} &= x(0)\exp(i\omega t) - i\left[\frac{p(0)}{m\omega}\right]\exp(i\omega t)
\tag{2.165}
\end{align}

Equating Hermitian and anti-Hermitian parts:
\begin{align}
x(t) &= x(0)\cos\omega t + \left[\frac{p(0)}{m\omega}\right]\sin\omega t \tag{2.166a} \\
p(t) &= -m\omega x(0)\sin\omega t + p(0)\cos\omega t \tag{2.166b}
\end{align}

These look the same as classical equations of motion. The $x$ and $p$ operators "oscillate" just like their classical analogues.

\subsubsection{Alternative Derivation: Baker-Hausdorff Lemma}

For pedagogical reasons, we present an alternative derivation of (2.166a) by evaluating:
\begin{equation}
x(t) = \exp\left(\frac{iHt}{\hbar}\right)x(0)\exp\left(\frac{-iHt}{\hbar}\right)
\tag{2.167}
\end{equation}

\textbf{Baker-Hausdorff lemma:}
\begin{equation}
\exp(iG\lambda)A\exp(-iG\lambda) = A + i\lambda[G,A] + \left(\frac{i^2\lambda^2}{2!}\right)[G,[G,A]] + \cdots + \left(\frac{i^n\lambda^n}{n!}\right)[G,[G,[G,\ldots[G,A]]]\ldots] + \cdots
\tag{2.168}
\end{equation}

where $G$ is Hermitian and $\lambda$ is real parameter.

Applying to (2.167):
\begin{equation}
\exp\left(\frac{iHt}{\hbar}\right)x(0)\exp\left(\frac{-iHt}{\hbar}\right) = x(0) + \left(\frac{it}{\hbar}\right)[H,x(0)] + \left(\frac{i^2t^2}{2!\hbar^2}\right)[H,[H,x(0)]] + \cdots
\tag{2.169}
\end{equation}

Each term can be reduced to $x$ or $p$ using:
\begin{align}
[H,x(0)] &= \frac{-i\hbar p(0)}{m} \tag{2.170a} \\
[H,p(0)] &= i\hbar m\omega^2 x(0) \tag{2.170b}
\end{align}

Thus:
\begin{align}
\exp\left(\frac{iHt}{\hbar}\right)x(0)\exp\left(\frac{-iHt}{\hbar}\right) &= x(0) + \left[\frac{p(0)}{m}\right]t - \left(\frac{1}{2!}\right)t^2\omega^2 x(0) \notag \\
&\quad - \left(\frac{1}{3!}\right)\frac{t^3\omega^2 p(0)}{m} + \cdots \notag \\
&= x(0)\cos\omega t + \left[\frac{p(0)}{m\omega}\right]\sin\omega t
\tag{2.171}
\end{align}

in agreement with (2.166a).

\subsubsection{Classical vs Quantum Oscillations}

From (2.166), one may be tempted to conclude that $\langle x\rangle$ and $\langle p\rangle$ always oscillate with angular frequency $\omega$. However, this is \textbf{not correct}.

For any energy eigenstate characterized by definite $n$, the expectation value $\langle n|x(t)|n\rangle$ vanishes because operators $x(0)$ and $p(0)$ change $n$ by $\pm 1$ and $|n\rangle$ and $|n\pm 1\rangle$ are orthogonal.

This is consistent with earlier conclusion (Section 2.1) that expectation value of an observable taken with respect to a stationary state does not vary with time.

\textbf{To observe classical-like oscillations:} Must look at a \textit{superposition} of energy eigenstates such as:
\begin{equation}
|\alpha\rangle = c_0|0\rangle + c_1|1\rangle
\tag{2.172}
\end{equation}

The expectation value of $x(t)$ taken with respect to (2.172) does oscillate.

\subsection{Coherent States}

We have seen that an energy eigenstate does not behave like the classical oscillator - in the sense of oscillating expectation values for $x$ and $p$ - no matter how large $n$ may be.

\textbf{Question:} How can we construct a superposition of energy eigenstates that most closely imitates the classical oscillator? In wave function language, we want a wave packet that bounces back and forth without spreading in shape.

\textbf{Answer:} A \textit{coherent state} defined by the eigenvalue equation for the non-Hermitian annihilation operator $a$:
\begin{equation}
a|\lambda\rangle = \lambda|\lambda\rangle
\tag{2.173}
\end{equation}

with complex eigenvalue $\lambda$ does the desired job. The coherent state has many remarkable properties:

\textbf{1. Poisson distribution:}
When expressed as superposition of energy eigenstates:
\begin{equation}
|\lambda\rangle = \sum_{n=0}^{\infty}f(n)|n\rangle
\tag{2.174}
\end{equation}

the distribution of $|f(n)|^2$ with respect to $n$ is Poisson type about mean value $\bar{n}$:
\begin{equation}
|f(n)|^2 = \left(\frac{\bar{n}^n}{n!}\right)\exp(-\bar{n})
\tag{2.175}
\end{equation}

\textbf{2.} Can be obtained by translating oscillator ground state by some finite distance.

\textbf{3.} Satisfies minimum uncertainty product relation at all times.

A systematic study of coherent states, pioneered by R. Glauber, is very rewarding; the reader is urged to work out Exercise 2.21 on this subject.

% Add this after the Coherent States section (after equation 2.175)

\newpage
\section{Schrödinger's Wave Equation}

\subsection{Time-Dependent Wave Equation}

We now turn to the Schrödinger picture and examine the time evolution of $|\alpha, t_0; t\rangle$ in the $\mathbf{x}$-representation. In other words, our task is to study the behavior of the wave function

\begin{equation}
\boxed{\psi(\mathbf{x}', t) = \langle \mathbf{x}'|\alpha, t_0; t\rangle}
\tag{2.176}
\end{equation}

as a function of time, where $|\alpha, t_0; t\rangle$ is a state ket in the Schrödinger picture at time $t$, and $\langle \mathbf{x}'|$ is a time-independent position eigenbra with eigenvalue $\mathbf{x}'$.

\textbf{Hamiltonian operator:}
\begin{equation}
H = \frac{\mathbf{p}^2}{2m} + V(\mathbf{x})
\tag{2.177}
\end{equation}

The potential $V(\mathbf{x})$ is a Hermitian operator; it is also local in the sense that in the $\mathbf{x}$-representation we have

\begin{equation}
\langle \mathbf{x}''|V(\mathbf{x})|\mathbf{x}'\rangle = V(\mathbf{x}')\delta^3(\mathbf{x}' - \mathbf{x}'')
\tag{2.178}
\end{equation}

where $V(\mathbf{x}')$ is a real function of $\mathbf{x}'$. Later we will consider more complicated Hamiltonians: a time-dependent potential $V(\mathbf{x}, t)$; a nonlocal but separable potential where the right-hand side of (2.178) is replaced by $v_1(\mathbf{x}'')v_2(\mathbf{x}')$; a momentum-dependent interaction of the form $\mathbf{p} \cdot \mathbf{A} + \mathbf{A} \cdot \mathbf{p}$, where $\mathbf{A}$ is the vector potential in electrodynamics; and so on.

\subsubsection{Derivation of the Wave Equation}

We now derive Schrödinger's time-dependent wave equation. We first write the Schrödinger equation for a state ket (2.27) in the $\mathbf{x}$-representation:

\begin{equation}
i\hbar \frac{\partial}{\partial t}\langle \mathbf{x}'|\alpha, t_0; t\rangle = \langle \mathbf{x}'|H|\alpha, t_0; t\rangle
\tag{2.179}
\end{equation}

where we have used the fact that the position eigenbras in the Schrödinger picture do not change with time.

\textbf{Kinetic energy contribution:}
Using the momentum representation of the kinetic energy operator, we can write the kinetic-energy contribution to the right-hand side of (2.179) as

\begin{equation}
\left\langle \mathbf{x}'\left|\frac{\mathbf{p}^2}{2m}\right|\alpha, t_0; t\right\rangle = -\left(\frac{\hbar^2}{2m}\right)\nabla'^2 \langle \mathbf{x}'|\alpha, t_0; t\rangle
\tag{2.180}
\end{equation}

\textbf{Potential energy contribution:}
As for $V(\mathbf{x})$, we simply use

\begin{equation}
\langle \mathbf{x}'|V(\mathbf{x}) = \langle \mathbf{x}'|V(\mathbf{x}')
\tag{2.181}
\end{equation}

where $V(\mathbf{x}')$ is no longer an operator.

\textbf{Combining everything:}
\begin{equation}
i\hbar \frac{\partial}{\partial t}\langle \mathbf{x}'|\alpha, t_0; t\rangle = -\left(\frac{\hbar^2}{2m}\right)\nabla'^2 \langle \mathbf{x}'|\alpha, t_0; t\rangle + V(\mathbf{x}')\langle \mathbf{x}'|\alpha, t_0; t\rangle
\tag{2.182}
\end{equation}

which we recognize to be the celebrated time-dependent wave equation of E. Schrödinger, usually written as

\begin{equation}
\boxed{i\hbar \frac{\partial}{\partial t}\psi(\mathbf{x}', t) = -\left(\frac{\hbar^2}{2m}\right)\nabla'^2 \psi(\mathbf{x}', t) + V(\mathbf{x}')\psi(\mathbf{x}', t)}
\tag{2.183}
\end{equation}

The quantum mechanics based on wave equation (2.183) is known as \textbf{wave mechanics}. This equation is, in fact, the starting point of many textbooks on quantum mechanics. In our formalism, however, this is just the Schrödinger equation for a state ket written explicitly in the $x$-basis when the Hamiltonian operator is taken to be (2.177).

\subsection{The Time-Independent Wave Equation}

We now derive the partial differential equation satisfied by energy eigenfunctions. We showed in Section 2.1 that the time dependence of a stationary state is given by $\exp(-iE_{a'}t/\hbar)$. This enables us to write its wave function as

\begin{equation}
\langle \mathbf{x}'|a', t_0; t\rangle = \langle \mathbf{x}'|a'\rangle \exp\left(\frac{-iE_{a'}t}{\hbar}\right)
\tag{2.184}
\end{equation}

where it is understood that initially the system is prepared in a simultaneous eigenstate of $A$ and $H$ with eigenvalues $a'$ and $E_{a'}$, respectively.

\subsubsection{Derivation from Time-Dependent Equation}

Let us now substitute (2.184) into the time-dependent Schrödinger equation (2.182). We are then led to

\begin{equation}
-\left(\frac{\hbar^2}{2m}\right)\nabla'^2 \langle \mathbf{x}'|a'\rangle + V(\mathbf{x}')\langle \mathbf{x}'|a'\rangle = E_{a'}\langle \mathbf{x}'|a'\rangle
\tag{2.185}
\end{equation}

This partial differential equation is satisfied by the energy eigenfunction $\langle \mathbf{x}'|a'\rangle$ with energy eigenvalue $E_{a'}$.

Actually, in wave mechanics where the Hamiltonian operator is given as a function of $\mathbf{x}$ and $\mathbf{p}$, as in (2.177), it is not necessary to refer explicitly to observable $A$ that commutes with $H$ because we can always choose $A$ to be that function of the observables $\mathbf{x}$ and $\mathbf{p}$ which coincides with $H$ itself. We may therefore omit reference to $a'$ and simply write (2.185) as the partial differential equation to be satisfied by the energy eigenfunction $u_E(\mathbf{x}')$:

\begin{equation}
\boxed{-\left(\frac{\hbar^2}{2m}\right)\nabla'^2 u_E(\mathbf{x}') + V(\mathbf{x}')u_E(\mathbf{x}') = E u_E(\mathbf{x}')}
\tag{2.186}
\end{equation}

This is the \textbf{time-independent wave equation} of E. Schrödinger, announced in the first of four monumental papers, all written in the first half of 1926, that laid the foundations of wave mechanics. In the same paper, Schrödinger immediately applied (2.186) to derive the energy spectrum of the hydrogen atom.

\subsubsection{Boundary Conditions and Energy Quantization}

To solve (2.186) some boundary condition has to be imposed. Suppose we seek a solution to (2.186) with

\begin{equation}
E < \lim_{|\mathbf{x}'| \to \infty} V(\mathbf{x}')
\tag{2.187}
\end{equation}

where the inequality relation is to hold for $|\mathbf{x}'| \to \infty$ in any direction. The appropriate boundary condition to be used in this case is

\begin{equation}
\boxed{u_E(\mathbf{x}') \to 0 \quad \text{as} \quad |\mathbf{x}'| \to \infty}
\tag{2.188}
\end{equation}

Physically this means that the particle is bound or confined within a finite region of space. We know from the theory of partial differential equations that (2.186) subject to boundary condition (2.188) allows nontrivial solutions only for a discrete set of values of $E$. It is in this sense that the time-independent Schrödinger equation (2.186) yields the \textit{quantization of energy levels}. Once the partial differential equation (2.186) is written, the problem of finding the energy levels of microscopic physical systems is as straightforward as that of finding the characteristic frequencies of vibrating strings or membranes. In both cases we solve boundary-value problems in mathematical physics.

\subsubsection{Historical Remarks}

A short digression on the history of quantum mechanics is in order here. The fact that exactly soluble eigenvalue problems in the theory of partial differential equations can also be treated using matrix methods was already known to mathematicians in the first quarter of the twentieth century. Furthermore, theoretical physicists like M. Born frequently consulted great mathematicians of the day -- D. Hilbert and H. Weyl, in particular. Yet when matrix mechanics was born in the summer of 1925, it did not immediately occur to the theoretical physicists or to the mathematicians to reformulate it using the language of partial differential equations. Six months after Heisenberg's pioneering paper, wave mechanics was proposed by Schrödinger. However, a close inspection of his papers shows that he was not at all influenced by the earlier works of Heisenberg, Born, and Jordan. Instead, the train of reasoning that led Schrödinger to formulate wave mechanics has its roots in W. R. Hamilton's analogy between optics and mechanics, and the particle-wave hypothesis of L. de Broglie. Once wave mechanics was formulated, many people, including Schrödinger himself, showed the equivalence between wave mechanics and matrix mechanics.

\subsubsection{General Properties of Energy Eigenfunctions}

It is assumed that the reader has some experience in solving the time-dependent and time-independent wave equations, including:
\begin{itemize}[leftmargin=*]
\item Time evolution of a Gaussian wave packet in a force-free region
\item One-dimensional transmission-reflection problems involving a rectangular potential barrier
\item Simple solutions of the time-independent wave equation: particle in a box, particle in a square well, the simple harmonic oscillator, the hydrogen atom
\end{itemize}

\textbf{Key properties of energy eigenfunctions and eigenvalues:}

\textbf{1. Discrete vs. continuous spectrum:} The energy levels exhibit a discrete or continuous spectrum depending on whether or not condition (2.187) is satisfied.

\textbf{2. Sinusoidal vs. damped behavior:} The energy eigenfunction in one dimension is \textit{sinusoidal} or \textit{damped} depending on whether $E - V(\mathbf{x}')$ is positive or negative.

\subsection{Interpretations of the Wave Function}

We now turn to discussions of the physical interpretations of the wave function. In Section 1.7 we commented on the probabilistic interpretation of $|\psi|^2$ that follows from the fact that $\langle \mathbf{x}'|\alpha, t_0; t\rangle$ is to be regarded as an expansion coefficient of $|\alpha, t_0; t\rangle$ in terms of the position eigenkets $\{|\mathbf{x}'\rangle\}$.

\subsubsection{Probability Density}

The quantity $\rho(\mathbf{x}', t)$ defined by

\begin{equation}
\boxed{\rho(\mathbf{x}', t) = |\psi(\mathbf{x}', t)|^2 = |\langle \mathbf{x}'|\alpha, t_0; t\rangle|^2}
\tag{2.189}
\end{equation}

is therefore regarded as the \textbf{probability density} in wave mechanics. Specifically, when we use a detector that ascertains the presence of the particle within a small volume element $d^3x'$ around $\mathbf{x}'$, the probability of recording a positive result at time $t$ is given by $\rho(\mathbf{x}', t)d^3x'$.

In the remainder of this section we use $\mathbf{x}$ for $\mathbf{x}'$ because the position operator will not appear.

\subsubsection{Continuity Equation and Probability Flux}

Using Schrödinger's time-dependent wave equation, it is straightforward to derive the continuity equation

\begin{equation}
\boxed{\frac{\partial \rho}{\partial t} + \nabla \cdot \mathbf{j} = 0}
\tag{2.190}
\end{equation}

where $\rho(\mathbf{x}, t)$ stands for $|\psi|^2$ as before, and $\mathbf{j}(\mathbf{x}, t)$, known as the \textbf{probability flux}, is given by

\begin{align}
\mathbf{j}(\mathbf{x}, t) &= -\left(\frac{i\hbar}{2m}\right)[\psi^*\nabla\psi - (\nabla\psi^*)\psi] \notag \\
&= \left(\frac{\hbar}{m}\right)\text{Im}(\psi^*\nabla\psi)
\tag{2.191}
\end{align}

The reality of the potential $V$ (or the Hermiticity of the $V$ operator) has played a crucial role in our obtaining this result. Conversely, a complex potential can phenomenologically account for the disappearance of a particle; such a potential is often used for nuclear reactions where incident particles get absorbed by nuclei.

\textbf{Relation to momentum:}
We may intuitively expect that the probability flux $\mathbf{j}$ is related to momentum. This is indeed the case for $\mathbf{j}$ \textit{integrated over all space}. From (2.191) we obtain

\begin{equation}
\int d^3x\, \mathbf{j}(\mathbf{x}, t) = \frac{\langle \mathbf{p}\rangle_t}{m}
\tag{2.192}
\end{equation}

where $\langle \mathbf{p}\rangle_t$ is the expectation value of the momentum operator at time $t$.

Equation (2.190) is reminiscent of the continuity equation in fluid dynamics that characterizes a hydrodynamic flow of a fluid in a source-free, sink-free region. Indeed, historically Schrödinger was first led to interpret $|\psi|^2$ as the actual matter density, or $e|\psi|^2$ as the actual electric charge density. If we adopt such a view, we are led to face some bizarre consequences.

\subsubsection{Statistical Interpretation}

A typical argument for a position measurement might go as follows. An atomic electron is to be regarded as a continuous distribution of matter filling up a finite region of space around the nucleus; yet, when a measurement is made to make sure that the electron is at some particular point, this continuous distribution of matter suddenly shrinks to a pointlike particle with no spatial extension. The more satisfactory \textit{statistical interpretation} of $|\psi|^2$ as the probability density was first given by M. Born.

\subsubsection{Physical Significance of the Phase}

To understand the physical significance of the wave function, let us write it as

\begin{equation}
\psi(\mathbf{x}, t) = \sqrt{\rho(\mathbf{x}, t)}\exp\left[\frac{iS(\mathbf{x}, t)}{\hbar}\right]
\tag{2.193}
\end{equation}

with $S$ real and $\rho > 0$, which can always be done for any complex function of $\mathbf{x}$ and $t$. The meaning of $\rho$ has already been given. What is the physical interpretation of $S$? Noting

\begin{equation}
\psi^*\nabla\psi = \sqrt{\rho}\,\nabla(\sqrt{\rho}) + \left(\frac{i}{\hbar}\right)\rho\nabla S
\tag{2.194}
\end{equation}

we can write the probability flux as [see (2.191)]

\begin{equation}
\boxed{\mathbf{j} = \frac{\rho\nabla S}{m}}
\tag{2.195}
\end{equation}

We now see that there is more to the wave function than the fact that $|\psi|^2$ is the probability density; the gradient of the phase $S$ contains a vital piece of information. From (2.195) we see that the \textit{spatial variation of the phase} of the wave function characterizes the probability flux; the stronger the phase variation, the more intense the flux. The direction of $\mathbf{j}$ at some point $\mathbf{x}$ is seen to be normal to the surface of constant phase that goes through that point.

\textbf{Example: Plane wave (momentum eigenfunction)}
\begin{equation}
\psi(\mathbf{x}, t) \propto \exp\left(\frac{i\mathbf{p}\cdot\mathbf{x}}{\hbar} - \frac{iEt}{\hbar}\right)
\tag{2.196}
\end{equation}

where $\mathbf{p}$ stands for the eigenvalue of the momentum operator. All this is evident because

\begin{equation}
\nabla S = \mathbf{p}
\tag{2.197}
\end{equation}

\subsubsection{Hydrodynamic Analogy and Its Limitations}

More generally, it is tempting to regard $\nabla S/m$ as some kind of ``velocity,''

\begin{equation}
\text{``}\mathbf{v}\text{''} = \frac{\nabla S}{m}
\tag{2.198}
\end{equation}

and to write the continuity equation (2.190) as

\begin{equation}
\frac{\partial \rho}{\partial t} + \nabla \cdot (\rho\text{``}\mathbf{v}\text{''}) = 0
\tag{2.199}
\end{equation}

just as in fluid dynamics. However, we would like to caution the reader against a too literal interpretation of $\mathbf{j}$ as $\rho$ times the velocity defined at every point in space, because a simultaneous precision measurement of position and velocity would necessarily violate the uncertainty principle.

\newpage
\section{Elementary Solutions to Schrödinger's Wave Equation}

It is both instructive and useful to look at some relatively elementary solutions to (2.186) for particular choices of the potential energy function $V(\mathbf{x})$. In this section we choose some particular examples to illustrate contemporary physics and/or which will be useful in later chapters of this textbook.

\subsection{Free Particle in Three Dimensions}

The case $V(\mathbf{x}) = 0$ has fundamental significance. We will consider the solution to Schrödinger's equation here in three dimensions using Cartesian coordinates. The solution in spherical coordinates will be left until our treatment of angular momentum is presented in the next chapter. Equation (2.186) becomes

\begin{equation}
\nabla^2 u_E(\mathbf{x}) = -\frac{2mE}{\hbar^2}u_E(\mathbf{x})
\tag{2.205}
\end{equation}

Define a vector $\mathbf{k}$ where

\begin{equation}
\mathbf{k}^2 = k_x^2 + k_y^2 + k_z^2 \equiv \frac{2mE}{\hbar^2} = \frac{\mathbf{p}^2}{\hbar^2}
\tag{2.206}
\end{equation}

that is, $\mathbf{p} = \hbar\mathbf{k}$.

\subsubsection{Separation of Variables}

Differential equation (2.205) is easily solved using the technique known as ``separation of variables.'' Writing

\begin{equation}
u_E(\mathbf{x}) = u_x(x)u_y(y)u_z(z)
\tag{2.207}
\end{equation}

we arrive at

\begin{equation}
\left[\frac{1}{u_x}\frac{d^2u_x}{dx^2} + k_x^2\right] + \left[\frac{1}{u_y}\frac{d^2u_y}{dy^2} + k_y^2\right] + \left[\frac{1}{u_z}\frac{d^2u_z}{dz^2} + k_z^2\right] = 0
\tag{2.208}
\end{equation}

This leads to individual plane wave solutions $u_w(w) = c_w e^{ik_w w}$ for $w = x, y, z$. Note that one gets the same energy $E$ for values $\pm k_w$.

Collecting these solutions, and combining the normalization constants, we obtain

\begin{equation}
u_E(\mathbf{x}) = c_x c_y c_z e^{ik_x x + ik_y y + ik_z z} = Ce^{i\mathbf{k}\cdot\mathbf{x}}
\tag{2.209}
\end{equation}

\subsubsection{Box Normalization}

The normalization constant $C$ presents the usual difficulties, which are generally handled by using a $\delta$-function normalization condition. It is convenient in many cases, however, to use a ``big box'' normalization, where all space is contained within a cube of side length $L$. We impose periodic boundary conditions on the box, and thereby obtain a finite normalization constant $C$. For any real calculation, we simply let the size $L \to \infty$ at the end of the calculation.

Imposing the condition $u_x(x + L) = u_x(x)$ we have $k_x L = 2\pi n_x$ where $n_x$ is an integer. That is

\begin{equation}
k_x = \frac{2\pi}{L}n_x \qquad k_y = \frac{2\pi}{L}n_y \qquad k_z = \frac{2\pi}{L}n_z
\tag{2.210}
\end{equation}

and the normalization criterion becomes

\begin{equation}
1 = \int_0^L dx \int_0^L dy \int_0^L dz\, u_E^*(\mathbf{x})u_E(\mathbf{x}) = L^3|C|^2
\tag{2.211}
\end{equation}

in which case $C = 1/L^{3/2}$ and

\begin{equation}
\boxed{u_E(\mathbf{x}) = \frac{1}{L^{3/2}}e^{i\mathbf{k}\cdot\mathbf{x}}}
\tag{2.212}
\end{equation}

\textbf{Energy eigenvalue:}
\begin{equation}
\boxed{E = \frac{\mathbf{p}^2}{2m} = \frac{\hbar^2\mathbf{k}^2}{2m} = \frac{\hbar^2}{2m}\left(\frac{2\pi}{L}\right)^2(n_x^2 + n_y^2 + n_z^2)}
\tag{2.213}
\end{equation}

\subsubsection{Density of States}

The eightfold degeneracy we mentioned earlier corresponds to the eight combinations of $(\pm n_x, \pm n_y, \pm n_z)$, but the degeneracy can actually be much larger since, in some cases, there are various combinations of $n_x$, $n_y$, and $n_z$ which can give the same $E$. In fact, in the (realistic) limit where $L$ is very large, there can be a large number of states $N$ which have an energy between $E$ and $E + dE$. This ``density of states'' $dN/dE$ is an important quantity for calculations of processes which include free particles.

To calculate the density of states, imagine a spherical shell in $\mathbf{k}$ space with radius $|\mathbf{k}| = 2\pi|\mathbf{n}|/L$ and thickness $d|\mathbf{k}| = 2\pi d|\mathbf{n}|/L$. All states within this shell have energy $E = \hbar^2\mathbf{k}^2/2m$. The number of states $dN$ within this shell is $4\pi n^2 d|\mathbf{n}|$. Therefore

\begin{align}
\frac{dN}{dE} &= \frac{4\pi n^2 d|\mathbf{n}|}{\hbar^2|\mathbf{k}|d|\mathbf{k}|/m} = \frac{4\pi}{\hbar^2}m\left(\frac{L}{2\pi}\right)^2|\mathbf{k}|\frac{L}{2\pi} \notag \\
&= \boxed{\frac{m^{3/2}E^{1/2}L^3}{\sqrt{2}\pi^2\hbar^3}}
\tag{2.214}
\end{align}

In a typical ``real'' calculation, the density of states will be multiplied by some probability that involves $u_E^*(\mathbf{x})u_E(\mathbf{x})$. In this case, the factors of $L^3$ will cancel explicitly, so the limit $L \to \infty$ is trivial.

\subsubsection{Probability Flux for Free Particle}

This ``big box'' normalization also yields the correct answer for the probability flux. Rewriting (2.196) with this normalization, we have

\begin{equation}
\psi(\mathbf{x}, t) = \frac{1}{L^{3/2}}\exp\left(\frac{i\mathbf{p}\cdot\mathbf{x}}{\hbar} - \frac{iEt}{\hbar}\right)
\tag{2.215}
\end{equation}

in which case we find

\begin{equation}
\mathbf{j}(\mathbf{x}, t) = \frac{\hbar}{m}\text{Im}(\psi^*\nabla\psi) = \frac{\hbar\mathbf{k}}{m}\frac{1}{L^3} = \mathbf{v}\rho
\tag{2.216}
\end{equation}

where $\rho = 1/L^3$ is indeed the probability density.

\subsection{The Simple Harmonic Oscillator}

We saw an elegant solution for the case $V(x) = m\omega^2 x^2/2$ in Section 2.3, which yielded the energy eigenvalues, eigenstates, and wave functions. Here, we demonstrate a different approach which solves the differential equation

\begin{equation}
-\frac{\hbar^2}{2m}\frac{d^2}{dx^2}u_E(x) + \frac{1}{2}m\omega^2 x^2 u_E(x) = Eu_E(x)
\tag{2.217}
\end{equation}

Our approach will introduce the concept of \textit{generating functions}, a generally useful technique which arises in many treatments of differential eigenvalue problems.

\subsubsection{Dimensionless Form}

First, transform (2.217) using the dimensionless position $y \equiv x/x_0$ where $x_0 \equiv \sqrt{\hbar/m\omega}$. Also introduce a dimensionless energy variable $\varepsilon \equiv 2E/\hbar\omega$. The differential equation we need to solve becomes therefore

\begin{equation}
\frac{d^2}{dy^2}u(y) + (\varepsilon - y^2)u(y) = 0
\tag{2.218}
\end{equation}

For $y \to \pm\infty$, the solution must tend to zero, otherwise the wave function will not be normalizable and hence unphysical. The differential equation $w''(y) - y^2 w(y) = 0$ has solutions $w(y) \propto \exp(\pm y^2/2)$, so we would have to choose the minus sign. We then ``remove'' the asymptotic behavior of the wave function by writing

\begin{equation}
u(y) = h(y)e^{-y^2/2}
\tag{2.219}
\end{equation}

where the function $h(y)$ satisfies the differential equation

\begin{equation}
\frac{d^2 h}{dy^2} - 2y\frac{dh}{dy} + (\varepsilon - 1)h(y) = 0
\tag{2.220}
\end{equation}

To this point, we have followed the traditional solution of the simple harmonic oscillator as found in many textbooks. Typically, one would now look for a series solution for $h(y)$ and discover that a normalizable solution is only possible if the series terminates. (In fact, we use this approach for the three-dimensional isotropic harmonic oscillator in this book. See Section 3.7.) One forces this termination by imposing the condition that $\varepsilon - 1$ be an even, nonnegative integer $2n$, $n = 0, 1, 2, \ldots$. The solutions are then written using the resulting polynomials $h_n(y)$. Of course, $\varepsilon - 1 = 2n$ is equivalent to $E = (n + \frac{1}{2})\hbar\omega$, the quantization relation (2.143).

\subsubsection{Generating Function Approach}

Let us take a different approach. Consider the ``Hermite polynomials'' $H_n(x)$ defined by the ``generating function'' $g(x, t)$ through

\begin{align}
g(x, t) &\equiv e^{-t^2 + 2tx} \tag{2.221a} \\
&\equiv \sum_{n=0}^{\infty} H_n(x)\frac{t^n}{n!} \tag{2.221b}
\end{align}

Some properties of the $H_n(x)$ are immediately obvious. For example, $H_0(x) = 1$. Also, since

\begin{equation}
g(0, t) = e^{-t^2} = \sum_{n=0}^{\infty}\frac{(-1)^n}{n!}t^{2n}
\tag{2.222}
\end{equation}

it is clear that $H_n(0) = 0$ if $n$ is odd, since this series only involves even powers of $t$. On the other hand, if we restrict ourselves to even values of $n$, we have

\begin{equation}
g(0, t) = e^{-t^2} = \sum_{n=0}^{\infty}\frac{(-1)^{(n/2)}}{(n/2)!}t^n = \sum_{n=0}^{\infty}\frac{(-1)^{(n/2)}}{(n/2)!}\frac{n!}{n!}t^n
\tag{2.223}
\end{equation}

and so $H_n(0) = (-1)^{n/2}n!/(n/2)!$. Also, since $g(-x, t)$ reverses the sign only on terms with odd powers of $t$, $H_n(-x) = (-1)^n H_n(x)$.

\subsubsection{Recursion Relations from Derivatives}

We can take derivatives of $g(x, t)$ to build the Hermite polynomials using recursion relations between them and their derivatives. The trick is that we can differentiate the analytic form of the generating function (2.221a) or the series form (2.221b) and then compare results.

\textbf{Derivative with respect to $x$:}
If we take the derivative using (2.221a), then

\begin{equation}
\frac{\partial g}{\partial x} = 2tg(x, t) = \sum_{n=0}^{\infty}2H_n(x)\frac{t^{n+1}}{n!} = \sum_{n=0}^{\infty}2(n+1)H_n(x)\frac{t^{n+1}}{(n+1)!}
\tag{2.224}
\end{equation}

where we insert the series definition of the generating function after taking the derivative. On the other hand, we can take the derivative of (2.221b) directly, in which case

\begin{equation}
\frac{\partial g}{\partial x} = \sum_{n=0}^{\infty}H_n'(x)\frac{t^n}{n!}
\tag{2.225}
\end{equation}

Comparing (2.224) and (2.225) shows that

\begin{equation}
\boxed{H_n'(x) = 2nH_{n-1}(x)}
\tag{2.226}
\end{equation}

This is enough information for us to build the Hermite polynomials:
\begin{align*}
H_0(x) &= 1 \\
\text{so } H_1'(x) &= 2, \text{ therefore } H_1(x) = 2x \\
\text{so } H_2'(x) &= 8x, \text{ therefore } H_2(x) = 4x^2 - 2 \\
\text{so } H_3'(x) &= 24x^2 - 12, \text{ therefore } H_3(x) = 8x^3 - 12x \\
&\vdots
\end{align*}

\textbf{Derivative with respect to $t$:}
To see why this is relevant to the simple harmonic oscillator, consider the derivative of the generating function with respect to $t$. If we start with (2.221a) then

\begin{align}
\frac{\partial g}{\partial t} &= -2tg(x, t) + 2xg(x, t) \notag \\
&= -\sum_{n=0}^{\infty}2H_n(x)\frac{t^{n+1}}{n!} + \sum_{n=0}^{\infty}2xH_n(x)\frac{t^n}{n!} \notag \\
&= -\sum_{n=0}^{\infty}2nH_{n-1}(x)\frac{t^n}{n!} + \sum_{n=0}^{\infty}2xH_n(x)\frac{t^n}{n!}
\tag{2.227}
\end{align}

Or, if we differentiate (2.221b) then we have

\begin{equation}
\frac{\partial g}{\partial t} = \sum_{n=0}^{\infty}nH_n(x)\frac{t^{n-1}}{n!} = \sum_{n=0}^{\infty}H_{n+1}(x)\frac{t^n}{n!}
\tag{2.228}
\end{equation}

Comparing (2.227) and (2.228) gives us the recursion relation

\begin{equation}
\boxed{H_{n+1}(x) = 2xH_n(x) - 2nH_{n-1}(x)}
\tag{2.229}
\end{equation}

\subsubsection{Hermite Differential Equation}

We combine (2.229) with (2.226) to find

\begin{align}
H_n''(x) &= 2n \cdot 2(n-1)H_{n-2}(x) \notag \\
&= 2n[2xH_{n-1}(x) - H_n(x)] \notag \\
&= 2xH_n'(x) - 2nH_n(x)
\tag{2.230}
\end{align}

In other words, the Hermite polynomials satisfy the differential equation

\begin{equation}
\boxed{H_n''(x) - 2xH_n'(x) + 2nH_n(x) = 0}
\tag{2.231}
\end{equation}

where $n$ is a nonnegative integer. This, however, is the same as the Schrödinger equation written as (2.220) since $\varepsilon - 1 = 2n$.

\subsubsection{Wave Functions and Orthogonality}

That is, the wave functions for the simple harmonic oscillator are given by

\begin{equation}
\boxed{u_n(x) = c_n H_n\left(x\sqrt{\frac{m\omega}{\hbar}}\right)e^{-m\omega x^2/2\hbar}}
\tag{2.232}
\end{equation}

up to some normalization constant $c_n$. This constant can be determined from the orthogonality relationship

\begin{equation}
\int_{-\infty}^{\infty}H_n(x)H_m(x)e^{-x^2}\,dx = \pi^{1/2}2^n n!\,\delta_{nm}
\tag{2.233}
\end{equation}

which is easily proved using the generating function. See Problem 2.25 at the end of this chapter.

Generating functions have a usefulness that far outreaches our limited application here. Among other things, many of the orthogonal polynomials which arise from solving the Schrödinger equation for different potentials can be derived from generating functions. See, for example, Problem 3.30 in Chapter 3.

\subsection{The Linear Potential}

Perhaps the first potential energy function, with bound states, to come to mind is the linear potential, namely

\begin{equation}
V(x) = k|x|
\tag{2.234}
\end{equation}

where $k$ is an arbitrary positive constant. Given a total energy $E$, this potential has a classical turning point at a value $x = a$ where $E = ka$. This point will be important for understanding the quantum behavior of a particle of mass $m$ bound by this potential.

The Schrödinger equation becomes

\begin{equation}
-\frac{\hbar^2}{2m}\frac{d^2 u_E}{dx^2} + k|x|u_E(x) = Eu_E(x)
\tag{2.235}
\end{equation}

\subsubsection{Symmetry and Boundary Conditions}

It is easiest to deal with the absolute value by restricting our attention to $x \geq 0$. We can do this because $V(-x) = V(x)$, so there are two types of solutions, namely $u_E(-x) = \pm u_E(x)$. In either case, we need $u_E(x)$ to tend towards zero as $x \to \infty$.

\textbf{Odd parity:} If $u_E(-x) = -u_E(x)$, then we need $u_E(0) = 0$.

\textbf{Even parity:} If $u_E(-x) = +u_E(x)$, then we have $u_E'(0) = 0$, since $u_E'(\varepsilon) - u_E'(-\varepsilon) \equiv 0$, even for $\varepsilon \to 0$.

\subsubsection{Dimensionless Form and the Airy Equation}

Once again, we write the differential equation in terms of dimensionless variables, based on appropriate scales for length and energy. In this case, the dimensionless length scale is $x_0 = (\hbar^2/mk)^{1/3}$ and the dimensionless energy scale is $E_0 = kx_0 = (\hbar^2 k^2/m)^{1/3}$. Defining $y \equiv x/x_0$ and $\varepsilon \equiv E/E_0$ allows us to rewrite (2.235) as

\begin{equation}
\frac{d^2 u_E}{dy^2} - 2(y - \varepsilon)u_E(y) = 0, \qquad y \geq 0
\tag{2.236}
\end{equation}

Notice that $y = \varepsilon$ when $x = E/k$, i.e. the classical turning point $x = a$. In fact, defining a translated position variable $z \equiv 2^{1/3}(y - \varepsilon)$, (2.236) becomes

\begin{equation}
\boxed{\frac{d^2 u_E}{dz^2} - zu_E(z) = 0}
\tag{2.237}
\end{equation}

This is the \textbf{Airy equation}, and the solution is the Airy function $\text{Ai}(z)$.

\subsubsection{Properties of the Airy Function}

The Airy function has a peculiar behavior: oscillatory for negative values of the argument, and decreasing rapidly towards zero for positive values. This is exactly the behavior we expect for the wave function, since $z = 0$ is the classical turning point.

\subsubsection{Energy Quantization}

The boundary conditions at $x = 0$ translate into zeros for either $\text{Ai}'(z)$ or $\text{Ai}(z)$ where $z = -2^{1/3}\varepsilon$. In other words, the zeros of the Airy function or its derivative determine the quantized energies:

\textbf{Even parity} (from $u_E'(0) = 0$):
\begin{equation}
\text{Ai}'(z) = 0 \qquad \text{for } z = -1.019, -3.249, -4.820, \ldots
\tag{2.238}
\end{equation}

\textbf{Odd parity} (from $u_E(0) = 0$):
\begin{equation}
\text{Ai}(z) = 0 \qquad \text{for } z = -2.338, -4.088, -5.521, \ldots
\tag{2.239}
\end{equation}

\textbf{Ground state energy:}
\begin{equation}
\boxed{E = \frac{1.019}{2^{1/3}}\left(\frac{\hbar^2 k^2}{m}\right)^{1/3}}
\end{equation}

\subsection{The WKB (Semiclassical) Approximation}

The WKB solution, after G. Wentzel, A. Kramers, and L. Brillouin, is an approximation technique based on making use of regions where the wavelength is much shorter than the typical distance over which the potential energy varies. Such is \textit{never} the case near classical turning points, but this is where the linear potential solution can be used to join the solutions on either side of them.

\subsubsection{Setup}

Again restricting ourselves to one dimension, we write Schrödinger's wave equation as

\begin{equation}
\frac{d^2 u_E}{dx^2} + \frac{2m}{\hbar^2}(E - V(x))u_E(x) = 0
\tag{2.241}
\end{equation}

Define the quantities

\begin{align}
k(x) &\equiv \left[\frac{2m}{\hbar^2}(E - V(x))\right]^{1/2} \qquad \text{for } E > V(x) \tag{2.242a} \\
k(x) &\equiv -i\kappa(x) \equiv -i\left[\frac{2m}{\hbar^2}(V(x) - E)\right]^{1/2} \qquad \text{for } E < V(x) \tag{2.242b}
\end{align}

and so (2.241) becomes

\begin{equation}
\frac{d^2 u_E}{dx^2} + [k(x)]^2 u_E(x) = 0
\tag{2.243}
\end{equation}

Now, if $V(x)$ were not changing with $x$, then $k(x)$ would be a constant, and $u(x) \propto \exp(\pm ikx)$ would solve (2.243). Consequently, if we assume that $V(x)$ varies only ``slowly'' with $x$, then we are tempted to try a solution of the form

\begin{equation}
u_E(x) \equiv \exp[iW(x)/\hbar]
\tag{2.244}
\end{equation}

In this case, (2.243) becomes

\begin{equation}
i\hbar\frac{d^2 W}{dx^2} - \left(\frac{dW}{dx}\right)^2 + \hbar^2[k(x)]^2 = 0
\tag{2.245}
\end{equation}

\subsubsection{The WKB Condition}

We consider a solution to this equation under the condition that

\begin{equation}
\boxed{\hbar\left|\frac{d^2 W}{dx^2}\right| \ll \left|\frac{dW}{dx}\right|^2}
\tag{2.246}
\end{equation}

This quantifies our notion of a ``slowly varying'' potential $V(x)$.

\subsubsection{Derivation of the WKB Wave Function}

Using condition (2.246), the lowest-order approximation for $W(x)$ is

\begin{equation}
W_0'(x) = \pm\hbar k(x)
\tag{2.247}
\end{equation}

leading to a first-order approximation for $W(x)$, based on

\begin{equation}
\left(\frac{dW_1}{dx}\right)^2 = \hbar^2[k(x)]^2 + i\hbar W_0''(x) = \hbar^2[k(x)]^2 \pm i\hbar^2 k'(x)
\tag{2.248}
\end{equation}

where the second term in (2.248) is much smaller than the first, so that

\begin{align}
W(x) &\approx W_1(x) = \pm\hbar\int^x dx'\,[k^2(x') \pm ik'(x')]^{1/2} \notag \\
&\approx \pm\hbar\int^x dx'\,k(x')\left[1 \pm \frac{i}{2}\frac{k'(x')}{k^2(x')}\right] \notag \\
&= \pm\hbar\int^x dx'\,k(x') + \frac{i}{2}\hbar\ln[k(x)]
\tag{2.249}
\end{align}

\textbf{The WKB wave function:}
\begin{equation}
\boxed{u_E(x) \approx \exp[iW(x)/\hbar] = \frac{1}{[k(x)]^{1/2}}\exp\left[\pm i\int^x dx'\,k(x')\right]}
\tag{2.250}
\end{equation}

Note that this specifies a choice of two solutions ($\pm$) either in the region where $E > V(x)$, with $k(x)$ from (2.242a), or in the region where $E < V(x)$, with $k(x)$ from (2.242b).

\subsubsection{Connection Formulas}

For a potential well with two turning points $x_1$ and $x_2$, the wave function must behave like (2.250), with $k(x)$ given by (2.242a) in region II (classically allowed), and by (2.242b) in regions I and III (classically forbidden). The solutions in the neighborhood of the turning points are given by Airy functions.

The asymptotic dependences of the Airy function are

\begin{align}
\text{Ai}(z) &\to \frac{1}{2\sqrt{\pi}}z^{-1/4}\exp\left(-\frac{2}{3}z^{3/2}\right) \qquad z \to +\infty \tag{2.251a} \\
\text{Ai}(z) &\to \frac{1}{\sqrt{\pi}}|z|^{-1/4}\cos\left(\frac{2}{3}|z|^{3/2} - \frac{\pi}{4}\right) \qquad z \to -\infty \tag{2.251b}
\end{align}

\textbf{Connection formula from region I to region II:}
\begin{equation}
\left\{\frac{1}{[V(x) - E]^{1/4}}\right\}\exp\left[-\frac{1}{\hbar}\int_x^{x_1}dx'\sqrt{2m[V(x') - E]}\right] \to \left\{\frac{2}{[E - V(x)]^{1/4}}\right\}\cos\left[\frac{1}{\hbar}\int_{x_1}^x dx'\sqrt{2m[E - V(x')]} - \frac{\pi}{4}\right]
\tag{2.252}
\end{equation}

\textbf{Connection formula from region III to region II:}
\begin{equation}
\left\{\frac{1}{[V(x) - E]^{1/4}}\right\}\exp\left[-\frac{1}{\hbar}\int_{x_2}^x dx'\sqrt{2m[V(x') - E]}\right] \to \left\{\frac{2}{[E - V(x)]^{1/4}}\right\}\cos\left[-\frac{1}{\hbar}\int_x^{x_2} dx'\sqrt{2m[E - V(x')]} + \frac{\pi}{4}\right]
\tag{2.253}
\end{equation}

\subsubsection{WKB Quantization Condition}

Of course, we must obtain the same form for the wave function in region II, regardless of which turning point is analyzed. This implies that the arguments of the cosine in (2.252) and (2.253) must differ at most by an integer multiple of $\pi$. In this way we obtain the \textbf{WKB quantization condition}:

\begin{equation}
\boxed{\int_{x_1}^{x_2} dx\sqrt{2m[E - V(x)]} = \left(n + \frac{1}{2}\right)\pi\hbar \quad (n = 0, 1, 2, 3, \ldots)}
\tag{2.254}
\end{equation}

Apart from the difference between $n + \frac{1}{2}$ and $n$, this equation is simply the quantization condition of the old quantum theory due to A. Sommerfeld and W. Wilson, originally written in 1915 as

\begin{equation}
\oint p\,dq = nh
\tag{2.255}
\end{equation}

where $h$ is Planck's $h$, not Dirac's $\hbar$, and the integral is evaluated over one whole period of classical motion, from $x_1$ to $x_2$ and back.

\subsubsection{Example: Bouncing Ball}

Consider the energy spectrum of a ball bouncing up and down over a hard surface:

\begin{equation}
V = \begin{cases} mgx & \text{for } x > 0 \\ \infty & \text{for } x < 0 \end{cases}
\tag{2.256}
\end{equation}

A better approach is to consider the odd-parity solutions of a modified problem:

\begin{equation}
V(x) = mg|x| \quad (-\infty < x < \infty)
\tag{2.258}
\end{equation}

whose turning points are $x_1 = -E/(mg)$ and $x_2 = E/(mg)$.

The quantization condition then becomes

\begin{equation}
\int_0^{E/mg} dx\sqrt{2m(E - mgx)} = \left(n - \frac{1}{4}\right)\pi\hbar \quad (n = 1, 2, 3, 4, \ldots)
\tag{2.261}
\end{equation}

This integral is elementary, and we obtain

\begin{equation}
\boxed{E_n = \left\{\frac{\left[3\left(n - \frac{1}{4}\right)\pi\right]^{2/3}}{2}\right\}(mg^2\hbar^2)^{1/3}}
\tag{2.262}
\end{equation}

for the quantized energy levels of the bouncing ball.

\subsubsection{Physical Interpretation: Semiclassical Limit}

The time-dependent wave function becomes

\begin{equation}
\psi(x, t) \propto u_E(x)\exp(-iEt/\hbar) = \exp(iW(x)/\hbar - iEt/\hbar)
\tag{2.263}
\end{equation}

Comparing this to (2.193) we see that $W(x)$ corresponds directly to Hamilton's characteristic function. Indeed, condition (2.246) is the condition for reaching the classical limit. For these reasons, the WKB approximation is frequently referred to as a ``semiclassical'' approximation.

\textbf{Validity condition:}
Condition (2.246) is equivalent to $|k'(x)| \ll |k^2(x)|$. In terms of the de Broglie wavelength divided by $2\pi$:

\begin{equation}
\boxed{\lambdabar = \frac{\hbar}{\sqrt{2m[E - V(x)]}} \ll \frac{2[E - V(x)]}{|dV/dx|}}
\tag{2.264}
\end{equation}

In other words, $\lambdabar$ must be small compared with the characteristic distance over which the potential varies appreciably. Roughly speaking, the potential must be essentially constant over many wavelengths. Thus we see that the semiclassical picture is reliable \textit{in the short-wavelength limit}.

\newpage
\section{Propagators and Feynman Path Integrals}

\subsection{Propagators in Wave Mechanics}

In Section 2.1 we showed how the most general time-evolution problem with a time-independent Hamiltonian can be solved once we expand the initial ket in terms of the eigenkets of an observable that commutes with $H$. Let us translate this statement into the language of wave mechanics.

We start with

\begin{equation}
|\alpha, t_0; t\rangle = \exp\left[\frac{-iH(t - t_0)}{\hbar}\right]|\alpha, t_0\rangle = \sum_{a'} |a'\rangle\langle a'|\alpha, t_0\rangle \exp\left[\frac{-iE_{a'}(t - t_0)}{\hbar}\right]
\tag{2.265}
\end{equation}

Multiplying both sides by $\langle \mathbf{x}'|$ on the left, we have

\begin{equation}
\langle \mathbf{x}'|\alpha, t_0; t\rangle = \sum_{a'} \langle \mathbf{x}'|a'\rangle\langle a'|\alpha, t_0\rangle \exp\left[\frac{-iE_{a'}(t - t_0)}{\hbar}\right]
\tag{2.266}
\end{equation}

which is of the form

\begin{equation}
\psi(\mathbf{x}', t) = \sum_{a'} c_{a'}(t_0) u_{a'}(\mathbf{x}') \exp\left[\frac{-iE_{a'}(t - t_0)}{\hbar}\right]
\tag{2.267}
\end{equation}

with

\begin{equation}
u_{a'}(\mathbf{x}') = \langle \mathbf{x}'|a'\rangle
\tag{2.268}
\end{equation}

standing for the eigenfunction of operator $A$ with eigenvalue $a'$. Note also that

\begin{equation}
\langle a'|\alpha, t_0\rangle = \int d^3x'\,\langle a'|\mathbf{x}'\rangle\langle \mathbf{x}'|\alpha, t_0\rangle
\tag{2.269}
\end{equation}

which we recognize as the usual rule in wave mechanics for getting the expansion coefficients of the initial state:

\begin{equation}
c_{a'}(t_0) = \int d^3x'\,u_{a'}^*(\mathbf{x}')\psi(\mathbf{x}', t_0)
\tag{2.270}
\end{equation}

\subsubsection{Definition of the Propagator}

All this should be straightforward and familiar. Now (2.266) together with (2.269) can also be visualized as some kind of integral operator acting on the initial wave function to yield the final wave function:

\begin{equation}
\boxed{\psi(\mathbf{x}'', t) = \int d^3x'\,K(\mathbf{x}'', t; \mathbf{x}', t_0)\psi(\mathbf{x}', t_0)}
\tag{2.271}
\end{equation}

Here the kernel of the integral operator, known as the \textbf{propagator} in wave mechanics, is given by

\begin{equation}
\boxed{K(\mathbf{x}'', t; \mathbf{x}', t_0) = \sum_{a'} \langle \mathbf{x}''|a'\rangle\langle a'|\mathbf{x}'\rangle \exp\left[\frac{-iE_{a'}(t - t_0)}{\hbar}\right]}
\tag{2.272}
\end{equation}

In any given problem the propagator depends only on the potential and is independent of the initial wave function. It can be constructed once the energy eigenfunctions and their eigenvalues are given.

\subsubsection{Physical Interpretation}

Clearly, the time evolution of the wave function is completely predicted if $K(\mathbf{x}'', t; \mathbf{x}', t_0)$ is known and $\psi(\mathbf{x}', t_0)$ is given initially. In this sense Schrödinger's wave mechanics is a \textit{perfectly causal theory}. The time development of a wave function subjected to some potential is as ``deterministic'' as anything else in classical mechanics \textit{provided that the system is left undisturbed}. The only peculiar feature, if any, is that when a measurement intervenes, the wave function changes abruptly, in an uncontrollable way, into one of the eigenfunctions of the observable being measured.

\subsubsection{Properties of the Propagator}

\textbf{Property 1:} For $t > t_0$, $K(\mathbf{x}'', t; \mathbf{x}', t_0)$ satisfies Schrödinger's time-dependent wave equation in the variables $\mathbf{x}''$ and $t$, with $\mathbf{x}'$ and $t_0$ fixed. This is evident from (2.272) because $\langle \mathbf{x}''|a'\rangle\exp[-iE_{a'}(t - t_0)/\hbar]$, being the wave function corresponding to $\mathcal{U}(t, t_0)|a'\rangle$, satisfies the wave equation.

\textbf{Property 2:}
\begin{equation}
\lim_{t \to t_0} K(\mathbf{x}'', t; \mathbf{x}', t_0) = \delta^3(\mathbf{x}'' - \mathbf{x}')
\tag{2.273}
\end{equation}

which is also obvious; as $t \to t_0$, because of the completeness of $\{|a'\rangle\}$, sum (2.272) just reduces to $\langle \mathbf{x}''|\mathbf{x}'\rangle$.

Because of these two properties, the propagator (2.272), regarded as a function of $\mathbf{x}''$, is simply the wave function at $t$ of a particle which was localized \textit{precisely} at $\mathbf{x}'$ at some earlier time $t_0$ ($< t$). Indeed, this interpretation follows, perhaps more elegantly, from noting that (2.272) can also be written as

\begin{equation}
\boxed{K(\mathbf{x}'', t; \mathbf{x}', t_0) = \langle \mathbf{x}''|\exp\left[\frac{-iH(t - t_0)}{\hbar}\right]|\mathbf{x}'\rangle}
\tag{2.274}
\end{equation}

where the time-evolution operator acting on $|\mathbf{x}'\rangle$ is just the state ket at $t$ of a system that was localized precisely at $\mathbf{x}'$ at time $t_0$ ($< t$).

If we wish to solve a more general problem where the initial wave function extends over a finite region of space, all we have to do is multiply $\psi(\mathbf{x}', t_0)$ by the propagator $K(\mathbf{x}'', t; \mathbf{x}', t_0)$ and integrate over all space (that is, over $\mathbf{x}'$). In this manner we can add the various contributions from different positions ($\mathbf{x}'$). This situation is analogous to one in electrostatics; if we wish to find the electrostatic potential due to a general charge distribution $\rho(\mathbf{x}')$, we first solve the point-charge problem, multiply the point-charge solution with the charge distribution, and integrate:

\begin{equation}
\phi(\mathbf{x}) = \int d^3x'\,\frac{\rho(\mathbf{x}')}{|\mathbf{x} - \mathbf{x}'|}
\tag{2.275}
\end{equation}

\subsubsection{Propagator as Green's Function}

The reader familiar with the theory of Green functions must have recognized by this time that the propagator is simply the Green function for the time-dependent wave equation satisfying

\begin{equation}
\left[-\left(\frac{\hbar^2}{2m}\right)\nabla''^2 + V(\mathbf{x}'') - i\hbar\frac{\partial}{\partial t}\right]K(\mathbf{x}'', t; \mathbf{x}', t_0) = -i\hbar\delta^3(\mathbf{x}'' - \mathbf{x}')\delta(t - t_0)
\tag{2.276}
\end{equation}

with the boundary condition

\begin{equation}
K(\mathbf{x}'', t; \mathbf{x}', t_0) = 0, \quad \text{for } t < t_0
\tag{2.277}
\end{equation}

The delta function $\delta(t - t_0)$ is needed on the right-hand side of (2.276) because $K$ varies discontinuously at $t = t_0$.

\subsubsection{Free Particle Propagator}

The particular form of the propagator is, of course, dependent on the particular potential to which the particle is subjected. Consider, as an example, a free particle in one dimension. The obvious observable that commutes with $H$ is momentum; $|p'\rangle$ is a simultaneous eigenket of the operators $p$ and $H$:

\begin{equation}
p|p'\rangle = p'|p'\rangle \qquad H|p'\rangle = \left(\frac{p'^2}{2m}\right)|p'\rangle
\tag{2.278}
\end{equation}

The momentum eigenfunction is just the transformation function of Section 1.7 [see (1.264)] which is of the plane wave form. Combining everything, we have

\begin{equation}
K(x'', t; x', t_0) = \left(\frac{1}{2\pi\hbar}\right)\int_{-\infty}^{\infty} dp'\,\exp\left[\frac{ip'(x'' - x')}{\hbar} - \frac{ip'^2(t - t_0)}{2m\hbar}\right]
\tag{2.279}
\end{equation}

The integral can be evaluated by completing the square in the exponent. Here we simply record the result:

\begin{equation}
\boxed{K(x'', t; x', t_0) = \sqrt{\frac{m}{2\pi i\hbar(t - t_0)}}\exp\left[\frac{im(x'' - x')^2}{2\hbar(t - t_0)}\right]}
\tag{2.280}
\end{equation}

This expression may be used, for example, to study how a Gaussian wave packet spreads out as a function of time.

\subsubsection{Simple Harmonic Oscillator Propagator}

For the simple harmonic oscillator, where the wave function of an energy eigenstate is given by

\begin{equation}
u_n(x)\exp\left(\frac{-iE_n t}{\hbar}\right) = \left(\frac{1}{2^{n/2}\sqrt{n!}}\right)\left(\frac{m\omega}{\pi\hbar}\right)^{1/4}\exp\left(\frac{-m\omega x^2}{2\hbar}\right) \times H_n\left(\sqrt{\frac{m\omega}{\hbar}}x\right)\exp\left[-i\omega\left(n + \frac{1}{2}\right)t\right]
\tag{2.281}
\end{equation}

the propagator is given by

\begin{equation}
\boxed{K(x'', t; x', t_0) = \sqrt{\frac{m\omega}{2\pi i\hbar\sin[\omega(t - t_0)]}}\exp\left\{\frac{im\omega}{2\hbar\sin[\omega(t - t_0)]}\left\{(x''^2 + x'^2)\cos[\omega(t - t_0)] - 2x''x'\right\}\right\}}
\tag{2.282}
\end{equation}

One way to prove this is to use

\begin{equation}
\left(\frac{1}{\sqrt{1 - \zeta^2}}\right)\exp\left[\frac{-(\xi^2 + \eta^2 - 2\xi\eta\zeta)}{(1 - \zeta^2)}\right] = \exp[-(\xi^2 + \eta^2)]\sum_{n=0}^{\infty}\left(\frac{\zeta^n}{2^n n!}\right)H_n(\xi)H_n(\eta)
\tag{2.283}
\end{equation}

Notice that (2.282) is, up to an overall sign, a periodic function of $t$ with angular frequency $\omega$, the classical oscillator frequency ($K \to -K$ under $t \to t + 2\pi/\omega$, so $K$ itself has period $4\pi/\omega$, while $|K|^2$ has period $2\pi/\omega$). This means, among other things, that a particle initially localized precisely at $x'$ will return to its original position with certainty at $2\pi/\omega$ ($4\pi/\omega$, and so forth) later.

\subsubsection{Connection to Statistical Mechanics}

Certain space and time integrals derivable from $K(\mathbf{x}'', t; \mathbf{x}', t_0)$ are of considerable interest. Without loss of generality we set $t_0 = 0$ in the following. The first integral we consider is obtained by setting $\mathbf{x}'' = \mathbf{x}'$ and integrating over all space. We have

\begin{align}
G(t) &\equiv \int d^3x'\,K(\mathbf{x}', t; \mathbf{x}', 0) \notag \\
&= \int d^3x'\sum_{a'}|\langle \mathbf{x}'|a'\rangle|^2\exp\left(\frac{-iE_{a'}t}{\hbar}\right) \notag \\
&= \sum_{a'}\exp\left(\frac{-iE_{a'}t}{\hbar}\right)
\tag{2.284}
\end{align}

This result is anticipated; recalling (2.274), we observe that setting $\mathbf{x}' = \mathbf{x}''$ and integrating are equivalent to taking the trace of the time-evolution operator in the $\mathbf{x}$-representation. But the trace is independent of representations; it can be evaluated more readily using the $\{|a'\rangle\}$ basis where the time-evolution operator is diagonal, which immediately leads to the last line of (2.284).

Now we see that (2.284) is just the ``sum over states,'' reminiscent of the partition function in statistical mechanics. In fact, if we analytically continue in the $t$ variable and make $t$ purely imaginary, with $\beta$ defined by

\begin{equation}
\beta = \frac{it}{\hbar}
\tag{2.285}
\end{equation}

real and positive, we can identify (2.284) with the partition function itself:

\begin{equation}
\boxed{Z = \sum_{a'}\exp(-\beta E_{a'})}
\tag{2.286}
\end{equation}

For this reason some of the techniques encountered in studying propagators in quantum mechanics are also useful in statistical mechanics.

\subsubsection{Energy Spectrum from the Propagator}

Next, let us consider the Laplace--Fourier transform of $G(t)$:

\begin{align}
\tilde{G}(E) &\equiv -i\int_0^{\infty} dt\,G(t)\exp(iEt/\hbar)/\hbar \notag \\
&= -i\int_0^{\infty} dt\sum_{a'}\exp(-iE_{a'}t/\hbar)\exp(iEt/\hbar)
\tag{2.287}
\end{align}

The integrand here oscillates indefinitely. But we can make the integral meaningful by letting $E$ acquire a small positive imaginary part:

\begin{equation}
E \to E + i\varepsilon
\tag{2.288}
\end{equation}

We then obtain, in the limit $\varepsilon \to 0$,

\begin{equation}
\boxed{\tilde{G}(E) = \sum_{a'}\frac{1}{E - E_{a'}}}
\tag{2.289}
\end{equation}

Observe now that the complete energy spectrum is exhibited as simple poles of $\tilde{G}(E)$ in the complex $E$-plane. If we wish to know the energy spectrum of a physical system, it is sufficient to study the analytic properties of $\tilde{G}(E)$.

\subsection{Propagator as a Transition Amplitude}

To gain further insight into the physical meaning of the propagator, we wish to relate it to the concept of transition amplitudes introduced in Section 2.2. But first, recall that the wave function which is the inner product of the fixed position bra $\langle \mathbf{x}'|$ with the moving state ket $|\alpha, t_0; t\rangle$ can also be regarded as the inner product of the Heisenberg picture position bra $\langle \mathbf{x}', t|$, which moves ``oppositely'' with time, with the Heisenberg picture state ket $|\alpha, t_0\rangle$, which is fixed in time.

Likewise, the propagator can also be written as

\begin{align}
K(\mathbf{x}'', t; \mathbf{x}', t_0) &= \sum_{a'} \langle \mathbf{x}''|a'\rangle\langle a'|\mathbf{x}'\rangle \exp\left[\frac{-iE_{a'}(t - t_0)}{\hbar}\right] \notag \\
&= \sum_{a'} \langle \mathbf{x}''|\exp\left(\frac{-iHt}{\hbar}\right)|a'\rangle\langle a'|\exp\left(\frac{iHt_0}{\hbar}\right)|\mathbf{x}'\rangle \notag \\
&= \langle \mathbf{x}'', t|\mathbf{x}', t_0\rangle
\tag{2.290}
\end{align}

where $|\mathbf{x}', t_0\rangle$ and $\langle \mathbf{x}'', t|$ are to be understood as an eigenket and an eigenbra of the position operator in the Heisenberg picture.

\subsubsection{Interpretation as Transition Amplitude}

In Section 2.2 we showed that $\langle b', t|a'\rangle$, in the Heisenberg picture notation, is the probability amplitude for a system originally prepared to be an eigenstate of $A$ with eigenvalue $a'$ at some initial time $t_0 = 0$ to be found at a later time $t$ in an eigenstate of $B$ with eigenvalue $b'$, and we called it the transition amplitude for going from state $|a'\rangle$ to state $|b'\rangle$.

Because there is nothing special about the choice of $t_0$, only the time difference $t - t_0$ is relevant, we can identify $\langle \mathbf{x}'', t|\mathbf{x}', t_0\rangle$ as the probability amplitude for the particle prepared at $t_0$ with position eigenvalue $\mathbf{x}'$ to be found at a later time $t$ at $\mathbf{x}''$. Roughly speaking, $\langle \mathbf{x}'', t|\mathbf{x}', t_0\rangle$ is the amplitude for the particle to go from a space-time point $(\mathbf{x}', t_0)$ to another space-time point $(\mathbf{x}'', t)$, so the term \textbf{transition amplitude} for this expression is quite appropriate. This interpretation is, of course, in complete accord with the interpretation we gave earlier for $K(\mathbf{x}'', t; \mathbf{x}', t_0)$.

\subsubsection{Alternative Interpretation: Unitary Transformation}

Yet another way to interpret $\langle \mathbf{x}'', t|\mathbf{x}', t_0\rangle$ is as follows. As we emphasized earlier, $|\mathbf{x}', t_0\rangle$ is the position eigenket at $t_0$ with the eigenvalue $\mathbf{x}'$ in the Heisenberg picture. Because at any given time the Heisenberg picture eigenkets of an observable can be chosen as base kets, we can regard $\langle \mathbf{x}'', t|\mathbf{x}', t_0\rangle$ as the transformation function that connects the two sets of base kets at \textit{different} times.

So in the Heisenberg picture, time evolution can be viewed as a \textit{unitary transformation}, in the sense of changing bases, that connects one set of base kets formed by $\{|\mathbf{x}', t_0\rangle\}$ to another formed by $\{|\mathbf{x}'', t\rangle\}$. This is reminiscent of classical physics, in which the time development of a classical dynamic variable such as $\mathbf{x}(t)$ is viewed as a canonical (or contact) transformation generated by the classical Hamiltonian.

\subsubsection{Symmetric Notation and Composition Property}

It turns out to be convenient to use a notation that treats the space and time coordinates more symmetrically. To this end we write $\langle \mathbf{x}'', t''|\mathbf{x}', t'\rangle$ in place of $\langle \mathbf{x}'', t|\mathbf{x}', t_0\rangle$. Because at any given time the position kets in the Heisenberg picture form a complete set, it is legitimate to insert the identity operator written as

\begin{equation}
\int d^3x''|\mathbf{x}'', t''\rangle\langle \mathbf{x}'', t''| = 1
\tag{2.291}
\end{equation}

at any place we desire.

\textbf{Composition property:}
For example, consider the time evolution from $t'$ to $t'''$; by dividing the time interval $(t', t''')$ into two parts, $(t', t'')$ and $(t'', t''')$, we have

\begin{equation}
\boxed{\langle \mathbf{x}''', t'''|\mathbf{x}', t'\rangle = \int d^3x''\,\langle \mathbf{x}''', t'''|\mathbf{x}'', t''\rangle\langle \mathbf{x}'', t''|\mathbf{x}', t'\rangle \qquad (t''' > t'' > t')}
\tag{2.292}
\end{equation}

We call this the \textbf{composition property} of the transition amplitude. Clearly, we can divide the time interval into as many smaller subintervals as we wish. We have

\begin{equation}
\langle \mathbf{x}'''', t''''|\mathbf{x}', t'\rangle = \int d^3x'''\int d^3x''\,\langle \mathbf{x}'''', t''''|\mathbf{x}''', t'''\rangle\langle \mathbf{x}''', t'''|\mathbf{x}'', t''\rangle \times \langle \mathbf{x}'', t''|\mathbf{x}', t'\rangle \quad (t'''' > t''' > t'' > t')
\tag{2.293}
\end{equation}

and so on. If we somehow guess the form of $\langle \mathbf{x}'', t''|\mathbf{x}', t'\rangle$ for an \textit{infinitesimal} time interval (between $t'$ and $t'' = t' + dt$), we should be able to obtain the amplitude $\langle \mathbf{x}'', t''|\mathbf{x}', t'\rangle$ for a finite time interval by compounding the appropriate transition amplitudes for infinitesimal time intervals in a manner analogous to (2.293). This kind of reasoning leads to an \textit{independent formulation} of quantum mechanics due to R. P. Feynman, published in 1948, to which we now turn our attention.

\subsection{Path Integrals as the Sum over Paths}

Without loss of generality we restrict ourselves to one-dimensional problems. Also, we avoid awkward expressions like $x''''...$ ($N$ times) by using notation such as $x_N$.

With this notation we consider the transition amplitude for a particle going from the initial space-time point $(x_1, t_1)$ to the final space-time point $(x_N, t_N)$. The entire time interval between $t_1$ and $t_N$ is divided into $N - 1$ equal parts:

\begin{equation}
t_j - t_{j-1} = \Delta t = \frac{(t_N - t_1)}{(N - 1)}
\tag{2.294}
\end{equation}

Exploiting the composition property, we obtain

\begin{equation}
\langle x_N, t_N|x_1, t_1\rangle = \int dx_{N-1}\int dx_{N-2}\cdots\int dx_2\,\langle x_N, t_N|x_{N-1}, t_{N-1}\rangle \times \langle x_{N-1}, t_{N-1}|x_{N-2}, t_{N-2}\rangle\cdots\langle x_2, t_2|x_1, t_1\rangle
\tag{2.295}
\end{equation}

To visualize this pictorially, we consider a space-time plane. The initial and final space-time points are fixed to be $(x_1, t_1)$ and $(x_N, t_N)$, respectively. For each time segment, say between $t_{n-1}$ and $t_n$, we are instructed to consider the transition amplitude to go from $(x_{n-1}, t_{n-1})$ to $(x_n, t_n)$; we then integrate over $x_2, x_3, \ldots, x_{N-1}$. This means that we must \textit{sum over all possible paths} in the space-time plane with the end points fixed.

\subsubsection{Classical Mechanics Review}

Before proceeding further, it is profitable to review here how paths appear in classical mechanics. Suppose we have a particle subjected to a force field derivable from a potential $V(x)$. The \textit{classical} Lagrangian is written as

\begin{equation}
L_{\text{classical}}(x, \dot{x}) = \frac{m\dot{x}^2}{2} - V(x)
\tag{2.296}
\end{equation}

Given this Lagrangian with the end points $(x_1, t_1)$ and $(x_N, t_N)$ specified, we do \textit{not} consider just any path joining $(x_1, t_1)$ and $(x_N, t_N)$ in classical mechanics. On the contrary, there exists a \textit{unique path} that corresponds to the actual motion of the classical particle.

More generally, according to Hamilton's principle, the unique path is that which makes the action stationary (an extremum, not necessarily a minimum: for the oscillator, paths longer than half a period are saddle points), defined as the time integral of the classical Lagrangian:

\begin{equation}
\delta\int_{t_1}^{t_2} dt\,L_{\text{classical}}(x, \dot{x}) = 0
\tag{2.299}
\end{equation}

from which Lagrange's equation of motion can be obtained.

\subsection{Feynman's Formulation}

The basic difference between classical mechanics and quantum mechanics should now be apparent. In classical mechanics a definite path in the $xt$-plane is associated with the particle's motion; in contrast, in quantum mechanics all possible paths must play roles including those which do not bear any resemblance to the classical path. Yet we must somehow be able to reproduce classical mechanics in a smooth manner in the limit $\hbar \to 0$. How are we to accomplish this?

As a young graduate student at Princeton University, R. P. Feynman tried to attack this problem. In looking for a possible clue, he was said to be intrigued by a mysterious remark in Dirac's book which, in our notation, amounts to the following statement:

\begin{equation*}
\exp\left[i\int_{t_1}^{t_2}\frac{dt\,L_{\text{classical}}(x, \dot{x})}{\hbar}\right] \quad \text{corresponds to} \quad \langle x_2, t_2|x_1, t_1\rangle
\end{equation*}

Feynman attempted to make sense out of this remark. Is ``corresponds to'' the same thing as ``is equal to'' or ``is proportional to''? In so doing he was led to formulate a space-time approach to quantum mechanics based on \textit{path integrals}.

\subsubsection{The Classical Action}

In Feynman's formulation the classical action plays a very important role. For compactness, we introduce a new notation:

\begin{equation}
S(n, n-1) \equiv \int_{t_{n-1}}^{t_n} dt\,L_{\text{classical}}(x, \dot{x})
\tag{2.300}
\end{equation}

Because $L_{\text{classical}}$ is a function of $x$ and $\dot{x}$, $S(n, n-1)$ is defined only after a definite path is specified along which the integration is to be carried out. So even though the path dependence is not explicit in this notation, it is understood that we are considering a particular path in evaluating the integral.

Imagine now that we are following some prescribed path. We concentrate our attention on a small segment along that path, say between $(x_{n-1}, t_{n-1})$ and $(x_n, t_n)$. According to Dirac, we are instructed to associate $\exp[iS(n, n-1)/\hbar]$ with that segment. Going along the definite path we are set to follow, we successively multiply expressions of this type to obtain

\begin{equation}
\prod_{n=2}^{N}\exp\left[\frac{iS(n, n-1)}{\hbar}\right] = \exp\left[\left(\frac{i}{\hbar}\right)\sum_{n=2}^{N}S(n, n-1)\right] = \exp\left[\frac{iS(N, 1)}{\hbar}\right]
\tag{2.301}
\end{equation}

This does not yet give $\langle x_N, t_N|x_1, t_1\rangle$; rather, this equation is the contribution to $\langle x_N, t_N|x_1, t_1\rangle$ arising from the particular path we have considered. We must still integrate over $x_2, x_3, \ldots, x_{N-1}$. At the same time, exploiting the composition property, we let the time interval between $t_{n-1}$ and $t_n$ be infinitesimally small. Thus our candidate expression for $\langle x_N, t_N|x_1, t_1\rangle$ may be written, in some loose sense, as

\begin{equation}
\langle x_N, t_N|x_1, t_1\rangle \sim \sum_{\text{all paths}}\exp\left[\frac{iS(N, 1)}{\hbar}\right]
\tag{2.302}
\end{equation}

where the sum is to be taken over an innumerably infinite set of paths!

\subsubsection{The Classical Limit}

Before presenting a more precise formulation, let us see whether considerations along this line make sense in the classical limit. As $\hbar \to 0$, the exponential in (2.302) oscillates very violently, so there is a tendency for cancellation among various contributions from neighboring paths. This is because $\exp[iS/\hbar]$ for some definite path and $\exp[iS/\hbar]$ for a slightly different path have very different phases because of the smallness of $\hbar$. So most paths do \textit{not} contribute when $\hbar$ is regarded as a small quantity. However, there is an important exception.

Suppose that we consider a path that satisfies

\begin{equation}
\delta S(N, 1) = 0
\tag{2.303}
\end{equation}

where the change in $S$ is due to a slight deformation of the path with the end points fixed. This is precisely the classical path by virtue of Hamilton's principle. We denote the $S$ that satisfies (2.303) by $S_{\min}$. We now attempt to deform the path a little bit from the classical path. The resulting $S$ is still equal to $S_{\min}$ to first order in deformation. This means that the phase of $\exp[iS/\hbar]$ does not vary very much as we deviate slightly from the classical path even if $\hbar$ is small. As a result, \textbf{constructive interference between neighboring paths is possible}. In the $\hbar \to 0$ limit, the major contributions must then arise from a very narrow strip (or a tube in higher dimensions) containing the classical path.

Our (or Feynman's) guess based on Dirac's mysterious remark makes good sense because the classical path is singled out in the $\hbar \to 0$ limit.

\subsubsection{Infinitesimal Time Interval}

To formulate Feynman's conjecture more precisely, let us go back to $\langle x_n, t_n|x_{n-1}, t_{n-1}\rangle$, where the time difference $t_n - t_{n-1}$ is assumed to be infinitesimally small. We write

\begin{equation}
\langle x_n, t_n|x_{n-1}, t_{n-1}\rangle = \left[\frac{1}{w(\Delta t)}\right]\exp\left[\frac{iS(n, n-1)}{\hbar}\right]
\tag{2.304}
\end{equation}

where we evaluate $S(n, n-1)$ in a moment in the $\Delta t \to 0$ limit. Notice that we have inserted a weight factor, $1/w(\Delta t)$, which is assumed to depend only on the time interval $t_n - t_{n-1}$ and not on $V(x)$. That such a factor is needed is clear from dimensional considerations; according to the way we normalized our position eigenkets, $\langle x_n, t_n|x_{n-1}, t_{n-1}\rangle$ must have the dimension of 1/length.

We now look at the exponential in (2.304). Our task is to evaluate the $\Delta t \to 0$ limit of $S(n, n-1)$. Because the time interval is so small, it is legitimate to make a straight-line approximation to the path joining $(x_{n-1}, t_{n-1})$ and $(x_n, t_n)$ as follows:

\begin{equation}
S(n, n-1) = \int_{t_{n-1}}^{t_n} dt\left[\frac{m\dot{x}^2}{2} - V(x)\right] = \Delta t\left\{\left(\frac{m}{2}\right)\left[\frac{(x_n - x_{n-1})}{\Delta t}\right]^2 - V\left(\frac{(x_n + x_{n-1})}{2}\right)\right\}
\tag{2.305}
\end{equation}

\textbf{Free particle case:}
As an example, we consider specifically the free-particle case, $V = 0$. Equation (2.304) now becomes

\begin{equation}
\langle x_n, t_n|x_{n-1}, t_{n-1}\rangle = \left[\frac{1}{w(\Delta t)}\right]\exp\left[\frac{im(x_n - x_{n-1})^2}{2\hbar\Delta t}\right]
\tag{2.306}
\end{equation}

We see that the exponent appearing here is completely identical to the one in the expression for the free-particle propagator (2.280).

\subsubsection{The Weight Factor}

We remarked earlier that the weight factor $1/w(\Delta t)$ appearing in (2.304) is assumed to be independent of $V(x)$, so we may as well evaluate it for the free particle. Noting the orthonormality, in the sense of $\delta$-function, of Heisenberg picture position eigenkets at equal times,

\begin{equation}
\langle x_n, t_n|x_{n-1}, t_{n-1}\rangle|_{t_n = t_{n-1}} = \delta(x_n - x_{n-1})
\tag{2.307}
\end{equation}

we obtain

\begin{equation}
\boxed{\frac{1}{w(\Delta t)} = \sqrt{\frac{m}{2\pi i\hbar\Delta t}}}
\tag{2.308}
\end{equation}

where we have used

\begin{equation}
\int_{-\infty}^{\infty} d\xi\,\exp\left(\frac{im\xi^2}{2\hbar\Delta t}\right) = \sqrt{\frac{2\pi i\hbar\Delta t}{m}}
\tag{2.309a}
\end{equation}

and

\begin{equation}
\lim_{\Delta t \to 0}\sqrt{\frac{m}{2\pi i\hbar\Delta t}}\exp\left(\frac{im\xi^2}{2\hbar\Delta t}\right) = \delta(\xi)
\tag{2.309b}
\end{equation}

This weight factor is, of course, anticipated from the expression for the free-particle propagator (2.280).

\subsubsection{The Path Integral}

To summarize, as $\Delta t \to 0$, we are led to

\begin{equation}
\langle x_n, t_n|x_{n-1}, t_{n-1}\rangle = \sqrt{\frac{m}{2\pi i\hbar\Delta t}}\exp\left[\frac{iS(n, n-1)}{\hbar}\right]
\tag{2.310}
\end{equation}

The final expression for the transition amplitude with $t_N - t_1$ finite is

\begin{equation}
\langle x_N, t_N|x_1, t_1\rangle = \lim_{N \to \infty}\left(\frac{m}{2\pi i\hbar\Delta t}\right)^{(N-1)/2} \times \int dx_{N-1}\int dx_{N-2}\cdots\int dx_2\prod_{n=2}^{N}\exp\left[\frac{iS(n, n-1)}{\hbar}\right]
\tag{2.311}
\end{equation}

where the $N \to \infty$ limit is taken with $x_N$ and $t_N$ fixed. It is customary here to define a new kind of multidimensional (in fact, infinite-dimensional) integral operator

\begin{equation}
\int_{x_1}^{x_N}\mathscr{D}[x(t)] \equiv \lim_{N \to \infty}\left(\frac{m}{2\pi i\hbar\Delta t}\right)^{(N-1)/2}\int dx_{N-1}\int dx_{N-2}\cdots\int dx_2
\tag{2.312}
\end{equation}

and write (2.311) as

\begin{equation}
\boxed{\langle x_N, t_N|x_1, t_1\rangle = \int_{x_1}^{x_N}\mathscr{D}[x(t)]\exp\left[i\int_{t_1}^{t_N}dt\,\frac{L_{\text{classical}}(x, \dot{x})}{\hbar}\right]}
\tag{2.313}
\end{equation}

This expression is known as \textbf{Feynman's path integral}. Its meaning as the sum over all possible paths should be apparent from (2.311).

\subsubsection{Equivalence to Schrödinger's Wave Mechanics}

Our steps leading to (2.313) are not meant to be a derivation. Rather, we (or Feynman) have attempted a new formulation of quantum mechanics based on the concept of paths, motivated by Dirac's mysterious remark. The only ideas we borrowed from the conventional form of quantum mechanics are (1) the superposition principle (used in summing the contributions from various alternate paths), (2) the composition property of the transition amplitude, and (3) classical correspondence in the $\hbar \to 0$ limit.

Even though we obtained the same result as the conventional theory for the free-particle case, it is not obvious, from what we have done so far, that Feynman's formulation is completely equivalent to Schrödinger's wave mechanics. We conclude this section by proving that $\langle x_N, t_N|x_1, t_1\rangle$ constructed according to Feynman's prescription indeed satisfies Schrödinger's time-dependent wave equation in the variables $x_N$, $t_N$, just as the propagator defined by (2.272).

We start with

\begin{align}
\langle x_N, t_N|x_1, t_1\rangle &= \int dx_{N-1}\,\langle x_N, t_N|x_{N-1}, t_{N-1}\rangle\langle x_{N-1}, t_{N-1}|x_1, t_1\rangle \notag \\
&= \int_{-\infty}^{\infty} dx_{N-1}\sqrt{\frac{m}{2\pi i\hbar\Delta t}}\exp\left[\left(\frac{im}{2\hbar}\right)\frac{(x_N - x_{N-1})^2}{\Delta t} - \frac{iV\Delta t}{\hbar}\right] \times \langle x_{N-1}, t_{N-1}|x_1, t_1\rangle
\tag{2.314}
\end{align}

where we have assumed $t_N - t_{N-1}$ to be infinitesimal. Introducing

\begin{equation}
\xi = x_N - x_{N-1}
\tag{2.315}
\end{equation}

and letting $x_N \to x$ and $t_N \to t + \Delta t$, we obtain

\begin{equation}
\langle x, t + \Delta t|x_1, t_1\rangle = \sqrt{\frac{m}{2\pi i\hbar\Delta t}}\int_{-\infty}^{(\infty)} d\xi\,\exp\left(\frac{im\xi^2}{2\hbar\Delta t} - \frac{iV\Delta t}{\hbar}\right)\langle x - \xi, t|x_1, t_1\rangle
\tag{2.316}
\end{equation}

As is evident from (2.309b), in the limit $\Delta t \to 0$, the major contribution to this integral comes from the $\xi \simeq 0$ region. It is therefore legitimate to expand $\langle x - \xi, t|x_1, t_1\rangle$ in powers of $\xi$. We also expand $\langle x, t + \Delta t|x_1, t_1\rangle$ and $\exp(-iV\Delta t/\hbar)$ in powers of $\Delta t$, so

\begin{align}
\langle x, t|x_1, t_1\rangle + \Delta t\frac{\partial}{\partial t}\langle x, t|x_1, t_1\rangle &= \sqrt{\frac{m}{2\pi i\hbar\Delta t}}\int_{-\infty}^{\infty} d\xi\,\exp\left(\frac{im\xi^2}{2\hbar\Delta t}\right)\left(1 - \frac{iV\Delta t}{\hbar} + \cdots\right) \notag \\
&\quad \times \left[\langle x, t|x_1, t_1\rangle + \left(\frac{\xi^2}{2}\right)\frac{\partial^2}{\partial x^2}\langle x, t|x_1, t_1\rangle + \cdots\right]
\tag{2.317}
\end{align}

where we have dropped a term linear in $\xi$ because it vanishes when integrated with respect to $\xi$. The $\langle x, t|x_1, t_1\rangle$ term on the left-hand side just matches the leading term on the right-hand side because of (2.309a). Collecting terms first order in $\Delta t$, we obtain

\begin{align}
\Delta t\frac{\partial}{\partial t}\langle x, t|x_1, t_1\rangle &= \left(\sqrt{\frac{m}{2\pi i\hbar\Delta t}}\right)(\sqrt{2\pi})\left(\frac{i\hbar\Delta t}{m}\right)^{3/2}\frac{1}{2}\frac{\partial^2}{\partial x^2}\langle x, t|x_1, t_1\rangle \notag \\
&\quad - \left(\frac{i}{\hbar}\right)\Delta t V\langle x, t|x_1, t_1\rangle
\tag{2.318}
\end{align}

where we have used

\begin{equation}
\int_{-\infty}^{\infty} d\xi\,\xi^2\exp\left(\frac{im\xi^2}{2\hbar\Delta t}\right) = \sqrt{2\pi}\left(\frac{i\hbar\Delta t}{m}\right)^{3/2}
\tag{2.319}
\end{equation}

obtained by differentiating (2.309a) with respect to $\Delta t$. In this manner we see that $\langle x, t|x_1, t_1\rangle$ satisfies Schrödinger's time-dependent wave equation:

\begin{equation}
\boxed{i\hbar\frac{\partial}{\partial t}\langle x, t|x_1, t_1\rangle = -\left(\frac{\hbar^2}{2m}\right)\frac{\partial^2}{\partial x^2}\langle x, t|x_1, t_1\rangle + V\langle x, t|x_1, t_1\rangle}
\tag{2.320}
\end{equation}

Thus we can conclude that $\langle x, t|x_1, t_1\rangle$ constructed according to Feynman's prescription is the same as the propagator in Schrödinger's wave mechanics.

\subsubsection{Significance of the Path Integral Formulation}

Feynman's space-time approach based on path integrals is not too convenient for attacking practical problems in nonrelativistic quantum mechanics. Even for the simple harmonic oscillator it is rather cumbersome to evaluate explicitly the relevant path integral. However, his approach is extremely gratifying from a conceptual point of view. By imposing a certain set of sensible requirements on a physical theory, we are inevitably led to a formalism equivalent to the usual formulation of quantum mechanics. It makes us wonder whether it is at all possible to construct a sensible alternative theory that is equally successful in accounting for microscopic phenomena.

Methods based on path integrals have been found to be very powerful in other branches of modern physics, such as quantum field theory and statistical mechanics.

\newpage
\section*{Theory of Angular Momentum}
\addcontentsline{toc}{section}{Theory of Angular Momentum}

This chapter is concerned with a systematic treatment of angular momentum and related topics. The importance of angular momentum in modern physics can hardly be overemphasized. A thorough understanding of angular momentum is essential in molecular, atomic, and nuclear spectroscopy; angular momentum considerations play an important role in scattering and collision problems as well as in bound-state problems. Furthermore, angular momentum concepts have important generalizations -- isospin in nuclear physics, SU(3), SU(2)$\otimes$U(1) in particle physics, and so forth.

\section{Rotations and Angular Momentum Commutation Relations}

\subsection{Finite Versus Infinitesimal Rotations}

We recall from elementary physics that rotations about the same axis commute, whereas rotations about different axes do not. For instance, a 30° rotation about the $z$-axis followed by a 60° rotation about the same $z$-axis is obviously equivalent to a 60° rotation followed by a 30° rotation, both about the same axis. However, let us consider a 90° rotation about the $z$-axis, denoted by $R_z(\pi/2)$, followed by a 90° rotation about the $x$-axis, denoted by $R_x(\pi/2)$; compare this with a 90° rotation about the $x$-axis followed by a 90° rotation about the $z$-axis. The net results are different.

Our first basic task is to work out quantitatively the manner in which rotations about different axes \textit{fail} to commute.

\subsubsection{Rotation Matrices}

To this end, we first recall how to represent rotations in three dimensions by $3 \times 3$ real, orthogonal matrices. Consider a vector $\mathbf{V}$ with components $V_x$, $V_y$, and $V_z$. When we rotate, the three components become some other set of numbers, $V_x'$, $V_y'$, and $V_z'$. The old and new components are related via a $3 \times 3$ orthogonal matrix $R$:

\begin{align}
\begin{pmatrix} V_x' \\ V_y' \\ V_z' \end{pmatrix} &= (R)\begin{pmatrix} V_x \\ V_y \\ V_z \end{pmatrix} \tag{3.1a} \\
RR^T &= R^T R = 1 \tag{3.1b}
\end{align}

where the superscript $T$ stands for a transpose of a matrix. It is a property of orthogonal matrices that

\begin{equation}
\sqrt{V_x^2 + V_y^2 + V_z^2} = \sqrt{V_x'^2 + V_y'^2 + V_z'^2}
\tag{3.2}
\end{equation}

is automatically satisfied.

\subsubsection{Rotation About the $z$-Axis}

To be definite, we consider a rotation about the $z$-axis by angle $\phi$. The convention we follow throughout this book is that a rotation operation affects a physical system itself, while the coordinate axes remain unchanged. The angle $\phi$ is taken to be positive when the rotation in question is counterclockwise in the $xy$-plane, as viewed from the positive $z$-side. If we associate a right-handed screw with such a rotation, a positive $\phi$ rotation around the $z$-axis means that the screw is advancing in the positive $z$-direction.

With this convention, we easily verify that

\begin{equation}
\boxed{R_z(\phi) = \begin{pmatrix} \cos\phi & -\sin\phi & 0 \\ \sin\phi & \cos\phi & 0 \\ 0 & 0 & 1 \end{pmatrix}}
\tag{3.3}
\end{equation}

Had we adopted a different convention, in which a physical system remained fixed but the coordinate axes rotated, this same matrix with a positive $\phi$ would have represented a \textit{clockwise} rotation of the $x$- and $y$-axes, when viewed from the positive $z$-side. It is obviously important not to mix the two conventions! Some authors distinguish the two approaches by using ``active rotations'' for physical systems rotated and ``passive rotations'' for the coordinate axes rotated.

\subsubsection{Infinitesimal Rotations}

We are particularly interested in an infinitesimal form of $R_z$:

\begin{equation}
R_z(\varepsilon) = \begin{pmatrix} 1 - \frac{\varepsilon^2}{2} & -\varepsilon & 0 \\ \varepsilon & 1 - \frac{\varepsilon^2}{2} & 0 \\ 0 & 0 & 1 \end{pmatrix}
\tag{3.4}
\end{equation}

where terms of order $\varepsilon^3$ and higher are ignored. Likewise, we have

\begin{equation}
R_x(\varepsilon) = \begin{pmatrix} 1 & 0 & 0 \\ 0 & 1 - \frac{\varepsilon^2}{2} & -\varepsilon \\ 0 & \varepsilon & 1 - \frac{\varepsilon^2}{2} \end{pmatrix}
\tag{3.5a}
\end{equation}

and

\begin{equation}
R_y(\varepsilon) = \begin{pmatrix} 1 - \frac{\varepsilon^2}{2} & 0 & \varepsilon \\ 0 & 1 & 0 \\ -\varepsilon & 0 & 1 - \frac{\varepsilon^2}{2} \end{pmatrix}
\tag{3.5b}
\end{equation}

which may be read from (3.4) by cyclic permutations of $x$, $y$, $z$, that is, $x \to y$, $y \to z$, $z \to x$.

\subsubsection{Non-Commutativity of Rotations}

Compare now the effect of a $y$-axis rotation followed by an $x$-axis rotation with that of an $x$-axis rotation followed by a $y$-axis rotation. Elementary matrix manipulations lead to

\begin{equation}
R_x(\varepsilon)R_y(\varepsilon) = \begin{pmatrix} 1 - \frac{\varepsilon^2}{2} & 0 & \varepsilon \\ \varepsilon^2 & 1 - \frac{\varepsilon^2}{2} & -\varepsilon \\ -\varepsilon & \varepsilon & 1 - \varepsilon^2 \end{pmatrix}
\tag{3.6a}
\end{equation}

and

\begin{equation}
R_y(\varepsilon)R_x(\varepsilon) = \begin{pmatrix} 1 - \frac{\varepsilon^2}{2} & \varepsilon^2 & \varepsilon \\ 0 & 1 - \frac{\varepsilon^2}{2} & -\varepsilon \\ -\varepsilon & \varepsilon & 1 - \varepsilon^2 \end{pmatrix}
\tag{3.6b}
\end{equation}

From (3.6a) and (3.6b) we have the first important result: \textbf{Infinitesimal rotations about different axes do commute if terms of order $\varepsilon^2$ and higher are ignored.}

The second and even more important result concerns the manner in which rotations about different axes \textit{fail} to commute when terms of order $\varepsilon^2$ are kept:

\begin{equation}
R_x(\varepsilon)R_y(\varepsilon) - R_y(\varepsilon)R_x(\varepsilon) = \begin{pmatrix} 0 & -\varepsilon^2 & 0 \\ \varepsilon^2 & 0 & 0 \\ 0 & 0 & 0 \end{pmatrix} = R_z(\varepsilon^2) - 1
\tag{3.7}
\end{equation}

where all terms of order higher than $\varepsilon^2$ have been ignored throughout this derivation. We also have

\begin{equation}
1 = R_{\text{any}}(0)
\tag{3.8}
\end{equation}

where \textit{any} stands for any rotation axis. Thus the final result can be written as

\begin{equation}
\boxed{R_x(\varepsilon)R_y(\varepsilon) - R_y(\varepsilon)R_x(\varepsilon) = R_z(\varepsilon^2) - R_{\text{any}}(0)}
\tag{3.9}
\end{equation}

This is an example of the commutation relations between rotation operations about different axes, which we will use later in deducing the angular-momentum commutation relations in quantum mechanics.

\subsection{Infinitesimal Rotations in Quantum Mechanics}

So far we have not used quantum-mechanical concepts. The matrix $R$ is just a $3 \times 3$ orthogonal matrix acting on a vector $\mathbf{V}$ written in column matrix form. We must now understand how to characterize rotations in quantum mechanics.

\subsubsection{The Rotation Operator}

Because rotations affect physical systems, the state ket corresponding to a rotated system is expected to look different from the state ket corresponding to the original unrotated system. Given a rotation operation $R$, characterized by a $3 \times 3$ orthogonal matrix $R$, we associate an operator $\mathscr{D}(R)$ in the appropriate ket space such that

\begin{equation}
\boxed{|\alpha\rangle_R = \mathscr{D}(R)|\alpha\rangle}
\tag{3.10}
\end{equation}

where $|\alpha\rangle_R$ and $|\alpha\rangle$ stand for the kets of the rotated and original system, respectively.

The matrix representation of $\mathscr{D}(R)$ depends on the dimensionality $N$ of the particular ket space in question. For $N = 2$, appropriate for describing a spin $\frac{1}{2}$ system with no other degrees of freedom, $\mathscr{D}(R)$ is represented by a $2 \times 2$ matrix; for a spin 1 system, the appropriate representation is a $3 \times 3$ unitary matrix, and so on.

\subsubsection{Angular Momentum as the Generator of Rotations}

By analogy with translations and time evolution, the appropriate infinitesimal operators can be written as

\begin{equation}
U_\varepsilon = 1 - iG\varepsilon
\tag{3.11}
\end{equation}

with a Hermitian operator $G$. Specifically:

\textbf{Translation:}
\begin{equation}
G \to \frac{p_x}{\hbar}, \quad \varepsilon \to dx'
\tag{3.12}
\end{equation}

\textbf{Time evolution:}
\begin{equation}
G \to \frac{H}{\hbar}, \quad \varepsilon \to dt
\tag{3.13}
\end{equation}

\textbf{Rotation:}
We \textit{define} the angular-momentum operator $J_k$ such that the operator for an infinitesimal rotation around the $k$th axis by angle $d\phi$ can be obtained by letting

\begin{equation}
G \to \frac{J_k}{\hbar}, \quad \varepsilon \to d\phi
\tag{3.14}
\end{equation}

With $J_k$ taken to be Hermitian, the infinitesimal rotation operator is guaranteed to be unitary and reduces to the identity operator in the limit $d\phi \to 0$.

\textbf{General infinitesimal rotation:}
\begin{equation}
\boxed{\mathscr{D}(\hat{\mathbf{n}}, d\phi) = 1 - i\left(\frac{\mathbf{J}\cdot\hat{\mathbf{n}}}{\hbar}\right)d\phi}
\tag{3.15}
\end{equation}

for a rotation about the direction characterized by a unit vector $\hat{\mathbf{n}}$ by an infinitesimal angle $d\phi$.

\textbf{Important:} In this book we do \textit{not} define the angular-momentum operator to be $\mathbf{x} \times \mathbf{p}$. This is important because spin angular momentum, to which our general formalism also applies, has nothing to do with $x_i$ and $p_j$. In quantum mechanics we \textit{define} $\mathbf{J}$ so that the operator for an infinitesimal rotation takes the form (3.15).

\subsection{Finite Rotations in Quantum Mechanics}

A finite rotation can be obtained by compounding successively infinitesimal rotations about the same axis. For instance, if we are interested in a finite rotation about the $z$-axis by angle $\phi$, we consider

\begin{equation}
\boxed{\mathscr{D}_z(\phi) = \lim_{N \to \infty}\left[1 - i\left(\frac{J_z}{\hbar}\right)\left(\frac{\phi}{N}\right)\right]^N = \exp\left(\frac{-iJ_z\phi}{\hbar}\right) = 1 - \frac{iJ_z\phi}{\hbar} - \frac{J_z^2\phi^2}{2\hbar^2} + \cdots}
\tag{3.16}
\end{equation}

\subsubsection{Group Properties}

In order to obtain the angular-momentum commutation relations, we need one more concept. For every rotation $R$ represented by a $3 \times 3$ orthogonal matrix $R$, there exists a rotation operator $\mathscr{D}(R)$ in the appropriate ket space. We postulate that $\mathscr{D}(R)$ has the same group properties as $R$:

\begin{align}
\text{Identity:} \qquad & R \cdot 1 = R \Rightarrow \mathscr{D}(R) \cdot 1 = \mathscr{D}(R) \tag{3.17a} \\
\text{Closure:} \qquad & R_1 R_2 = R_3 \Rightarrow \mathscr{D}(R_1)\mathscr{D}(R_2) = \mathscr{D}(R_3) \tag{3.17b} \\
\text{Inverses:} \qquad & RR^{-1} = 1 \Rightarrow \mathscr{D}(R)\mathscr{D}^{-1}(R) = 1 \tag{3.17c} \\
\text{Associativity:} \qquad & R_1(R_2 R_3) = (R_1 R_2)R_3 \Rightarrow \mathscr{D}(R_1)[\mathscr{D}(R_2)\mathscr{D}(R_3)] = [\mathscr{D}(R_1)\mathscr{D}(R_2)]\mathscr{D}(R_3) \tag{3.17d}
\end{align}

\subsection{Commutation Relations for Angular Momentum}

Let us now return to the fundamental commutation relations for rotation operations (3.9) written in terms of the $R$ matrices. Its rotation operator analogue would read

\begin{equation}
\left(1 - \frac{iJ_x\varepsilon}{\hbar} - \frac{J_x^2\varepsilon^2}{2\hbar^2}\right)\left(1 - \frac{iJ_y\varepsilon}{\hbar} - \frac{J_y^2\varepsilon^2}{2\hbar^2}\right) - \left(1 - \frac{iJ_y\varepsilon}{\hbar} - \frac{J_y^2\varepsilon^2}{2\hbar^2}\right)\left(1 - \frac{iJ_x\varepsilon}{\hbar} - \frac{J_x^2\varepsilon^2}{2\hbar^2}\right) = 1 - \frac{iJ_z\varepsilon^2}{\hbar} - 1
\tag{3.18}
\end{equation}

Terms of order $\varepsilon$ automatically drop out. Equating terms of order $\varepsilon^2$ on both sides of (3.18), we obtain

\begin{equation}
[J_x, J_y] = i\hbar J_z
\tag{3.19}
\end{equation}

Repeating this kind of argument with rotations about other axes, we obtain

\begin{equation}
\boxed{[J_i, J_j] = i\hbar\varepsilon_{ijk}J_k}
\tag{3.20}
\end{equation}

known as the \textbf{fundamental commutation relations of angular momentum}.

\textbf{Key points:}
\begin{itemize}[leftmargin=*]
\item Because of (3.20), the rotation group in three dimensions is \textbf{non-Abelian}.
\item In contrast, the translation group in three dimensions is Abelian because $p_i$ and $p_j$ commute even with $i \neq j$.
\item Commutation relations (3.20) summarize in a compact manner \textit{all} the basic properties of rotations in three dimensions.
\end{itemize}

\subsubsection{Transformation of Angular Momentum Operators}

The fundamental commutation relations make it possible to show that the angular-momentum operators themselves transform as expected under rotations. Consider a rotation by a finite angle $\phi$ about the $z$-axis. If the ket of the system before rotation is given by $|\alpha\rangle$, the ket after rotation is given by

\begin{equation}
|\alpha\rangle_R = \mathscr{D}_z(\phi)|\alpha\rangle
\tag{3.21}
\end{equation}

Under rotation, the expectation value $\langle J_x \rangle$ changes as follows:

\begin{equation}
\langle J_x \rangle \to {}_R\langle\alpha|J_x|\alpha\rangle_R = \langle\alpha|\mathscr{D}_z^\dagger(\phi)J_x\mathscr{D}_z(\phi)|\alpha\rangle
\tag{3.22}
\end{equation}

We must therefore compute

\begin{equation}
\exp\left(\frac{iJ_z\phi}{\hbar}\right)J_x\exp\left(\frac{-iJ_z\phi}{\hbar}\right)
\tag{3.23}
\end{equation}

Using the Baker-Hausdorff lemma (2.168) to evaluate (3.23), we have

\begin{align}
\exp\left(\frac{iJ_z\phi}{\hbar}\right)J_x\exp\left(\frac{-iJ_z\phi}{\hbar}\right) &= J_x + \left(\frac{i\phi}{\hbar}\right)\underbrace{[J_z, J_x]}_{i\hbar J_y} + \frac{1}{2!}\left(\frac{i\phi}{\hbar}\right)^2\underbrace{[J_z,[J_z,J_x]]}_{\hbar^2 J_x} + \cdots \notag \\
&= J_x\left[1 - \frac{\phi^2}{2!} + \cdots\right] - J_y\left[\phi - \frac{\phi^3}{3!} + \cdots\right] \notag \\
&= \boxed{J_x\cos\phi - J_y\sin\phi}
\tag{3.24}
\end{align}

This confirms that angular momentum operators transform as vector components under rotations.

\section{Spin $\frac{1}{2}$ Systems and Finite Rotations}

\subsection{Rotation Operator for Spin $\frac{1}{2}$}

The lowest number of dimensions in which the angular-momentum commutation relations (3.20) are realized is $N = 2$. The operators defined by

\begin{equation}
\boxed{S_x = \frac{\hbar}{2}\{|+\rangle\langle-| + |-\rangle\langle+|\}, \quad S_y = \frac{i\hbar}{2}\{-|+\rangle\langle-| + |-\rangle\langle+|\}, \quad S_z = \frac{\hbar}{2}\{|+\rangle\langle+| - |-\rangle\langle-|\}}
\tag{3.25}
\end{equation}

satisfy commutation relations (3.20) with $J_k$ replaced by $S_k$.

\subsubsection{Transformation of Spin Expectation Values}

Under a rotation about the $z$-axis by angle $\phi$:
\begin{align}
\langle S_x \rangle &\to \langle S_x \rangle\cos\phi - \langle S_y \rangle\sin\phi \tag{3.26} \\
\langle S_y \rangle &\to \langle S_y \rangle\cos\phi + \langle S_x \rangle\sin\phi \tag{3.27} \\
\langle S_z \rangle &\to \langle S_z \rangle \tag{3.28}
\end{align}

In general:
\begin{equation}
\langle J_k \rangle \to \sum_l R_{kl}\langle J_l \rangle
\tag{3.30}
\end{equation}

\subsubsection{The $2\pi$ Rotation Problem}

For a general ket $|\alpha\rangle = |+\rangle\langle+|\alpha\rangle + |-\rangle\langle-|\alpha\rangle$:

\begin{equation}
\exp\left(\frac{-iS_z\phi}{\hbar}\right)|\alpha\rangle = e^{-i\phi/2}|+\rangle\langle+|\alpha\rangle + e^{i\phi/2}|-\rangle\langle-|\alpha\rangle
\tag{3.32}
\end{equation}

\textbf{Key result:} For a $2\pi$ rotation:
\begin{equation}
\boxed{|\alpha\rangle_{R_z(2\pi)} \to -|\alpha\rangle}
\tag{3.33}
\end{equation}

A $4\pi$ ($\phi = 4\pi$) rotation is needed to get back to the same ket with a plus sign!

\subsection{Spin Precession Revisited}

The Hamiltonian for spin in a magnetic field:
\begin{equation}
H = -\left(\frac{e}{m_e c}\right)\mathbf{S}\cdot\mathbf{B} = \omega S_z, \quad \text{where } \omega \equiv \frac{|e|B}{m_e c}
\tag{3.34, 3.35}
\end{equation}

The time-evolution operator:
\begin{equation}
\mathscr{U}(t, 0) = \exp\left(\frac{-iHt}{\hbar}\right) = \exp\left(\frac{-iS_z\omega t}{\hbar}\right)
\tag{3.36}
\end{equation}

\textbf{Spin precession equations:}
\begin{align}
\langle S_x \rangle_t &= \langle S_x \rangle_{t=0}\cos\omega t - \langle S_y \rangle_{t=0}\sin\omega t \tag{3.37a} \\
\langle S_y \rangle_t &= \langle S_y \rangle_{t=0}\cos\omega t + \langle S_x \rangle_{t=0}\sin\omega t \tag{3.37b} \\
\langle S_z \rangle_t &= \langle S_z \rangle_{t=0} \tag{3.37c}
\end{align}

\textbf{Important periods:}
\begin{equation}
\boxed{\tau_{\text{precession}} = \frac{2\pi}{\omega}, \qquad \tau_{\text{state ket}} = \frac{4\pi}{\omega}}
\tag{3.39}
\end{equation}

The period for the state ket is \textit{twice} as long as the period for spin precession.

\subsection{Pauli Two-Component Formalism}

\subsubsection{Spinor Representation}

Base kets and bras:
\begin{equation}
|+\rangle \doteq \begin{pmatrix} 1 \\ 0 \end{pmatrix} \equiv \chi_+, \quad |-\rangle \doteq \begin{pmatrix} 0 \\ 1 \end{pmatrix} \equiv \chi_-
\tag{3.44}
\end{equation}

A general state ket is represented by a \textbf{two-component spinor}:
\begin{equation}
\chi = \begin{pmatrix} \langle+|\alpha\rangle \\ \langle-|\alpha\rangle \end{pmatrix} \equiv \begin{pmatrix} c_+ \\ c_- \end{pmatrix} = c_+\chi_+ + c_-\chi_-
\tag{3.46}
\end{equation}

\subsubsection{Pauli Matrices}

\begin{equation}
\boxed{\sigma_1 = \begin{pmatrix} 0 & 1 \\ 1 & 0 \end{pmatrix}, \quad \sigma_2 = \begin{pmatrix} 0 & -i \\ i & 0 \end{pmatrix}, \quad \sigma_3 = \begin{pmatrix} 1 & 0 \\ 0 & -1 \end{pmatrix}}
\tag{3.50}
\end{equation}

where subscripts 1, 2, and 3 refer to $x$, $y$, and $z$, respectively. The spin operators are:
\begin{equation}
\mathbf{S} = \frac{\hbar}{2}\boldsymbol{\sigma}
\end{equation}

\subsubsection{Properties of Pauli Matrices}

\textbf{Anticommutation relations:}
\begin{equation}
\boxed{\{\sigma_i, \sigma_j\} = 2\delta_{ij}}
\tag{3.52}
\end{equation}

\textbf{Commutation relations:}
\begin{equation}
\boxed{[\sigma_i, \sigma_j] = 2i\varepsilon_{ijk}\sigma_k}
\tag{3.53}
\end{equation}

\textbf{Other properties:}
\begin{align}
\sigma_i^2 &= 1 \tag{3.51a} \\
\sigma_1\sigma_2 &= -\sigma_2\sigma_1 = i\sigma_3 \tag{3.54} \\
\sigma_i^\dagger &= \sigma_i, \quad \det(\sigma_i) = -1, \quad \text{Tr}(\sigma_i) = 0 \tag{3.55}
\end{align}

\subsubsection{Important Identity}

For $\boldsymbol{\sigma}\cdot\mathbf{a}$ where $\mathbf{a}$ is a vector:
\begin{equation}
\boldsymbol{\sigma}\cdot\mathbf{a} = \begin{pmatrix} a_3 & a_1 - ia_2 \\ a_1 + ia_2 & -a_3 \end{pmatrix}
\tag{3.56}
\end{equation}

\textbf{Key identity:}
\begin{equation}
\boxed{(\boldsymbol{\sigma}\cdot\mathbf{a})(\boldsymbol{\sigma}\cdot\mathbf{b}) = \mathbf{a}\cdot\mathbf{b} + i\boldsymbol{\sigma}\cdot(\mathbf{a}\times\mathbf{b})}
\tag{3.57}
\end{equation}

If $\mathbf{a}$ is real:
\begin{equation}
(\boldsymbol{\sigma}\cdot\mathbf{a})^2 = |\mathbf{a}|^2
\tag{3.59}
\end{equation}

\subsection{Rotations in the Two-Component Formalism}

The $2 \times 2$ matrix representation of the rotation operator:
\begin{equation}
\exp\left(\frac{-i\mathbf{S}\cdot\hat{\mathbf{n}}\phi}{\hbar}\right) \doteq \exp\left(\frac{-i\boldsymbol{\sigma}\cdot\hat{\mathbf{n}}\phi}{2}\right)
\tag{3.60}
\end{equation}

Using $(\boldsymbol{\sigma}\cdot\hat{\mathbf{n}})^n = 1$ for $n$ even and $\boldsymbol{\sigma}\cdot\hat{\mathbf{n}}$ for $n$ odd:

\begin{equation}
\boxed{\exp\left(\frac{-i\boldsymbol{\sigma}\cdot\hat{\mathbf{n}}\phi}{2}\right) = \mathbf{1}\cos\left(\frac{\phi}{2}\right) - i\boldsymbol{\sigma}\cdot\hat{\mathbf{n}}\sin\left(\frac{\phi}{2}\right)}
\tag{3.62}
\end{equation}

Explicitly in $2 \times 2$ form:
\begin{equation}
\boxed{\exp\left(\frac{-i\boldsymbol{\sigma}\cdot\hat{\mathbf{n}}\phi}{2}\right) = \begin{pmatrix} \cos\frac{\phi}{2} - in_z\sin\frac{\phi}{2} & (-in_x - n_y)\sin\frac{\phi}{2} \\ (-in_x + n_y)\sin\frac{\phi}{2} & \cos\frac{\phi}{2} + in_z\sin\frac{\phi}{2} \end{pmatrix}}
\tag{3.63}
\end{equation}

Under rotations, the spinor transforms as:
\begin{equation}
\chi \to \exp\left(\frac{-i\boldsymbol{\sigma}\cdot\hat{\mathbf{n}}\phi}{2}\right)\chi
\tag{3.64}
\end{equation}

\textbf{$2\pi$ rotation in matrix form:}
\begin{equation}
\exp\left(\frac{-i\boldsymbol{\sigma}\cdot\hat{\mathbf{n}}\phi}{2}\right)\bigg|_{\phi=2\pi} = -1, \quad \text{for any } \hat{\mathbf{n}}
\tag{3.67}
\end{equation}

\subsubsection{Eigenspinor of $\boldsymbol{\sigma}\cdot\hat{\mathbf{n}}$}

To find the spinor satisfying $\boldsymbol{\sigma}\cdot\hat{\mathbf{n}}\chi = \chi$ (eigenvalue $+1$), where $\hat{\mathbf{n}}$ has polar angle $\beta$ and azimuthal angle $\alpha$:

\begin{equation}
\boxed{\chi = \begin{pmatrix} \cos(\beta/2)e^{-i\alpha/2} \\ \sin(\beta/2)e^{i\alpha/2} \end{pmatrix}}
\tag{3.70}
\end{equation}

This is obtained by rotating the spin-up spinor $\binom{1}{0}$ first about the $y$-axis by $\beta$, then about the $z$-axis by $\alpha$.

\section{SO(3), SU(2), and Euler Rotations}

\subsection{Orthogonal Group}

A general rotation can be characterized by three real numbers: the polar and azimuthal angles of the unit vector $\hat{\mathbf{n}}$ and the rotation angle $\phi$ itself.

\subsubsection{Parameters of a $3 \times 3$ Orthogonal Matrix}

A real $3 \times 3$ matrix has 9 entries, but the orthogonality constraint
\begin{equation}
RR^T = 1
\tag{3.71}
\end{equation}
corresponds to 6 independent equations (since $RR^T$ is symmetric with 6 independent entries). Thus there are $9 - 6 = 3$ independent parameters in $R$.

\subsubsection{Group Properties of SO(3)}

The set of all multiplication operations with orthogonal matrices forms a group:

\begin{enumerate}
\item \textbf{Closure:} $(R_1 R_2)(R_1 R_2)^T = R_1 R_2 R_2^T R_1^T = 1$ \hfill (3.72)
\item \textbf{Associativity:} $R_1(R_2 R_3) = (R_1 R_2)R_3$ \hfill (3.73)
\item \textbf{Identity:} $R \cdot 1 = 1 \cdot R = R$ \hfill (3.74)
\item \textbf{Inverse:} $RR^{-1} = R^{-1}R = 1$ \hfill (3.75)
\end{enumerate}

This group is called \textbf{SO(3)}: S = special (det $R = +1$), O = orthogonal, 3 = three dimensions.

\subsection{Unitary Unimodular Group}

The $2 \times 2$ rotation matrix (3.63) is unitary, so $|c_+|^2 + |c_-|^2 = 1$ is preserved. Furthermore, matrix (3.63) is \textbf{unimodular} (determinant = 1).

The most general unitary unimodular matrix:
\begin{equation}
\boxed{U(a, b) = \begin{pmatrix} a & b \\ -b^* & a^* \end{pmatrix}}
\tag{3.77}
\end{equation}

where $a$ and $b$ are complex numbers satisfying the \textbf{unimodular condition}:
\begin{equation}
|a|^2 + |b|^2 = 1
\tag{3.78}
\end{equation}

\textbf{Verification of unitarity:}
\begin{equation}
U(a,b)^\dagger U(a,b) = \begin{pmatrix} a^* & -b \\ b^* & a \end{pmatrix}\begin{pmatrix} a & b \\ -b^* & a^* \end{pmatrix} = 1
\tag{3.79}
\end{equation}

\textbf{Cayley-Klein parameters:} Comparing (3.63) with (3.77):
\begin{equation}
\text{Re}(a) = \cos\frac{\phi}{2}, \quad \text{Im}(a) = -n_z\sin\frac{\phi}{2}, \quad \text{Re}(b) = -n_y\sin\frac{\phi}{2}, \quad \text{Im}(b) = -n_x\sin\frac{\phi}{2}
\tag{3.80}
\end{equation}

\textbf{Group multiplication:}
\begin{equation}
U(a_1, b_1)U(a_2, b_2) = U(a_1 a_2 - b_1 b_2^*, a_1 b_2 + a_2^* b_1)
\tag{3.81}
\end{equation}

\textbf{Inverse:}
\begin{equation}
U^{-1}(a, b) = U(a^*, -b)
\tag{3.83}
\end{equation}

This group is called \textbf{SU(2)}: S = special (det = 1), U = unitary, 2 = dimensionality 2.

\subsubsection{Relationship Between SO(3) and SU(2)}

\textbf{Important:} SO(3) and SU(2) are \textit{not} isomorphic. The correspondence is \textbf{two-to-one}: for a given $R$, the corresponding $U$ is double valued.

In SO(3): $2\pi$ and $4\pi$ rotations are both the $3 \times 3$ identity matrix.

In SU(2): $2\pi$ rotation gives $-1$ (times identity), $4\pi$ rotation gives $+1$ (identity).

More generally, $U(a, b)$ and $U(-a, -b)$ both correspond to a \textit{single} $3 \times 3$ matrix in SO(3). The two groups are \textbf{locally isomorphic}.

\subsection{Euler Rotations}

An arbitrary rotation of a rigid body can be accomplished in three steps, known as \textbf{Euler rotations}:

\begin{enumerate}
\item Rotate about the $z$-axis by angle $\alpha$
\item Rotate about the $y'$-axis (body $y$-axis after first rotation) by angle $\beta$
\item Rotate about the $z'$-axis (body $z$-axis after second rotation) by angle $\gamma$
\end{enumerate}

In terms of body-axis rotations:
\begin{equation}
R(\alpha, \beta, \gamma) \equiv R_{z'}(\gamma)R_{y'}(\beta)R_z(\alpha)
\tag{3.85}
\end{equation}

\subsubsection{Conversion to Fixed-Axis Rotations}

Body-axis rotations can be expressed in terms of fixed-axis rotations:
\begin{align}
R_{y'}(\beta) &= R_z(\alpha)R_y(\beta)R_z^{-1}(\alpha) \tag{3.86} \\
R_{z'}(\gamma) &= R_{y'}(\beta)R_z(\gamma)R_{y'}^{-1}(\beta) \tag{3.87}
\end{align}

\textbf{Key result:} Euler rotations in terms of \textit{fixed}-axis rotations:
\begin{equation}
\boxed{R(\alpha, \beta, \gamma) = R_z(\alpha)R_y(\beta)R_z(\gamma)}
\tag{3.89}
\end{equation}

\subsubsection{Euler Rotations for Spin $\frac{1}{2}$}

The rotation operator in terms of Euler angles:
\begin{equation}
\mathscr{D}(\alpha, \beta, \gamma) = \mathscr{D}_z(\alpha)\mathscr{D}_y(\beta)\mathscr{D}_z(\gamma)
\tag{3.90}
\end{equation}

The $2 \times 2$ matrix representation:
\begin{align}
&\exp\left(\frac{-i\sigma_3\alpha}{2}\right)\exp\left(\frac{-i\sigma_2\beta}{2}\right)\exp\left(\frac{-i\sigma_3\gamma}{2}\right) \notag \\
&= \begin{pmatrix} e^{-i\alpha/2} & 0 \\ 0 & e^{i\alpha/2} \end{pmatrix}\begin{pmatrix} \cos(\beta/2) & -\sin(\beta/2) \\ \sin(\beta/2) & \cos(\beta/2) \end{pmatrix}\begin{pmatrix} e^{-i\gamma/2} & 0 \\ 0 & e^{i\gamma/2} \end{pmatrix} \notag \\
&= \boxed{\begin{pmatrix} e^{-i(\alpha+\gamma)/2}\cos(\beta/2) & -e^{-i(\alpha-\gamma)/2}\sin(\beta/2) \\ e^{i(\alpha-\gamma)/2}\sin(\beta/2) & e^{i(\alpha+\gamma)/2}\cos(\beta/2) \end{pmatrix}}
\tag{3.91}
\end{align}

This is the $j = \frac{1}{2}$ \textbf{irreducible representation} of the rotation operator $\mathscr{D}(\alpha, \beta, \gamma)$.

\textbf{Matrix elements notation:}
\begin{equation}
\mathscr{D}_{m'm}^{(1/2)}(\alpha, \beta, \gamma) = \left\langle j = \frac{1}{2}, m'\left|\exp\left(\frac{-iJ_z\alpha}{\hbar}\right)\exp\left(\frac{-iJ_y\beta}{\hbar}\right)\exp\left(\frac{-iJ_z\gamma}{\hbar}\right)\right|j = \frac{1}{2}, m\right\rangle
\tag{3.92}
\end{equation}

\textbf{Note:} The middle rotation is about the $y$-axis (not $x$-axis as in some classical mechanics texts) so that the off-diagonal elements are purely real.

\section{Eigenvalues and Eigenstates of Angular Momentum}

In this section we study more general angular-momentum states. We work out the eigenvalues and eigenkets of $\mathbf{J}^2$ and $J_z$ and derive the expressions for matrix elements of angular-momentum operators.

\subsection{Commutation Relations and the Ladder Operators}

\subsubsection{The $\mathbf{J}^2$ Operator}

Define the operator:
\begin{equation}
\mathbf{J}^2 \equiv J_x J_x + J_y J_y + J_z J_z
\tag{3.153}
\end{equation}

This operator commutes with every component of $\mathbf{J}$:
\begin{equation}
\boxed{[\mathbf{J}^2, J_k] = 0 \quad (k = 1, 2, 3)}
\tag{3.154}
\end{equation}

\textbf{Proof for $k = 3$:}
\begin{align}
[J_x J_x + J_y J_y + J_z J_z, J_z] &= J_x[J_x, J_z] + [J_x, J_z]J_x + J_y[J_y, J_z] + [J_y, J_z]J_y \notag \\
&= J_x(-i\hbar J_y) + (-i\hbar J_y)J_x + J_y(i\hbar J_x) + (i\hbar J_x)J_y = 0
\tag{3.155}
\end{align}

Because $J_x$, $J_y$, and $J_z$ do not commute with each other, we can choose only one of them to be diagonalized simultaneously with $\mathbf{J}^2$. By convention we choose $J_z$.

\subsubsection{Simultaneous Eigenkets}

We look for simultaneous eigenkets of $\mathbf{J}^2$ and $J_z$, with eigenvalues $a$ and $b$:
\begin{align}
\mathbf{J}^2|a, b\rangle &= a|a, b\rangle \tag{3.156a} \\
J_z|a, b\rangle &= b|a, b\rangle \tag{3.156b}
\end{align}

\subsubsection{Ladder Operators}

Define the \textbf{ladder operators}:
\begin{equation}
\boxed{J_\pm \equiv J_x \pm iJ_y}
\tag{3.157}
\end{equation}

\textbf{Key commutation relations:}
\begin{align}
[J_+, J_-] &= 2\hbar J_z \tag{3.158a} \\
[J_z, J_\pm] &= \pm\hbar J_\pm \tag{3.158b} \\
[\mathbf{J}^2, J_\pm] &= 0 \tag{3.159}
\end{align}

\textbf{Physical meaning:} The ladder operators raise/lower the $J_z$ eigenvalue by $\hbar$:
\begin{equation}
J_z(J_\pm|a, b\rangle) = ([J_z, J_\pm] + J_\pm J_z)|a, b\rangle = (b \pm \hbar)(J_\pm|a, b\rangle)
\tag{3.160}
\end{equation}

Also, $J_\pm$ does not change the $\mathbf{J}^2$ eigenvalue:
\begin{equation}
\mathbf{J}^2(J_\pm|a, b\rangle) = J_\pm\mathbf{J}^2|a, b\rangle = a(J_\pm|a, b\rangle)
\tag{3.163}
\end{equation}

We may write:
\begin{equation}
J_\pm|a, b\rangle = c_\pm|a, b \pm \hbar\rangle
\tag{3.164}
\end{equation}

\subsection{Eigenvalues of $\mathbf{J}^2$ and $J_z$}

\subsubsection{Upper Bound on $b$}

There exists an upper limit to $b$ for a given $a$:
\begin{equation}
a \geq b^2
\tag{3.165}
\end{equation}

\textbf{Proof:}
\begin{equation}
\mathbf{J}^2 - J_z^2 = \frac{1}{2}(J_+ J_- + J_- J_+) = \frac{1}{2}(J_+ J_+^\dagger + J_+^\dagger J_+)
\tag{3.166}
\end{equation}

Since $J_+ J_+^\dagger$ and $J_+^\dagger J_+$ have nonnegative expectation values:
\begin{equation}
\langle a, b|(\mathbf{J}^2 - J_z^2)|a, b\rangle \geq 0
\tag{3.168}
\end{equation}

\subsubsection{Maximum and Minimum Values}

There must exist a $b_{\max}$ such that:
\begin{equation}
J_+|a, b_{\max}\rangle = 0
\tag{3.169}
\end{equation}

Using $J_- J_+ = \mathbf{J}^2 - J_z^2 - \hbar J_z$:
\begin{equation}
(J_- J_+)|a, b_{\max}\rangle = (\mathbf{J}^2 - J_z^2 - \hbar J_z)|a, b_{\max}\rangle = 0
\tag{3.170, 3.171, 3.172}
\end{equation}

This gives:
\begin{equation}
a = b_{\max}(b_{\max} + \hbar)
\tag{3.174}
\end{equation}

Similarly, there exists a $b_{\min}$ such that $J_-|a, b_{\min}\rangle = 0$, giving:
\begin{equation}
a = b_{\min}(b_{\min} - \hbar)
\tag{3.177}
\end{equation}

Comparing (3.174) and (3.177):
\begin{equation}
b_{\max} = -b_{\min}
\tag{3.178}
\end{equation}

with $b_{\max}$ positive, and:
\begin{equation}
-b_{\max} \leq b \leq b_{\max}
\tag{3.179}
\end{equation}

\subsubsection{Quantization Condition}

Since we reach $|a, b_{\max}\rangle$ from $|a, b_{\min}\rangle$ by applying $J_+$ a finite number of times:
\begin{equation}
b_{\max} = b_{\min} + n\hbar
\tag{3.180}
\end{equation}

where $n$ is a non-negative integer. This gives:
\begin{equation}
b_{\max} = \frac{n\hbar}{2}
\tag{3.181}
\end{equation}

\subsubsection{Standard Notation}

Define $j \equiv b_{\max}/\hbar$ and $m \equiv b/\hbar$:
\begin{equation}
j = \frac{n}{2}
\tag{3.182}
\end{equation}

The eigenvalue of $\mathbf{J}^2$:
\begin{equation}
a = \hbar^2 j(j+1)
\tag{3.183}
\end{equation}

\textbf{Key results:}
\begin{align}
\boxed{\mathbf{J}^2|j, m\rangle = j(j+1)\hbar^2|j, m\rangle} \tag{3.186a} \\
\boxed{J_z|j, m\rangle = m\hbar|j, m\rangle} \tag{3.186b}
\end{align}

with $j$ either an integer or a half-integer, and:
\begin{equation}
\boxed{m = -j, -j+1, \ldots, j-1, j \quad \text{($2j+1$ states)}}
\tag{3.185}
\end{equation}

\subsection{Matrix Elements of Angular-Momentum Operators}

From the eigenvalue equations:
\begin{align}
\langle j', m'|\mathbf{J}^2|j, m\rangle &= j(j+1)\hbar^2\delta_{j'j}\delta_{m'm} \tag{3.187a} \\
\langle j', m'|J_z|j, m\rangle &= m\hbar\delta_{j'j}\delta_{m'm} \tag{3.187b}
\end{align}

\subsubsection{Matrix Elements of $J_\pm$}

To find $|c_{jm}^+|^2$, consider:
\begin{equation}
\langle j, m|J_+^\dagger J_+|j, m\rangle = \langle j, m|(\mathbf{J}^2 - J_z^2 - \hbar J_z)|j, m\rangle = \hbar^2[j(j+1) - m^2 - m]
\tag{3.188}
\end{equation}

This gives:
\begin{equation}
|c_{jm}^+|^2 = \hbar^2(j-m)(j+m+1)
\tag{3.190}
\end{equation}

Choosing $c_{jm}^+$ to be real and positive by convention:

\begin{equation}
\boxed{J_+|j, m\rangle = \sqrt{(j-m)(j+m+1)}\hbar|j, m+1\rangle}
\tag{3.191}
\end{equation}

\begin{equation}
\boxed{J_-|j, m\rangle = \sqrt{(j+m)(j-m+1)}\hbar|j, m-1\rangle}
\tag{3.192}
\end{equation}

\textbf{General matrix elements:}
\begin{equation}
\boxed{\langle j', m'|J_\pm|j, m\rangle = \sqrt{(j \mp m)(j \pm m + 1)}\hbar\delta_{j'j}\delta_{m', m\pm 1}}
\tag{3.193}
\end{equation}

\subsection{Representations of the Rotation Operator}

\subsubsection{Wigner $\mathscr{D}$-Functions}

The matrix elements of the rotation operator $\mathscr{D}(R)$ are called \textbf{Wigner functions}:
\begin{equation}
\boxed{\mathscr{D}_{m'm}^{(j)}(R) = \langle j, m'|\exp\left(\frac{-i\mathbf{J}\cdot\hat{\mathbf{n}}\phi}{\hbar}\right)|j, m\rangle}
\tag{3.194}
\end{equation}

\textbf{Important:} Rotations cannot change the $j$-value:
\begin{equation}
\mathbf{J}^2\mathscr{D}(R)|j, m\rangle = \mathscr{D}(R)\mathbf{J}^2|j, m\rangle = j(j+1)\hbar^2[\mathscr{D}(R)|j, m\rangle]
\tag{3.195}
\end{equation}

The $(2j+1) \times (2j+1)$ matrix formed by $\mathscr{D}_{m'm}^{(j)}(R)$ is the \textbf{$(2j+1)$-dimensional irreducible representation} of the rotation operator.

\subsubsection{Physical Interpretation}

When we rotate a state $|j, m\rangle$:
\begin{equation}
|j, m\rangle \to \mathscr{D}(R)|j, m\rangle
\tag{3.200}
\end{equation}

The amplitude to be found in $|j, m'\rangle$ after rotation:
\begin{equation}
\mathscr{D}(R)|j, m\rangle = \sum_{m'}|j, m'\rangle\langle j, m'|\mathscr{D}(R)|j, m\rangle = \sum_{m'}|j, m'\rangle\mathscr{D}_{m'm}^{(j)}(R)
\tag{3.201}
\end{equation}

\subsubsection{Euler Angle Representation}

Using Euler angles $(\alpha, \beta, \gamma)$:
\begin{equation}
\mathscr{D}_{m'm}^{(j)}(\alpha, \beta, \gamma) = e^{-i(m'\alpha + m\gamma)}\langle j, m'|\exp\left(\frac{-iJ_y\beta}{\hbar}\right)|j, m\rangle
\tag{3.202}
\end{equation}

Define the \textbf{reduced rotation matrix} (Wigner $d$-function):
\begin{equation}
\boxed{d_{m'm}^{(j)}(\beta) \equiv \langle j, m'|\exp\left(\frac{-iJ_y\beta}{\hbar}\right)|j, m\rangle}
\tag{3.203}
\end{equation}

so that:
\begin{equation}
\mathscr{D}_{m'm}^{(j)}(\alpha, \beta, \gamma) = e^{-im'\alpha}d_{m'm}^{(j)}(\beta)e^{-im\gamma}
\end{equation}

\subsubsection{Examples}

\textbf{For $j = 1/2$:}
\begin{equation}
d^{(1/2)} = \begin{pmatrix} \cos(\beta/2) & -\sin(\beta/2) \\ \sin(\beta/2) & \cos(\beta/2) \end{pmatrix}
\tag{3.204}
\end{equation}

\textbf{For $j = 1$:}

The $J_y$ matrix representation:
\begin{equation}
J_y^{(j=1)} = \frac{\hbar}{2}\begin{pmatrix} 0 & -\sqrt{2}i & 0 \\ \sqrt{2}i & 0 & -\sqrt{2}i \\ 0 & \sqrt{2}i & 0 \end{pmatrix}
\tag{3.206}
\end{equation}

For $j = 1$ only, $(J_y^{(j=1)}/\hbar)^3 = J_y^{(j=1)}/\hbar$, so:
\begin{equation}
\exp\left(\frac{-iJ_y\beta}{\hbar}\right) \to 1 - \left(\frac{J_y}{\hbar}\right)^2(1 - \cos\beta) - i\left(\frac{J_y}{\hbar}\right)\sin\beta
\tag{3.208}
\end{equation}

\begin{equation}
\boxed{d^{(1)}(\beta) = \begin{pmatrix} \frac{1}{2}(1+\cos\beta) & -\frac{1}{\sqrt{2}}\sin\beta & \frac{1}{2}(1-\cos\beta) \\ \frac{1}{\sqrt{2}}\sin\beta & \cos\beta & -\frac{1}{\sqrt{2}}\sin\beta \\ \frac{1}{2}(1-\cos\beta) & \frac{1}{\sqrt{2}}\sin\beta & \frac{1}{2}(1+\cos\beta) \end{pmatrix}}
\tag{3.209}
\end{equation}

\section{Orbital Angular Momentum}

When spin angular momentum is zero or can be ignored, the angular momentum $\mathbf{J}$ for a single particle is the same as \textbf{orbital angular momentum}:
\begin{equation}
\mathbf{L} = \mathbf{x} \times \mathbf{p}
\tag{3.210}
\end{equation}

\subsection{Orbital Angular Momentum as Rotation Generator}

The orbital angular-momentum operator satisfies the angular-momentum commutation relations:
\begin{equation}
[L_i, L_j] = i\varepsilon_{ijk}\hbar L_k
\tag{3.211}
\end{equation}

\textbf{Proof:}
\begin{align}
[L_x, L_y] &= [yp_z - zp_y, zp_x - xp_z] \notag \\
&= [yp_z, zp_x] + [zp_y, xp_z] \notag \\
&= yp_x[p_z, z] + p_y x[z, p_z] \notag \\
&= i\hbar(xp_y - yp_x) = i\hbar L_z
\tag{3.212}
\end{align}

\subsubsection{$L_z$ as Generator of Rotation}

Let the infinitesimal rotation operator act on a position eigenket:
\begin{align}
\left[1 - i\left(\frac{\delta\phi}{\hbar}\right)L_z\right]|x', y', z'\rangle &= \left[1 - i\left(\frac{p_y}{\hbar}\right)(\delta\phi x') + i\left(\frac{p_x}{\hbar}\right)(\delta\phi y')\right]|x', y', z'\rangle \notag \\
&= |x' - y'\delta\phi, y' + x'\delta\phi, z'\rangle
\tag{3.214}
\end{align}

This is precisely what we expect if $L_z$ generates an infinitesimal rotation about the $z$-axis.

\subsubsection{Differential Operators in Spherical Coordinates}

\begin{equation}
\boxed{\langle\mathbf{x}'|L_z|\alpha\rangle = -i\hbar\frac{\partial}{\partial\phi}\langle\mathbf{x}'|\alpha\rangle}
\tag{3.218}
\end{equation}

\begin{equation}
\langle\mathbf{x}'|L_x|\alpha\rangle = -i\hbar\left(-\sin\phi\frac{\partial}{\partial\theta} - \cot\theta\cos\phi\frac{\partial}{\partial\phi}\right)\langle\mathbf{x}'|\alpha\rangle
\tag{3.220}
\end{equation}

\begin{equation}
\langle\mathbf{x}'|L_y|\alpha\rangle = -i\hbar\left(\cos\phi\frac{\partial}{\partial\theta} - \cot\theta\sin\phi\frac{\partial}{\partial\phi}\right)\langle\mathbf{x}'|\alpha\rangle
\tag{3.221}
\end{equation}

\textbf{Ladder operators:}
\begin{equation}
\boxed{\langle\mathbf{x}'|L_\pm|\alpha\rangle = -i\hbar e^{\pm i\phi}\left(\pm i\frac{\partial}{\partial\theta} - \cot\theta\frac{\partial}{\partial\phi}\right)\langle\mathbf{x}'|\alpha\rangle}
\tag{3.222}
\end{equation}

\textbf{$\mathbf{L}^2$ operator:}
Using $\mathbf{L}^2 = L_z^2 + \frac{1}{2}(L_+ L_- + L_- L_+)$:
\begin{equation}
\boxed{\langle\mathbf{x}'|\mathbf{L}^2|\alpha\rangle = -\hbar^2\left[\frac{1}{\sin^2\theta}\frac{\partial^2}{\partial\phi^2} + \frac{1}{\sin\theta}\frac{\partial}{\partial\theta}\left(\sin\theta\frac{\partial}{\partial\theta}\right)\right]\langle\mathbf{x}'|\alpha\rangle}
\tag{3.224}
\end{equation}

This is just the angular part of the Laplacian in spherical coordinates (apart from $1/r^2$).

\subsubsection{Important Operator Identity}

\begin{equation}
\boxed{\mathbf{L}^2 = \mathbf{x}^2\mathbf{p}^2 - (\mathbf{x}\cdot\mathbf{p})^2 + i\hbar\mathbf{x}\cdot\mathbf{p}}
\tag{3.225}
\end{equation}

This leads to the kinetic energy in spherical coordinates:
\begin{equation}
\frac{1}{2m}\langle\mathbf{x}'|\mathbf{p}^2|\alpha\rangle = -\frac{\hbar^2}{2m}\left[\left(\frac{\partial^2}{\partial r^2} + \frac{2}{r}\frac{\partial}{\partial r}\right)\langle\mathbf{x}'|\alpha\rangle - \frac{1}{\hbar^2 r^2}\langle\mathbf{x}'|\mathbf{L}^2|\alpha\rangle\right]
\tag{3.230}
\end{equation}

\subsection{Spherical Harmonics}

For a spherically symmetric potential, energy eigenfunctions can be written as:
\begin{equation}
\langle\mathbf{x}'|n, l, m\rangle = R_{nl}(r)Y_l^m(\theta, \phi)
\tag{3.231}
\end{equation}

Define the \textbf{direction eigenket} $|\hat{\mathbf{n}}\rangle$:
\begin{equation}
\langle\hat{\mathbf{n}}|l, m\rangle = Y_l^m(\theta, \phi) = Y_l^m(\hat{\mathbf{n}})
\tag{3.232}
\end{equation}

$Y_l^m(\theta, \phi)$ is the amplitude for a state characterized by $l$, $m$ to be found in the direction $\hat{\mathbf{n}}$ specified by $\theta$ and $\phi$.

\subsubsection{Eigenvalue Equations}

The eigenvalue equations translate to differential equations:

\begin{equation}
-i\hbar\frac{\partial}{\partial\phi}Y_l^m(\theta, \phi) = m\hbar Y_l^m(\theta, \phi)
\tag{3.235}
\end{equation}

implying $Y_l^m(\theta, \phi)$ has $\phi$-dependence $e^{im\phi}$.

\begin{equation}
\left[\frac{1}{\sin\theta}\frac{\partial}{\partial\theta}\left(\sin\theta\frac{\partial}{\partial\theta}\right) + \frac{1}{\sin^2\theta}\frac{\partial^2}{\partial\phi^2} + l(l+1)\right]Y_l^m = 0
\tag{3.237}
\end{equation}

\subsubsection{Orthogonality}

\begin{equation}
\int_0^{2\pi}d\phi\int_{-1}^{1}d(\cos\theta)\,Y_{l'}^{m'^*}(\theta, \phi)Y_l^m(\theta, \phi) = \delta_{ll'}\delta_{mm'}
\tag{3.239}
\end{equation}

with completeness relation:
\begin{equation}
\int d\Omega_{\hat{\mathbf{n}}}|\hat{\mathbf{n}}\rangle\langle\hat{\mathbf{n}}| = 1
\tag{3.240}
\end{equation}

\subsubsection{Construction of Spherical Harmonics}

Start with $m = l$ case. From $L_+|l, l\rangle = 0$ and (3.222):
\begin{equation}
\langle\hat{\mathbf{n}}|l, l\rangle = Y_l^l(\theta, \phi) = c_l e^{il\phi}\sin^l\theta
\tag{3.243}
\end{equation}

with normalization constant:
\begin{equation}
c_l = \frac{(-1)^l}{2^l l!}\sqrt{\frac{(2l+1)(2l)!}{4\pi}}
\tag{3.244}
\end{equation}

Apply $L_-$ successively to get other $m$ values:
\begin{equation}
\langle\hat{\mathbf{n}}|l, m-1\rangle = \frac{1}{\sqrt{(l+m)(l-m+1)}}e^{-i\phi}\left(-\frac{\partial}{\partial\theta} + i\cot\theta\frac{\partial}{\partial\phi}\right)\langle\hat{\mathbf{n}}|l, m\rangle
\tag{3.245}
\end{equation}

\textbf{Result for $m \geq 0$:}
\begin{equation}
\boxed{Y_l^m(\theta, \phi) = \frac{(-1)^l}{2^l l!}\sqrt{\frac{(2l+1)}{4\pi}\frac{(l+m)!}{(l-m)!}}\,e^{im\phi}\frac{1}{\sin^m\theta}\frac{d^{l-m}}{d(\cos\theta)^{l-m}}(\sin\theta)^{2l}}
\tag{3.246}
\end{equation}

\textbf{For negative $m$:}
\begin{equation}
\boxed{Y_l^{-m}(\theta, \phi) = (-1)^m[Y_l^m(\theta, \phi)]^*}
\tag{3.247}
\end{equation}

\textbf{For $m = 0$:}
\begin{equation}
Y_l^0(\theta, \phi) = \sqrt{\frac{2l+1}{4\pi}}P_l(\cos\theta)
\tag{3.248}
\end{equation}

where $P_l$ is the Legendre polynomial.

\subsubsection{Why $l$ Must Be an Integer}

Several arguments show that orbital angular momentum cannot have half-integer $l$:

\textbf{1. Single-valuedness:} For half-integer $m$, $e^{im(2\pi)} = -1$, so the wave function would not be single-valued under a $2\pi$ rotation.

\textbf{2. Direct calculation:} Under a $2\pi$ rotation about the $z$-axis:
\begin{equation}
\langle\mathbf{x}'|\exp\left(\frac{-iL_z 2\pi}{\hbar}\right)|\alpha\rangle = \langle x'\cos 2\pi + y'\sin 2\pi, y'\cos 2\pi - x'\sin 2\pi, z'|\alpha\rangle = \langle\mathbf{x}'|\alpha\rangle
\tag{3.250}
\end{equation}

No sign change! This is unlike spin-$\frac{1}{2}$ systems.

\textbf{3. Singularity at poles:} If $l = m = \frac{1}{2}$ were possible:
\begin{equation}
Y_{1/2}^{1/2}(\theta, \phi) = c_{1/2}e^{i\phi/2}\sqrt{\sin\theta}
\tag{3.251}
\end{equation}

Applying $L_-$ gives $Y_{1/2}^{-1/2} \propto \cot\theta\sqrt{\sin\theta}$, which is singular at $\theta = 0, \pi$.

\textbf{4. Completeness:} The integer-$l$ spherical harmonics form a complete set.

\subsection{Spherical Harmonics as Rotation Matrices}

We can construct the direction eigenket $|\hat{\mathbf{n}}\rangle$ by rotating $|\hat{\mathbf{z}}\rangle$:
\begin{equation}
|\hat{\mathbf{n}}\rangle = \mathscr{D}(R)|\hat{\mathbf{z}}\rangle
\tag{3.255}
\end{equation}

Using Euler angles $\alpha = \phi$, $\beta = \theta$, $\gamma = 0$:
\begin{equation}
\langle l, m'|\hat{\mathbf{n}}\rangle = \sum_m \mathscr{D}_{m'm}^{(l)}(\alpha = \phi, \beta = \theta, \gamma = 0)\langle l, m|\hat{\mathbf{z}}\rangle
\tag{3.258}
\end{equation}

Since $|\hat{\mathbf{z}}\rangle$ is an eigenket of $L_z$ with eigenvalue zero, $\langle l, m|\hat{\mathbf{z}}\rangle = 0$ for $m \neq 0$:
\begin{equation}
\langle l, m|\hat{\mathbf{z}}\rangle = \sqrt{\frac{2l+1}{4\pi}}\delta_{m0}
\tag{3.259}
\end{equation}

\textbf{Key result connecting spherical harmonics and rotation matrices:}
\begin{equation}
\boxed{Y_l^{m^*}(\theta, \phi) = \sqrt{\frac{2l+1}{4\pi}}\mathscr{D}_{m0}^{(l)}(\alpha = \phi, \beta = \theta, \gamma = 0)}
\tag{3.260}
\end{equation}

or equivalently:
\begin{equation}
\mathscr{D}_{m0}^{(l)}(\alpha, \beta, \gamma = 0) = \sqrt{\frac{4\pi}{2l+1}}Y_l^{m^*}(\theta, \phi)\bigg|_{\theta = \beta, \phi = \alpha}
\tag{3.261}
\end{equation}

\textbf{Special case $m = 0$:}
\begin{equation}
\boxed{d_{00}^{(l)}(\beta)\big|_{\beta = \theta} = P_l(\cos\theta)}
\tag{3.262}
\end{equation}

The Legendre polynomial is the $m = m' = 0$ element of the reduced rotation matrix.

\section{Schrödinger's Equation for Central Potentials}

Problems described by Hamiltonians of the form
\begin{equation}
H = \frac{\mathbf{p}^2}{2m} + V(r), \qquad r^2 = \mathbf{x}^2
\tag{3.263}
\end{equation}
are called ``central potential'' or ``central force'' problems. The fundamental importance lies in spherical symmetry:
\begin{equation}
[\mathbf{L}, \mathbf{p}^2] = [\mathbf{L}, \mathbf{x}^2] = 0
\tag{3.264}
\end{equation}
and therefore
\begin{equation}
\boxed{[\mathbf{L}, H] = [\mathbf{L}^2, H] = 0}
\tag{3.265}
\end{equation}

\subsection{The Radial Equation}

Equation (3.265) makes clear that we should search for energy eigenstates $|\alpha\rangle = |Elm\rangle$ where
\begin{align}
H|Elm\rangle &= E|Elm\rangle \tag{3.266} \\
\mathbf{L}^2|Elm\rangle &= l(l+1)\hbar^2|Elm\rangle \tag{3.267} \\
L_z|Elm\rangle &= m\hbar|Elm\rangle \tag{3.268}
\end{align}

Combining with (3.230) and (3.231), we arrive at the \textbf{radial equation}:
\begin{equation}
\boxed{\left[-\frac{\hbar^2}{2mr^2}\frac{d}{dr}\left(r^2\frac{d}{dr}\right) + \frac{l(l+1)\hbar^2}{2mr^2} + V(r)\right]R_{El}(r) = ER_{El}(r)}
\tag{3.269}
\end{equation}

\subsubsection{Substitution $u_{El}(r) = rR_{El}(r)$}

Making the substitution $R_{El}(r) = u_{El}(r)/r$:
\begin{equation}
\boxed{-\frac{\hbar^2}{2m}\frac{d^2 u_{El}}{dr^2} + \left[\frac{l(l+1)\hbar^2}{2mr^2} + V(r)\right]u_{El}(r) = Eu_{El}(r)}
\tag{3.271}
\end{equation}

This looks like a 1D Schrödinger equation with \textbf{effective potential}:
\begin{equation}
\boxed{V_{\text{eff}}(r) = V(r) + \frac{l(l+1)\hbar^2}{2mr^2}}
\tag{3.273}
\end{equation}

The term $l(l+1)\hbar^2/2mr^2$ is the \textbf{angular momentum barrier} (centrifugal barrier).

\subsubsection{Boundary Conditions}

\textbf{At $r \to 0$:} For well-behaved potentials with $\lim_{r\to 0}r^2 V(r) = 0$:
\begin{equation}
\frac{d^2 u_{El}}{dr^2} = \frac{l(l+1)}{r^2}u_{El}(r) \quad (r \to 0)
\tag{3.274}
\end{equation}

General solution: $u_{El}(r) = Ar^{l+1} + Br^{-l}$. Regularity requires $B = 0$:
\begin{equation}
\boxed{R_{El}(r) \to r^l \quad \text{as } r \to 0}
\tag{3.276}
\end{equation}

\textbf{At $r \to \infty$:} For bound states with $V(r) \to 0$:
\begin{equation}
\frac{d^2 u_E}{dr^2} = \kappa^2 u, \qquad \kappa^2 \equiv -2mE/\hbar^2 > 0
\tag{3.277}
\end{equation}

Solution: $u_E(r) \propto e^{-\kappa r}$.

\textbf{General ansatz for bound states:}
\begin{equation}
u_{El}(\rho) = \rho^{l+1}e^{-\rho}w(\rho)
\tag{3.278}
\end{equation}
where $\rho = \kappa r$ and $w(\rho)$ is ``well-behaved.''

\subsection{The Free Particle and Infinite Spherical Well}

For a free particle, $E = \hbar^2 k^2/2m$ and $\rho \equiv kr$. The radial equation becomes:
\begin{equation}
\frac{d^2 R}{d\rho^2} + \frac{2}{\rho}\frac{dR}{d\rho} + \left[1 - \frac{l(l+1)}{\rho^2}\right]R = 0
\tag{3.281}
\end{equation}

Solutions are \textbf{spherical Bessel functions}:
\begin{align}
j_l(\rho) &= (-\rho)^l\left[\frac{1}{\rho}\frac{d}{d\rho}\right]^l\left(\frac{\sin\rho}{\rho}\right) \tag{3.282a} \\
n_l(\rho) &= -(-\rho)^l\left[\frac{1}{\rho}\frac{d}{d\rho}\right]^l\left(\frac{\cos\rho}{\rho}\right) \tag{3.282b}
\end{align}

\textbf{Behavior:} As $\rho \to 0$: $j_l(\rho) \to \rho^l$, $n_l(\rho) \to \rho^{-l-1}$. Hence $j_l(\rho)$ are the physical solutions.

\textbf{First few spherical Bessel functions:}
\begin{align}
j_0(\rho) &= \frac{\sin\rho}{\rho} \tag{3.284} \\
j_1(\rho) &= \frac{\sin\rho}{\rho^2} - \frac{\cos\rho}{\rho} \tag{3.285} \\
j_2(\rho) &= \left[\frac{3}{\rho^3} - \frac{1}{\rho}\right]\sin\rho - \frac{3\cos\rho}{\rho^2} \tag{3.286}
\end{align}

\textbf{Infinite spherical well:} $V(r) = 0$ for $r < a$, $V = \infty$ for $r > a$.

Boundary condition: $j_l(ka) = 0$, i.e., $ka$ equals zeros of spherical Bessel functions.

\textbf{Energy levels:}
\begin{align}
E_{l=0} &= \frac{\hbar^2}{2ma^2}[\pi^2, (2\pi)^2, (3\pi)^2, \ldots] \tag{3.287} \\
E_{l=1} &= \frac{\hbar^2}{2ma^2}[4.49^2, 7.73^2, 10.90^2, \ldots] \tag{3.288} \\
E_{l=2} &= \frac{\hbar^2}{2ma^2}[5.76^2, 9.10^2, 12.32^2, \ldots] \tag{3.289}
\end{align}

No degeneracies in $l$ for the infinite spherical well.

\subsection{The Isotropic Harmonic Oscillator}

\begin{equation}
H = \frac{\mathbf{p}^2}{2m} + \frac{1}{2}m\omega^2 r^2
\tag{3.290}
\end{equation}

Introduce dimensionless variables:
\begin{equation}
E = \frac{1}{2}\hbar\omega\lambda, \qquad r = \left[\frac{\hbar}{m\omega}\right]^{1/2}\rho
\tag{3.291}
\end{equation}

Ansatz: $u(\rho) = \rho^{l+1}e^{-\rho^2/2}f(\rho)$, leading to:
\begin{equation}
\rho\frac{d^2 f}{d\rho^2} + 2[(l+1) - \rho^2]\frac{df}{d\rho} + [\lambda - (2l+3)]\rho f(\rho) = 0
\tag{3.294}
\end{equation}

Series solution with termination condition $2n + 2l + 3 - \lambda = 0$ gives:

\begin{equation}
\boxed{E_{ql} = \left(2q + l + \frac{3}{2}\right)\hbar\omega \equiv \left(N + \frac{3}{2}\right)\hbar\omega}
\tag{3.302}
\end{equation}

where $q = 0, 1, 2, \ldots$ (radial quantum number), $l = 0, 1, 2, \ldots$, and $N \equiv 2q + l$ is the \textbf{principal quantum number}.

\textbf{Key properties:}
\begin{itemize}
\item $q$ counts the number of nodes in the radial function
\item For even (odd) $N$, only even (odd) values of $l$ are allowed
\item The parity of the wave function is even or odd with the value of $N$
\end{itemize}

\textbf{Alternative basis:} Using $H = H_x + H_y + H_z$ where $H_i = a_i^\dagger a_i + \frac{1}{2}$:
\begin{equation}
E = \left(n_x + \frac{1}{2} + n_y + \frac{1}{2} + n_z + \frac{1}{2}\right)\hbar\omega = \left(N + \frac{3}{2}\right)\hbar\omega
\tag{3.304}
\end{equation}
where $N = n_x + n_y + n_z$. The degeneracy is the same regardless of basis.

\subsection{The Coulomb Potential}

\begin{equation}
V(\mathbf{x}) = -\frac{Ze^2}{r}
\tag{3.305}
\end{equation}

Define:
\begin{equation}
\rho_0 = \left[\frac{2m}{-E}\right]^{1/2}\frac{Ze^2}{\hbar} = \left[\frac{2mc^2}{-E}\right]^{1/2}Z\alpha
\tag{3.306}
\end{equation}
where $\alpha \equiv e^2/\hbar c \approx 1/137$ is the \textbf{fine-structure constant}.

The radial equation becomes \textbf{Kummer's equation}:
\begin{equation}
x\frac{d^2 F}{dx^2} + (c - x)\frac{dF}{dx} - aF = 0
\tag{3.308}
\end{equation}
with $x = 2\rho$, $c = 2(l+1)$, $2a = 2(l+1) - \rho_0$.

Solution: \textbf{confluent hypergeometric function}:
\begin{equation}
F(a; c; x) = 1 + \frac{a}{c}\frac{x}{1!} + \frac{a(a+1)}{c(c+1)}\frac{x^2}{2!} + \cdots
\tag{3.310}
\end{equation}

\textbf{Quantization:} Series must terminate, requiring $a + N = 0$ for some non-negative integer $N$:
\begin{equation}
\rho_0 = 2(N + l + 1), \qquad N = 0, 1, 2, \ldots, \quad l = 0, 1, 2, \ldots
\tag{3.312}
\end{equation}

Define the \textbf{principal quantum number}:
\begin{equation}
\boxed{n \equiv N + l + 1 = 1, 2, 3, \ldots \quad \text{where} \quad l = 0, 1, \ldots, n-1}
\tag{3.313}
\end{equation}

\textbf{Energy eigenvalues (Balmer formula):}
\begin{equation}
\boxed{E = -\frac{1}{2}mc^2\frac{Z^2\alpha^2}{n^2} = -13.6\text{ eV}\frac{Z^2}{n^2}}
\tag{3.315}
\end{equation}

\textbf{Bohr radius:}
\begin{equation}
\boxed{a_0 = \frac{\hbar}{mc\alpha} = \frac{\hbar^2}{me^2} = 0.53 \text{ \AA}}
\tag{3.317}
\end{equation}

\textbf{Degeneracy:}
\begin{equation}
\boxed{\text{Degeneracy} = \sum_{l=0}^{n-1}(2l+1) = n^2}
\tag{3.318}
\end{equation}

This degeneracy is not accidental but reflects a subtle symmetry of the Coulomb potential.

\textbf{Hydrogen wave functions:}
\begin{equation}
\psi_{nlm}(\mathbf{x}) = \langle\mathbf{x}|nlm\rangle = R_{nl}(r)Y_l^m(\theta, \phi)
\tag{3.319}
\end{equation}

\begin{equation}
R_{nl}(r) = \frac{1}{(2l+1)!}\left(\frac{2Zr}{na_0}\right)^l e^{-Zr/na_0}\left[\left(\frac{2Z}{na_0}\right)^3\frac{(n+l)!}{2n(n-l-1)!}\right]^{1/2} F(-n+l+1; 2l+2; 2Zr/na_0)
\tag{3.320}
\end{equation}

\textbf{Key properties:}
\begin{itemize}
\item Only $l = 0$ (s-states) wave functions are nonzero at the origin
\item There are $n - l - 1$ nodes in the radial wave function
\end{itemize}

\section{Tensor Operators}

\subsection{Vector Operators}

A \textbf{vector operator} $\mathbf{V}$ in quantum mechanics is defined by how its expectation value transforms under rotations. Under the state transformation $|\alpha\rangle \to \mathscr{D}(R)|\alpha\rangle$, we require:
\begin{equation}
\langle\alpha|V_i|\alpha\rangle \to \langle\alpha|\mathscr{D}^\dagger(R)V_i\mathscr{D}(R)|\alpha\rangle = \sum_j R_{ij}\langle\alpha|V_j|\alpha\rangle
\tag{3.445}
\end{equation}

This implies the \textbf{operator equation}:
\begin{equation}
\mathscr{D}^\dagger(R)V_i\mathscr{D}(R) = \sum_j R_{ij}V_j
\tag{3.446}
\end{equation}

For an infinitesimal rotation with $\mathscr{D}(R) = 1 - i\varepsilon\mathbf{J}\cdot\hat{\mathbf{n}}/\hbar$, this gives the \textbf{defining commutation relation for a vector operator}:
\begin{equation}
\boxed{[V_i, J_j] = i\varepsilon_{ijk}\hbar V_k}
\tag{3.451}
\end{equation}

\textbf{Special cases:} The angular momentum commutation relations are a special case with $V_i \to J_i$. Other examples:
\begin{equation}
[y, L_z] = i\hbar x, \quad [x, L_z] = -i\hbar y, \quad [p_x, L_z] = -i\hbar p_y, \quad [p_y, L_z] = i\hbar p_x
\tag{3.451'}
\end{equation}

\subsection{Cartesian Tensors Versus Irreducible Tensors}

A \textbf{Cartesian tensor} of rank $n$ transforms as:
\begin{equation}
T_{ijk\ldots} \to \sum_{i'}\sum_{j'}\sum_{k'}\cdots R_{ii'}R_{jj'}\cdots T_{i'j'k'\ldots}
\tag{3.454}
\end{equation}

The simplest rank-2 Cartesian tensor is a \textbf{dyadic} $T_{ij} = U_i V_j$. However, Cartesian tensors are \textit{reducible}. The dyadic decomposes as:
\begin{equation}
U_i V_j = \frac{\mathbf{U}\cdot\mathbf{V}}{3}\delta_{ij} + \frac{(U_i V_j - U_j V_i)}{2} + \left(\frac{U_i V_j + U_j V_i}{2} - \frac{\mathbf{U}\cdot\mathbf{V}}{3}\delta_{ij}\right)
\tag{3.456}
\end{equation}

The three terms are: (1) scalar (1 component, $l=0$), (2) antisymmetric tensor $\propto (\mathbf{U}\times\mathbf{V})_k$ (3 components, $l=1$), (3) symmetric traceless tensor (5 components, $l=2$). This gives $3 \times 3 = 1 + 3 + 5$.

\subsubsection{Spherical Tensors}

A \textbf{spherical tensor} of rank $k$ is constructed by replacing the unit vector $\hat{\mathbf{n}}$ in $Y_l^m(\hat{\mathbf{n}})$ with a vector operator $\mathbf{V}$:
\begin{equation}
T_q^{(k)} = Y_{l=k}^{m=q}(\mathbf{V})
\tag{3.458}
\end{equation}

\textbf{Rank 1 spherical tensor components:}
\begin{equation}
\boxed{T_0^{(1)} = \sqrt{\frac{3}{4\pi}}V_z, \qquad T_{\pm 1}^{(1)} = \mp\sqrt{\frac{3}{4\pi}}\frac{V_x \pm iV_y}{\sqrt{2}}}
\tag{3.459}
\end{equation}

\textbf{Rank 2 example:}
\begin{equation}
T_{\pm 2}^{(2)} = \sqrt{\frac{15}{32\pi}}(V_x \pm iV_y)^2
\tag{3.460}
\end{equation}

\subsubsection{Transformation Properties}

A spherical tensor operator of rank $k$ with $(2k+1)$ components transforms under rotation as:
\begin{equation}
\boxed{\mathscr{D}^\dagger(R)T_q^{(k)}\mathscr{D}(R) = \sum_{q'=-k}^{k}\mathscr{D}_{qq'}^{(k)*}(R)T_{q'}^{(k)}}
\tag{3.465a}
\end{equation}
or equivalently:
\begin{equation}
\mathscr{D}(R)T_q^{(k)}\mathscr{D}^\dagger(R) = \sum_{q'=-k}^{k}\mathscr{D}_{q'q}^{(k)}(R)T_{q'}^{(k)}
\tag{3.465b}
\end{equation}

\subsubsection{Commutation Relations for Spherical Tensors}

From the infinitesimal form, we obtain the \textbf{defining commutation relations}:
\begin{align}
\boxed{[J_z, T_q^{(k)}] = \hbar q T_q^{(k)}} \tag{3.468a} \\[6pt]
\boxed{[J_\pm, T_q^{(k)}] = \hbar\sqrt{(k \mp q)(k \pm q + 1)}\, T_{q\pm 1}^{(k)}} \tag{3.468b}
\end{align}

These are completely analogous to $[J_z, |jm\rangle] = \hbar m|jm\rangle$ and the ladder operator relations.

\subsection{Product of Tensors}

Tensor products of spherical tensors can be formed using Clebsch-Gordan coefficients:

\textbf{Theorem:} Let $X_{q_1}^{(k_1)}$ and $Z_{q_2}^{(k_2)}$ be irreducible spherical tensors of rank $k_1$ and $k_2$. Then
\begin{equation}
\boxed{T_q^{(k)} = \sum_{q_1}\sum_{q_2}\langle k_1 k_2; q_1 q_2|k_1 k_2; kq\rangle X_{q_1}^{(k_1)}Z_{q_2}^{(k_2)}}
\tag{3.470}
\end{equation}
is a spherical (irreducible) tensor of rank $k$.

\textbf{Explicit products from two vectors $\mathbf{U}$, $\mathbf{V}$:}
\begin{align}
T_0^{(0)} &= \frac{-\mathbf{U}\cdot\mathbf{V}}{3} = \frac{U_{+1}V_{-1} + U_{-1}V_{+1} - U_0 V_0}{3} \notag \\[4pt]
T_q^{(1)} &= \frac{(\mathbf{U}\times\mathbf{V})_q}{i\sqrt{2}} \notag \\[4pt]
T_{\pm 2}^{(2)} &= U_{\pm 1}V_{\pm 1}, \quad T_{\pm 1}^{(2)} = \frac{U_{\pm 1}V_0 + U_0 V_{\pm 1}}{\sqrt{2}} \tag{3.469} \\[4pt]
T_0^{(2)} &= \frac{U_{+1}V_{-1} + 2U_0 V_0 + U_{-1}V_{+1}}{\sqrt{6}} \notag
\end{align}
where $U_{\pm 1} = \mp(U_x \pm iU_y)/\sqrt{2}$ and $U_0 = U_z$.

\textbf{Normalization:} the $T^{(2)}$ components above are exactly the CG products (3.470) with $k_1 = k_2 = 1$. The $T^{(0)}$ and $T^{(1)}$ above (Sakurai's (3.469)) are spherical tensors but not the CG products; (3.470) gives
$$T_0^{(0)} = \frac{-\mathbf{U}\cdot\mathbf{V}}{\sqrt{3}}, \qquad T_q^{(1)} = \frac{i(\mathbf{U}\times\mathbf{V})_q}{\sqrt{2}},$$
which differ from the ones above by the constant factors $\sqrt{3}$ and $-1$.

\subsection{Matrix Elements of Tensor Operators: The Wigner-Eckart Theorem}

\subsubsection{Selection Rules}

\textbf{$m$-selection rule:}
\begin{equation}
\boxed{\langle\alpha', j'm'|T_q^{(k)}|\alpha, jm\rangle = 0 \quad \text{unless } m' = q + m}
\tag{3.471}
\end{equation}

\textit{Proof:} From $[J_z, T_q^{(k)}] = \hbar q T_q^{(k)}$, we have
\[
\langle\alpha', j'm'|\left([J_z, T_q^{(k)}] - \hbar q T_q^{(k)}\right)|\alpha, jm\rangle = [(m' - m)\hbar - \hbar q]\langle\alpha', j'm'|T_q^{(k)}|\alpha, jm\rangle = 0
\]

\textbf{Triangular relation:}
\begin{equation}
\boxed{|j - k| \leq j' \leq j + k}
\tag{3.475}
\end{equation}

This follows from the requirement that the Clebsch-Gordan coefficient be nonvanishing.

\subsubsection{The Wigner-Eckart Theorem}

\begin{tcolorbox}[colback=blue!5!white, colframe=blue!75!black, title=Wigner-Eckart Theorem]
The matrix elements of tensor operators with respect to angular-momentum eigenstates satisfy:
\begin{equation}
\boxed{\langle\alpha', j'm'|T_q^{(k)}|\alpha, jm\rangle = \langle jk; mq|jk; j'm'\rangle \frac{\langle\alpha' j'||T^{(k)}||\alpha j\rangle}{\sqrt{2j'+1}}}
\tag{3.474}
\end{equation}
where the \textbf{double-bar matrix element} (reduced matrix element) $\langle\alpha' j'||T^{(k)}||\alpha j\rangle$ is independent of $m$, $m'$, and $q$.
\end{tcolorbox}

\textbf{Physical interpretation:}
\begin{itemize}
\item The first factor (Clebsch-Gordan coefficient) depends only on the \textit{geometry} --- how the system is oriented with respect to the $z$-axis.
\item The second factor (reduced matrix element) contains all the \textit{dynamics} --- radial integrals, etc. It is independent of magnetic quantum numbers.
\item To evaluate matrix elements for all combinations of $m$, $m'$, $q$, it suffices to compute just \textit{one} of them; all others are related by Clebsch-Gordan coefficients.
\end{itemize}

\textit{Proof:} Using (3.468b), we have the recursion relation:
\begin{align}
&\sqrt{(j' \pm m')(j' \mp m' + 1)}\langle\alpha', j', m' \mp 1|T_q^{(k)}|\alpha, jm\rangle \notag \\
&= \sqrt{(j \mp m)(j \pm m + 1)}\langle\alpha', j'm'|T_q^{(k)}|\alpha, j, m \pm 1\rangle \notag \\
&\quad + \sqrt{(k \mp q)(k \pm q + 1)}\langle\alpha', j'm'|T_{q\pm 1}^{(k)}|\alpha, jm\rangle
\tag{3.477}
\end{align}

This has the same form as the Clebsch-Gordan recursion relation (3.369) under $j' \to j$, $m' \to m$, $j \to j_1$, $m \to m_1$, $k \to j_2$, $q \to m_2$. Since both satisfy the same recursion, their solutions are proportional:
\begin{equation}
\langle\alpha', j'm'|T_{q\pm 1}^{(k)}|\alpha, jm\rangle = (\text{constant independent of } m, q, m')\langle jk; m\, q\pm 1|jk; j'm'\rangle
\tag{3.480}
\end{equation}
which proves the theorem. $\square$

\subsubsection{Examples}

\textbf{Example 1: Scalar operator ($k = 0$)}

For $T_0^{(0)} = S$:
\begin{equation}
\langle\alpha', j'm'|S|\alpha, jm\rangle = \delta_{jj'}\delta_{mm'}\frac{\langle\alpha' j'||S||\alpha j\rangle}{\sqrt{2j'+1}}
\tag{3.481}
\end{equation}
A scalar operator cannot change $j$ or $m$ values.

\textbf{Example 2: Vector operator ($k = 1$) --- Dipole Selection Rules}

For a vector operator with spherical components $V_{q=\pm 1, 0}$:
\begin{equation}
\boxed{\Delta m \equiv m' - m = \pm 1, 0, \qquad \Delta j \equiv j' - j = \pm 1, 0}
\tag{3.482}
\end{equation}

In addition, the $0 \to 0$ transition is \textbf{forbidden}. This is the \textbf{electric dipole selection rule} in the long-wavelength limit.

\subsubsection{The Projection Theorem}

For $j = j'$, the Wigner-Eckart theorem applied to a vector operator takes a particularly simple form:

\begin{tcolorbox}[colback=green!5!white, colframe=green!75!black, title=Projection Theorem]
\begin{equation}
\boxed{\langle\alpha', jm'|V_q|\alpha, jm\rangle = \frac{\langle\alpha', jm|\mathbf{J}\cdot\mathbf{V}|\alpha, jm\rangle}{\hbar^2 j(j+1)}\langle jm'|J_q|jm\rangle}
\tag{3.483}
\end{equation}
\end{tcolorbox}

where $J_{\pm 1} = \mp(J_x \pm iJ_y)/\sqrt{2} = \mp J_\pm/\sqrt{2}$ and $J_0 = J_z$.

\textit{Proof:} By the Wigner-Eckart theorem applied to both $V_q$ and $J_q$:
\begin{equation}
\frac{\langle\alpha', jm'|V_q|\alpha, jm\rangle}{\langle\alpha, jm'|J_q|\alpha, jm\rangle} = \frac{\langle\alpha' j||\mathbf{V}||\alpha j\rangle}{\langle\alpha j||\mathbf{J}||\alpha j\rangle}
\tag{3.487}
\end{equation}

Since $\mathbf{J}\cdot\mathbf{V}$ is a scalar operator (independent of $m$), and using $\langle\alpha, jm|\mathbf{J}^2|\alpha, jm\rangle = j(j+1)\hbar^2$, we obtain (3.483). $\square$

\textbf{Physical meaning:} Within a fixed-$j$ subspace, the matrix elements of any vector operator are proportional to those of $\mathbf{J}$. This is the quantum analog of the classical result that the ``component of $\mathbf{V}$ along $\mathbf{J}$'' determines the behavior.

\end{document}
